% La Deuxième Guerre mondiale — Histoire Tle Tech — Cours signé OnBuch+
% Programme officiel MINESEC. Compiler avec : tectonic N01-deuxieme-guerre-mondiale.tex
\newcommand{\DOCMATIERE}{Histoire}
\newcommand{\DOCNIVEAU}{Tle Tech}
\newcommand{\DOCMODULE}{Histoire – classe de Terminale technique}
\newcommand{\DOCLECON}{N01}
\newcommand{\DOCTITRE}{La Deuxième Guerre mondiale}
% OnBuch+ — préambule « cours signés » ENRICHI (compilé avec tectonic / XeLaTeX)
% Système visuel « Pop » d'OnBuch+ : fond crème (jamais blanc), bordures encre
% épaisses, ombres « dures » décalées, titres Archivo Black en capitales.
% Version MATHÉMATIQUES : graphes (pgfplots/tikz), boîtes pédagogiques
% (définition, propriété, méthode, attention, le savais-tu, exercice résolu,
% exercice, TP, bilan), page de garde et signature OnBuch+ ancrée en bas.
%
% Macros attendues dans main.tex AVANT \input{preamble} :
%   \DOCMATIERE  \DOCNIVEAU  \DOCMODULE  \DOCLECON (numéro)  \DOCTITRE
\documentclass[11pt]{article}

\usepackage{fontspec}
\usepackage[french]{babel}
\usepackage[a4paper,margin=2cm,top=3.4cm,bottom=2.8cm,headheight=1.4cm,footskip=1.3cm]{geometry}
\usepackage{amsmath,amssymb,amsfonts}
\usepackage{mathtools} % \xmapsto, \coloneqq... souvent utilisés par les modèles
\usepackage{cancel} % simplifications barrées (bilans de forces)
% \abs{z} (module, valeur absolue) et \norm{u} (norme), utilisés par les modèles
\providecommand{\abs}{}\let\abs\relax\DeclarePairedDelimiter{\abs}{\lvert}{\rvert}
\providecommand{\norm}{}\let\norm\relax\DeclarePairedDelimiter{\norm}{\lVert}{\rVert}
\usepackage[table]{xcolor} % [table] : fournit aussi \rowcolors (alternance de lignes)
\usepackage{pagecolor}
\usepackage{tikz}
\usetikzlibrary{arrows.meta,positioning,calc,shapes.geometric,shapes.misc,shapes.arrows,decorations.pathmorphing,decorations.pathreplacing,decorations.markings,patterns,shadows,mindmap,trees,fit,backgrounds,matrix,intersections,angles,quotes}
\usepackage{pgfplots}
\usepackage[version=4]{mhchem} % équations nucléaires \ce{^{235}_{92}U}
\usepackage[european,siunitx]{circuitikz} % schémas électriques
\pgfplotsset{compat=1.18}
\usepgfplotslibrary{fillbetween,groupplots} % aires entre courbes (calcul intégral)
\usepackage{enumitem}
\usepackage{fancyhdr}
% [most] charge toutes les bibliothèques tcolorbox (listings, minted, raster,
% documentation...) et gonfle inutilement la mémoire TeX sur un document long
% (~300 boîtes) : on ne charge que ce qui est réellement utilisé.
\usepackage[skins,breakable]{tcolorbox}
\usepackage{graphicx}
\usepackage{colortbl}
\usepackage{booktabs}
\usepackage{tabularx}
\usepackage{multirow}
\usepackage{diagbox} % case d'en-tête diagonale des tableaux à double entrée
% Colonnes extensibles centrée / gauche / droite (C, L, R), souvent utilisées par les modèles
\newcolumntype{C}{>{\centering\arraybackslash}X}
\newcolumntype{L}{>{\raggedright\arraybackslash}X}
\newcolumntype{R}{>{\raggedleft\arraybackslash}X}
\usepackage{array}
\usepackage{multicol}
\usepackage{siunitx}
% parse-numbers=false : accepte aussi \SI{2,5 \times 10^{-3}}{...}
\sisetup{locale=FR,output-decimal-marker={,},group-minimum-digits=4,parse-numbers=false}
\DeclareSIUnit\ohm{\ensuremath{\symup{\Omega}}} % U+2126 absent de la police
\DeclareSIUnit\division{div} % réglages d'oscilloscope (ms/div, V/div)
\DeclareSIUnit\tr{tr} % tours (vitesse de rotation, ex. \SI{1500}{\tr\per\minute})
\DeclareSIUnit\molar{M} % concentration molaire (ex. \SI{3}{\molar}, \SI{1}{\micro\molar})
\DeclareSIUnit\an{an} % année (le modèle confond parfois avec \year, un compteur TeX) — ex. \SI{5}{\kilo\meter\per\an}
\DeclareSIUnit\FCFA{FCFA} % franc CFA (ex. \SI{500}{\FCFA\per\kilo\gram})
\DeclareSIUnit\degreeBrix{\textdegree Brix} % teneur en sucre d'un sirop/jus (réfractométrie)
\DeclareSIUnit\month{mois} % le modèle utilise parfois \month, un compteur TeX, comme unité de durée
\DeclareSIUnit\microliter{\micro\liter} % le modèle écrit parfois \microliter en un seul token
\DeclareSIUnit\million{millions} % dénombrement (ex. \SI{15}{\million\per\mL} spermatozoïdes)
\usepackage{qrcode}
% \degree : commande fréquemment utilisée par les modèles pour un angle ou une
% température en dehors de \SI{}{} (non fournie par les paquets chargés).
\providecommand{\degree}{^{\circ}}
% \qed (fin de démonstration), utilisé par les modèles sans amsthm
\providecommand{\qed}{\hfill$\square$}
% \barre : double barre des valeurs interdites dans un tableau de variations (tkz-tab)
\providecommand{\barre}{\ensuremath{\|}}
% environnement proof (amsthm non chargé) utilisé par les modèles
\ifdefined\proof\else\newenvironment{proof}[1][Démonstration]{\par\noindent\textit{#1.}\ }{\qed\par}\fi
\usepackage{lastpage}
\usepackage{hyperref}
\hypersetup{hidelinks}

% ============================================================
% Palette « Pop » OnBuch+ — mêmes hex que la classe P (plus_ui.dart)
% ============================================================
\definecolor{poporange}{HTML}{F59321}
\definecolor{popdark}{HTML}{E07A0C}
\definecolor{popcream}{HTML}{FFF6E8}
\definecolor{popink}{HTML}{1C1714}
\definecolor{poppurple}{HTML}{7A5AE0}
\definecolor{popgreen}{HTML}{2E9E6A}
\definecolor{poppink}{HTML}{E0457B}
\definecolor{popblue}{HTML}{2F7DE1}
\definecolor{popgold}{HTML}{C98A12}
\definecolor{popmuted}{HTML}{8A7F76}
\definecolor{popline}{HTML}{EADFD2}
\definecolor{popgray}{HTML}{8A7F76}
\colorlet{popgrey}{popgray}
% Alias tolérants (le modèle nomme parfois les couleurs différemment)
\colorlet{popwhite}{white}
\colorlet{popblack}{popink}
\colorlet{popred}{poppink}
\colorlet{popyellow}{popgold}
% Teintes claires (fonds de tableaux/encadrés) — le modèle nomme parfois
% les couleurs avec un suffixe L (ex. popblueL) sans les avoir définies.
\colorlet{poporangeL}{poporange!20}
\colorlet{popdarkL}{popdark!20}
\colorlet{poppurpleL}{poppurple!20}
\colorlet{popgreenL}{popgreen!20}
\colorlet{poppinkL}{poppink!20}
\colorlet{popblueL}{popblue!20}
\colorlet{popgoldL}{popgold!20}
\colorlet{popredL}{popred!20}
\colorlet{popgrayL}{popgray!20}
% Teintes claires pour fonds de tableaux / zones
\colorlet{popblueL}{popblue!10!white}
\colorlet{poporangeL}{poporange!14!white}
\colorlet{popgreenL}{popgreen!12!white}
\colorlet{poppurpleL}{poppurple!12!white}
\colorlet{poppinkL}{poppink!12!white}
\colorlet{popgoldL}{popgold!14!white}

\pagecolor{popcream}

% ============================================================
% Typographie
% ============================================================
\defaultfontfeatures{Ligatures=TeX}
\setmainfont{PlusJakartaSans-Medium.ttf}[Path=../../../fonts/,BoldFont=PlusJakartaSans-Bold.ttf]
% Maths en sans-serif (Fira Math) avec chiffres et lettres droites de Plus
% Jakarta, pour que les formules se fondent dans le texte.
\usepackage[math-style=ISO,bold-style=ISO]{unicode-math}
\setmathfont{FiraMath-Regular.otf}[Path=../../../fonts/]
\setmathfont{PlusJakartaSans-Medium.ttf}[Path=../../../fonts/,range={up/num,up/{Latin,latin},\mathup}]
\usepackage{icomma} % virgule décimale sans espace en mode math
% Caractères mathématiques Unicode absents des polices de texte (Plus Jakarta,
% Archivo Black) : rendus par leur équivalent mathématique, en texte comme en
% formule — sinon ils s'affichent en carré vide (ex. « Arithmétique dans ℤ »).
\usepackage{newunicodechar}
\newunicodechar{ℝ}{\ensuremath{\mathbb{R}}}
\newunicodechar{ℤ}{\ensuremath{\mathbb{Z}}}
\newunicodechar{ℂ}{\ensuremath{\mathbb{C}}}
\newunicodechar{ℕ}{\ensuremath{\mathbb{N}}}
\newunicodechar{ℚ}{\ensuremath{\mathbb{Q}}}
\newunicodechar{𝔻}{\ensuremath{\mathbb{D}}}
\newunicodechar{α}{\ensuremath{\alpha}}
\newunicodechar{β}{\ensuremath{\beta}}
\newunicodechar{γ}{\ensuremath{\gamma}}
\newunicodechar{δ}{\ensuremath{\delta}}
\newunicodechar{ε}{\ensuremath{\varepsilon}}
\newunicodechar{θ}{\ensuremath{\theta}}
\newunicodechar{λ}{\ensuremath{\lambda}}
\newunicodechar{μ}{\ensuremath{\mu}}
\newunicodechar{σ}{\ensuremath{\sigma}}
\newunicodechar{τ}{\ensuremath{\tau}}
\newunicodechar{φ}{\ensuremath{\varphi}}
\newunicodechar{ψ}{\ensuremath{\psi}}
\newunicodechar{ω}{\ensuremath{\omega}}
\newunicodechar{Σ}{\ensuremath{\Sigma}}
% majuscules produites par \MakeUppercase dans les titres (α-aminés -> Α)
\newunicodechar{Α}{\ensuremath{\alpha}}
\newunicodechar{Β}{\ensuremath{\beta}}
\newunicodechar{Γ}{\ensuremath{\Gamma}}
\newunicodechar{Δ}{\ensuremath{\Delta}}
\newunicodechar{Ε}{\ensuremath{\varepsilon}}
\newunicodechar{Θ}{\ensuremath{\Theta}}
\newunicodechar{Λ}{\ensuremath{\Lambda}}
\newunicodechar{Μ}{\ensuremath{\mu}}
\newunicodechar{Τ}{\ensuremath{\tau}}
\newunicodechar{Φ}{\ensuremath{\Phi}}
\newunicodechar{Ψ}{\ensuremath{\Psi}}
\newunicodechar{Ω}{\ensuremath{\Omega}}
\newunicodechar{↦}{\ensuremath{\mapsto}}
\newunicodechar{∈}{\ensuremath{\in}}
\newunicodechar{∉}{\ensuremath{\notin}}
\newunicodechar{∧}{\ensuremath{\wedge}}
\newunicodechar{⇔}{\ensuremath{\Leftrightarrow}}
\newunicodechar{⟺}{\ensuremath{\Longleftrightarrow}}
\newunicodechar{⇒}{\ensuremath{\Rightarrow}}
\newunicodechar{⟹}{\ensuremath{\Longrightarrow}}
\newunicodechar{∘}{\ensuremath{\circ}}
\newunicodechar{∪}{\ensuremath{\cup}}
\newunicodechar{∩}{\ensuremath{\cap}}
\newunicodechar{≡}{\ensuremath{\equiv}}
\newunicodechar{⊂}{\ensuremath{\subset}}
\newunicodechar{⊥}{\ensuremath{\perp}}
\newunicodechar{∃}{\ensuremath{\exists}}
\newunicodechar{∀}{\ensuremath{\forall}}
\newunicodechar{∛}{\ensuremath{\sqrt[3]{\,}}}
\newunicodechar{∜}{\ensuremath{\sqrt[4]{\,}}}
\newunicodechar{⃗}{\ensuremath{{}^{\to}}}
\newunicodechar{ⁿ}{\ifmmode^{n}\else\textsuperscript{\ensuremath{n}}\fi}
\newunicodechar{ˣ}{\ifmmode^{x}\else\textsuperscript{\ensuremath{x}}\fi}
\newunicodechar{ᵇ}{\ifmmode^{b}\else\textsuperscript{\ensuremath{b}}\fi}
\newunicodechar{ʳ}{\ifmmode^{r}\else\textsuperscript{\ensuremath{r}}\fi}
\newunicodechar{ᵘ}{\ifmmode^{u}\else\textsuperscript{\ensuremath{u}}\fi}
\newunicodechar{ʸ}{\ifmmode^{y}\else\textsuperscript{\ensuremath{y}}\fi}
\newunicodechar{ᵃ}{\ifmmode^{a}\else\textsuperscript{\ensuremath{a}}\fi}
\newunicodechar{ᵉ}{\ifmmode^{e}\else\textsuperscript{\ensuremath{e}}\fi}
\newunicodechar{ᵏ}{\ifmmode^{k}\else\textsuperscript{\ensuremath{k}}\fi}
\newunicodechar{ᵖ}{\ifmmode^{p}\else\textsuperscript{\ensuremath{p}}\fi}
\newunicodechar{ᴺ}{\ifmmode^{N}\else\textsuperscript{\ensuremath{N}}\fi}
\newunicodechar{ᶜ}{\ifmmode^{c}\else\textsuperscript{\ensuremath{c}}\fi}
\newunicodechar{ʰ}{\ifmmode^{h}\else\textsuperscript{\ensuremath{h}}\fi}
\newunicodechar{ᵠ}{\ifmmode^{q}\else\textsuperscript{\ensuremath{q}}\fi}
\newunicodechar{ₙ}{\ifmmode_{n}\else\textsubscript{\ensuremath{n}}\fi}
\newunicodechar{ₐ}{\ifmmode_{a}\else\textsubscript{\ensuremath{a}}\fi}
\newunicodechar{ᵢ}{\ifmmode_{i}\else\textsubscript{\ensuremath{i}}\fi}
\newunicodechar{ᵣ}{\ifmmode_{r}\else\textsubscript{\ensuremath{r}}\fi}
\newunicodechar{ₖ}{\ifmmode_{k}\else\textsubscript{\ensuremath{k}}\fi}
\newunicodechar{ᵦ}{\ifmmode_{\beta}\else\textsubscript{\ensuremath{\beta}}\fi}
\newunicodechar{ₓ}{\ifmmode_{x}\else\textsubscript{\ensuremath{x}}\fi}
\newfontfamily\popdisplay{ArchivoBlack-Regular.ttf}[Path=../../../fonts/]
\newfontfamily\poplogo{SpaceGrotesk-Bold.ttf}[Path=../../../fonts/]
\newfontfamily\popsemi{PlusJakartaSans-SemiBold.ttf}[Path=../../../fonts/]
\newfontfamily\popxbold{PlusJakartaSans-ExtraBold.ttf}[Path=../../../fonts/]
\newfontfamily\popxboldls{PlusJakartaSans-ExtraBold.ttf}[Path=../../../fonts/,LetterSpace=3.0]

\setlist[enumerate]{leftmargin=1.7em,itemsep=5pt,topsep=4pt}
\setlist[itemize]{leftmargin=1.7em,itemsep=4pt,topsep=4pt,label={\color{poporange}\textbullet}}
\setlength{\parindent}{0pt}
\setlength{\parskip}{5pt}
\renewcommand{\arraystretch}{1.25}

% pgfplots : style Pop par défaut
\pgfplotsset{
  every axis/.append style={axis line style={popink,line width=1pt},tick style={popink},
    grid style={popline,line width=0.5pt},label style={font=\small},tick label style={font=\footnotesize}},
  popcurve/.style={poporange,line width=1.8pt,no markers,smooth},
}
% Même style disponible dans un \draw TikZ simple (hors pgfplots)
\tikzset{popcurve/.style={poporange,line width=1.8pt}}
\tikzset{popgrid/.style={popline,very thin}}
\colorlet{popcurve}{poporange} % le nom du style est parfois utilisé comme couleur

% ============================================================
% Pill / squiggle
% ============================================================
\newcommand{\poppill}[3]{%
  \tikz[baseline=-0.55ex]\node[draw=popink,line width=1.2pt,rounded corners=2.6pt,%
    fill=#2,inner xsep=6.5pt,inner ysep=3pt]{%
    \popxboldls\fontsize{7}{7}\selectfont\color{#3}\MakeUppercase{#1}};%
}
\newcommand{\popsquiggle}[1]{%
  \tikz[baseline=-0.4ex]{
    \draw[#1,line width=2.6pt,line cap=round]
      plot [smooth] coordinates {(0,0.09) (0.34,0.01) (0.68,0.21) (1.02,0.01) (1.36,0.10)};
  }%
}
% Mot-clé surligné (marqueur orange sous le texte)
\newcommand{\cle}[1]{{\popxbold\color{popdark}#1}}

% ============================================================
% En-tête / pied
% ============================================================
\pagestyle{fancy}
\fancyhf{}
\renewcommand{\headrulewidth}{0pt}
\renewcommand{\footrulewidth}{0pt}
\lhead{\raisebox{-0.15ex}{\poplogo\fontsize{13}{13}\selectfont\color{popink}OnBuch\color{poporange}+}}
\rhead{\poppill{\DOCMATIERE\ \textbullet\ \DOCNIVEAU\ \textbullet\ Leçon \DOCLECON}{white}{popink}}
\lfoot{\popxbold\fontsize{7.6}{7.6}\selectfont\color{popmuted}\MakeUppercase{Cours signé OnBuch+}\ \textbullet\ {\popsemi\DOCTITRE}}
\rfoot{\tikz[baseline=-0.6ex]\node[rounded corners=8.5pt,fill=popink,minimum height=17pt,minimum width=17pt,inner xsep=5pt,inner sep=0]{\popxbold\fontsize{8}{8}\selectfont\color{popcream}\thepage\,/\,\pageref*{LastPage}};}

% ============================================================
% Titres
% ============================================================
\newcommand{\chaptitre}[2]{%
  \begin{tcolorbox}[enhanced,colback=poporange,colframe=popink,boxrule=2.6pt,arc=11pt,
      drop shadow=popink,left=15pt,right=15pt,top=12pt,bottom=11pt,before skip=3pt,after skip=18pt]
    \poppill{#2}{popink}{popcream}\par\vspace{9pt}
    {\raggedright\hyphenpenalty=10000\exhyphenpenalty=10000\popdisplay\fontsize{22}{24}\selectfont\color{popcream}\MakeUppercase{#1}\par}
  \end{tcolorbox}
}
% Section numérotée : pastille orange avec le numéro + titre Archivo
\newcounter{popsec}
\newcommand{\coursec}[1]{%
  \stepcounter{popsec}\setcounter{popsub}{0}%
  \par\vspace{16pt}\noindent
  \begin{minipage}{\linewidth}
  \tikz[baseline=-0.7ex]\node[draw=popink,line width=1.6pt,fill=poporange,rounded corners=4pt,minimum size=22pt,inner sep=0]{\popdisplay\fontsize{12}{12}\selectfont\color{popcream}\thepopsec};%
  \hspace{8pt}{\popdisplay\fontsize{15}{17}\selectfont\color{popink}\MakeUppercase{#1}}\par
  \vspace{4pt}\noindent{\color{poporange}\rule{36pt}{3.6pt}}
  \end{minipage}\par\nopagebreak\vspace{8pt}
}
\newcounter{popsub}
\newcommand{\courssub}[1]{%
  \stepcounter{popsub}%
  \par\vspace{9pt}\noindent{\popxbold\fontsize{11.5}{13}\selectfont\color{popink}{\color{poporange}\thepopsec.\thepopsub}\hspace{6pt}#1}\par\nopagebreak\vspace{4pt}
}

% ============================================================
% Boîtes pédagogiques — même recette : fond blanc, bordure épaisse colorée,
% ombre dure encre, badge plein. Titre optionnel [#1].
% ============================================================
\tcbset{popbase/.style={breakable,enhanced,colback=white,boxrule=2.3pt,arc=9pt,
  shadow={3pt}{-3pt}{0pt}{popink},left=14pt,right=14pt,top=11pt,bottom=11pt,before skip=11pt,after skip=15pt,
  fontupper=\popsemi\fontsize{10.6}{14.5}\selectfont\color{popink},
  fontlower=\popsemi\fontsize{10.6}{14.5}\selectfont\color{popink}}}
% Coche dessinée (le glyphe ✓ n'existe pas dans les polices du document)
\newcommand{\popcheck}[1]{\tikz[baseline=-0.5ex]\draw[#1,line width=1.6pt,line cap=round,line join=round](-0.13,0.0)--(-0.04,-0.09)--(0.13,0.1);}
\providecommand{\coche}{\popcheck{popgreen}}
\providecommand{\resultat}[1]{\par\smallskip\noindent\fcolorbox{popgreen}{popgreenL}{\parbox{\dimexpr\linewidth-2\fboxsep-2\fboxrule\relax}{\textbf{Résultat.} #1}}\par\smallskip}
\newcommand{\popboxhead}[3]{\poppill{#1}{#2}{white}\ifx\relax#3\relax\else\hspace{6pt}{\popxbold\fontsize{10.6}{12}\selectfont\color{popink}#3}\fi\par\vspace{7pt}}

\newtcolorbox{aretenir}[1][]{popbase,colframe=popgreen,before upper={\popboxhead{À retenir}{popgreen}{#1}}}
\newtcolorbox{exemplebox}[1][]{popbase,colframe=poporange,before upper={\popboxhead{Exemple}{poporange}{#1}}}
\newtcolorbox{definition}[1][]{popbase,colframe=popblue,before upper={\popboxhead{Définition}{popblue}{#1}}}
\newtcolorbox{propriete}[1][]{popbase,colframe=poppurple,before upper={\popboxhead{Propriété}{poppurple}{#1}}}
\newtcolorbox{methode}[1][]{popbase,colframe=popgold,before upper={\popboxhead{Méthode}{popgold}{#1}}}
\newtcolorbox{attention}[1][]{popbase,colframe=poppink,before upper={\popboxhead{Attention}{poppink}{#1}}}
\newtcolorbox{experience}[1][]{popbase,colframe=popblue,colback=popblueL,before upper={\popboxhead{Expérience}{popblue}{#1}}}
% « Le savais-tu ? » : carte violette pleine (contexte, culture, Cameroun)
\newtcolorbox{savaistu}[1][]{popbase,colback=poppurple,colframe=popink,
  fontupper=\popsemi\fontsize{10.4}{14.2}\selectfont\color{white},
  before upper={\poppill{Le savais-tu\,?}{white}{poppurple}\ifx\relax#1\relax\else\hspace{6pt}{\popxbold\color{white}#1}\fi\par\vspace{7pt}}}
% Exercice résolu : énoncé en haut, \tcblower puis la solution sur fond crème
\newcounter{popexo}\newcounter{popexr}
\newtcolorbox{exoresolu}[1][]{popbase,colframe=poporange,colbacklower=popcream,
  segmentation style={popink,line width=1.2pt,dashed},
  before upper={\stepcounter{popexr}\popboxhead{Exercice résolu \thepopexr}{poporange}{#1}},
  before lower={\poppill{Solution}{popink}{popcream}\par\vspace{6pt}}}
% Exercice (non résolu en place) — la correction est dans la partie Corrigés
\newtcolorbox{exercice}[1][]{popbase,colframe=popink,
  before upper={\stepcounter{popexo}\popboxhead{Exercice \thepopexo}{popink}{#1}}}
% Corrigé
\newtcolorbox{corrige}[1][]{popbase,colframe=popgreen,colback=popgreenL,
  before upper={\popboxhead{Corrigé}{popgreen}{#1}}}
% Objectifs / prérequis / bilan
\newtcolorbox{objectifs}{popbase,colframe=popink,colback=poporangeL,
  before upper={\popboxhead{Objectifs de la leçon}{popink}{}}}
\newtcolorbox{prerequis}{popbase,colframe=popink,colback=popblueL,
  before upper={\popboxhead{Prérequis}{popblue}{}}}
\newtcolorbox{bilan}{popbase,colframe=popink,colback=white,boxrule=2.8pt,
  before upper={\popboxhead{Fiche bilan}{poporange}{L'essentiel à connaître pour l'examen}}}
% Figure encadrée avec légende
\newtcolorbox{popfigure}[1][]{enhanced,breakable=false,colback=white,colframe=popline,boxrule=1.4pt,arc=8pt,
  left=10pt,right=10pt,top=10pt,bottom=8pt,before skip=10pt,after skip=12pt,halign=center,
  lower separated=false,halign lower=center,
  fontlower=\popsemi\fontsize{9}{11.5}\selectfont\color{popmuted},
  #1}
\newcommand{\legende}[1]{\par\vspace{6pt}{\popsemi\fontsize{9}{11.5}\selectfont\color{popmuted}#1}}

% ============================================================
% Signature OnBuch+
% ============================================================
\newcommand{\poplogoinline}[1]{{\poplogo\fontsize{#1}{#1}\selectfont\color{popink}OnBuch\color{poporange}+}}
\newcommand{\signatureonbuch}{%
  \begin{tcolorbox}[enhanced,colback=popink,colframe=popink,boxrule=2.2pt,arc=10pt,
      drop shadow=poporange,left=14pt,right=14pt,top=10pt,bottom=10pt,before skip=0pt,after skip=0pt]
    \tikz[baseline=-0.7ex]\node[circle,fill=popgreen,draw=popcream,line width=1.3pt,minimum size=22pt,inner sep=0]{\popcheck{white}};%
    \hspace{9pt}\begin{minipage}[c]{0.62\linewidth}
      {\popxbold\fontsize{12}{13}\selectfont\color{popcream}Signé\ \textbullet\ {\poplogo OnBuch\color{poporange}+}}\par\vspace{2pt}
      {\raggedright\popsemi\fontsize{8}{10.5}\selectfont\color{popline}\MakeUppercase{Équipe pédagogique OnBuch+}\\\MakeUppercase{Conforme au programme officiel MINESEC}\par}
    \end{minipage}\hfill
    {\poplogo\fontsize{22}{22}\selectfont\color{popcream}OnBuch\color{poporange}+}
  \end{tcolorbox}%
}

% ============================================================
% Page de garde
% #1 titre · #2 matière · #3 niveau · #4 module · #5 n° leçon · #6 durée ·
% #7 description / objectifs (texte ou itemize)
% ============================================================
\newcommand{\pagegarde}[7]{%
  \thispagestyle{empty}
  \newgeometry{a4paper,margin=1.7cm,top=1.6cm,bottom=1.4cm}%
  \begin{tikzpicture}[remember picture,overlay]
    \fill[poporange!18] ([xshift=-3.2cm,yshift=-3.6cm]current page.north east) circle (4.6cm);
    \fill[poppurple!14] ([xshift=2.4cm,yshift=7.6cm]current page.south west) circle (3.8cm);
  \end{tikzpicture}%
  {\poplogo\fontsize{30}{30}\selectfont\color{popink}OnBuch\color{poporange}+}\hfill
  \raisebox{0.6ex}{\poppill{Cours signé \textbullet\ Premium}{popink}{popcream}}\par
  \vspace{1pt}\popsquiggle{poppurple}\par
  \vspace{8pt}
  \begin{tcolorbox}[enhanced,colback=poporange,colframe=popink,boxrule=2.8pt,arc=13pt,
      drop shadow=popink,left=18pt,right=18pt,top=12pt,bottom=13pt,after skip=10pt]
    \poppill{#2 \textbullet\ #3}{popink}{popcream}\hspace{5pt}\poppill{#4}{white}{popink}\par\vspace{8pt}
    {\popdisplay\fontsize{14}{14}\selectfont\color{popink}LEÇON #5}\par\vspace{4pt}
    {\raggedright\hyphenpenalty=10000\exhyphenpenalty=10000\popdisplay\fontsize{25}{27}\selectfont\color{popcream}\MakeUppercase{#1}\par}
    \vspace{8pt}
    {\popxbold\fontsize{9.5}{9.5}\selectfont\color{popink}\MakeUppercase{Durée conseillée : #6}}
  \end{tcolorbox}
  \begin{tcolorbox}[enhanced,colback=white,colframe=popink,boxrule=2.2pt,arc=10pt,
      drop shadow=popink,left=16pt,right=16pt,top=10pt,bottom=10pt,after skip=8pt]
    \poppill{Dans cette leçon}{poppurple}{white}\par\vspace{5pt}
    \popsemi\fontsize{9.3}{12.6}\selectfont\color{popink}#7
  \end{tcolorbox}
  \noindent\begin{minipage}[t]{0.487\linewidth}
    \begin{tcolorbox}[enhanced,colback=popgreen,colframe=popink,boxrule=2.2pt,arc=10pt,
        drop shadow=popink,left=11pt,right=11pt,top=10pt,bottom=10pt,height=3.1cm,valign=center]
      \tcbox[colback=white,colframe=popink,boxrule=1.3pt,arc=5pt,boxsep=0pt,nobeforeafter,
        left=3pt,right=3pt,top=3pt,bottom=3pt,valign=center]{\qrcode[height=1.9cm]{https://play.google.com/store/apps/details?id=com.luvvix.onbuch&hl=fr}}%
      \hspace{8pt}\begin{minipage}[c]{0.5\linewidth}
        {\popdisplay\fontsize{11}{12.5}\selectfont\color{white}\MakeUppercase{Télécharge OnBuch}}\par\vspace{4pt}
        {\raggedright\popsemi\fontsize{8.4}{11}\selectfont\color{white}Gratuit \textbullet\ Google Play\\Tuteur IA, annales, concours, résultats\par}
      \end{minipage}
    \end{tcolorbox}
  \end{minipage}\hfill
  \begin{minipage}[t]{0.487\linewidth}
    \begin{tcolorbox}[enhanced,colback=poppurple,colframe=popink,boxrule=2.2pt,arc=10pt,
        drop shadow=popink,left=11pt,right=11pt,top=10pt,bottom=10pt,height=3.1cm,valign=center]
      \tikz[baseline=-0.7ex]\node[circle,fill=white,draw=popink,line width=1.3pt,minimum size=1.6cm,inner sep=0]{\popdisplay\fontsize{24}{24}\selectfont\color{poppurple}+};%
      \hspace{8pt}\begin{minipage}[c]{0.6\linewidth}
        {\popdisplay\fontsize{11}{12.5}\selectfont\color{white}\MakeUppercase{Passe à OnBuch+}}\par\vspace{4pt}
        {\raggedright\popsemi\fontsize{8.4}{11}\selectfont\color{white}Tous les cours signés, Tuteur IA illimité, fiches PDF — onglet OnBuch+ de l'app\par}
      \end{minipage}
    \end{tcolorbox}
  \end{minipage}
  \vspace{8pt}
  \popcontenu
  \vfill
  \signatureonbuch
  \restoregeometry
  \newpage
}

% Rangée « Ce cours contient » (page de garde)
\newcommand{\poptile}[3]{% #1 couleur #2 gros texte #3 légende
  \begin{tcolorbox}[enhanced,colback=#1,colframe=popink,boxrule=2pt,arc=9pt,shadow={2.5pt}{-2.5pt}{0pt}{popink},
      left=6pt,right=6pt,top=9pt,bottom=9pt,halign=center,valign=center,height=1.75cm,width=\linewidth,nobeforeafter]
    {\popdisplay\fontsize{12.5}{13}\selectfont\color{white}\MakeUppercase{#2}}\par\vspace{3pt}
    {\centering\popsemi\fontsize{7.8}{9.5}\selectfont\color{white}#3\par}
  \end{tcolorbox}}
\newcommand{\popcontenu}{%
  \par\noindent{\popxboldls\fontsize{7.5}{7.5}\selectfont\color{popmuted}CE COURS SIGNÉ CONTIENT}\par\vspace{7pt}
  \noindent\begin{minipage}[t]{0.235\linewidth}\poptile{popblue}{Cours}{complet, illustré, pas à pas}\end{minipage}\hfill
  \begin{minipage}[t]{0.235\linewidth}\poptile{poporange}{Méthodes}{et exercices résolus}\end{minipage}\hfill
  \begin{minipage}[t]{0.235\linewidth}\poptile{poppink}{Exercices}{type examen + corrigés}\end{minipage}\hfill
  \begin{minipage}[t]{0.235\linewidth}\poptile{popgreen}{Fiche bilan}{l'essentiel à retenir}\end{minipage}\par}

% Sommaire
\newtcolorbox{sommaire}{enhanced,colback=white,colframe=popink,boxrule=2.2pt,arc=10pt,
  drop shadow=popink,left=18pt,right=18pt,top=13pt,bottom=13pt,before skip=6pt,after skip=20pt,
  before upper={\poppill{Sommaire}{poppurple}{white}\par\vspace{9pt}\popsemi\fontsize{10.6}{14.5}\selectfont\color{popink}}}

% Signature de fin de cours — poussée tout en bas de la dernière page
\newcommand{\findecours}{%
  \par\vfill\nopagebreak
  \begin{center}\popsquiggle{poporange}\end{center}\vspace{4pt}
  \signatureonbuch
}
\AtBeginDocument{\def\llbracket{\mathopen{[\mkern-3.5mu[}}\def\rrbracket{\mathclose{]\mkern-3.5mu]}}} % intervalles d'entiers (glyphe ⟦ absent de la police)
% Glyphes absents de Fira Math : redéfinis à partir de glyphes disponibles
\AtBeginDocument{%
  \def\setminus{\mathbin{\backslash}}\let\smallsetminus\setminus
  \def\vdots{\mathord{\vbox{\baselineskip3.2pt\lineskiplimit0pt\kern4pt\hbox{.}\hbox{.}\hbox{.}}}}%
  \def\ddots{\mathinner{\mkern1mu\raise6.5pt\hbox{.}\mkern2mu\raise3.6pt\hbox{.}\mkern2mu\raise0.7pt\hbox{.}\mkern1mu}}%
  \def\triangle{\mathord{\scalebox{0.9}{$\symup{\Delta}$}}}%
  \def\nabla{\mathord{\raisebox{\depth}{\scalebox{1}[-1]{$\symup{\Delta}$}}}}%
  \def\checkmark{\popcheck{popdark}}%
  \def\star{\mathbin{\ast}}\def\bigstar{\mathord{\ast}}%
  \def\diamond{\mathbin{\scalebox{0.6}{$\Diamond$}}}%
  \def\bigcup{\mathop{\mathchoice{\raisebox{-0.3em}{\scalebox{1.7}{$\cup$}}}{\scalebox{1.25}{$\cup$}}{\cup}{\cup}}\displaylimits}%
  \def\bigcap{\mathop{\mathchoice{\raisebox{-0.3em}{\scalebox{1.7}{$\cap$}}}{\scalebox{1.25}{$\cap$}}{\cap}{\cap}}\displaylimits}%
  \def\lesseqqgtr{\mathrel{\lessgtr}}\def\bigoplus{\mathop{\oplus}\displaylimits}%
}
\providecommand{\ssi}{si et seulement si} % abréviation fréquente
\AtBeginDocument{\providecommand{\wideparen}{\overparen}} % arc de cercle

%% FIGFIX-BEGIN v5 — correctifs globaux des figures (docs/onbuch-generation/tools/figfix)
\usepackage{adjustbox}\usepackage{etoolbox}\tcbuselibrary{raster}
% toute figure TikZ/pgfplots plus large que la ligne est réduite à la largeur disponible
\BeforeBeginEnvironment{tikzpicture}{\begin{adjustbox}{max width=\linewidth}}
\AfterEndEnvironment{tikzpicture}{\end{adjustbox}}
% légendes pgfplots lisibles par-dessus les courbes
\pgfplotsset{every axis/.append style={legend style={fill=white,fill opacity=0.92,text opacity=1,draw=black!15,inner sep=3pt}}}
% étiquettes d'arêtes (syntaxe edge["..."]) avec fond blanc
\tikzset{every edge quotes/.append style={fill=white,inner sep=1.2pt}}
%% FIGFIX-END
\begin{document}

\pagegarde{La Deuxième Guerre mondiale}{Histoire}{Tle Tech}{Histoire – classe de Terminale technique}{N01}{4 séances (fiche de progression harmonisée)}{Cette leçon analyse la Seconde Guerre mondiale comme un conflit total et mondial, en mettant l'accent sur le basculement géopolitique de 1945 (naissance de l'ONU, guerre froide) et sur le rôle décisif mais souvent méconnu de l'Afrique et du Cameroun, à travers l'étude critique de sources variées (archives, chiffres, cartes, témoignages).
\vspace{4pt}
\begin{multicols}{2}\small
\begin{itemize}
\item À la fin de cette leçon, je sais expliquer les causes structurelles et conjoncturelles de la Seconde Guerre mondiale et identifier les phases majeures du conflit (1939-1945).
\item Je sais caractériser la nature totalitaire et génocidaire du IIIe Reich (Shoah, guerre d'anéantissement) et la mobilisation totale des sociétés (économie, propagande, résistance).
\item Je sais analyser les conséquences démographiques, territoriales, politiques et morales du conflit et expliquer la création de l'ONU et le nouvel ordre bipolaire.
\item Je sais exploiter des documents statistiques et cartographiques pour montrer l'apport militaire, économique et stratégique de l'Afrique et du Cameroun à la victoire alliée.
\item Je sais rédiger une réponse argumentée en mobilisant des connaissances précises (dates, chiffres, acteurs) et en justifiant ma démarche historique (critique de sources, mise en perspective).
\item Je sais réaliser un croquis commenté ou un diagramme à bandes pour représenter la contribution camerounaise (effectifs, productions, ralliement).
\end{itemize}\end{multicols}}

\begin{sommaire}
\begin{multicols}{2}
\begin{enumerate}
\item 1. Genèse et extension d'un conflit mondial (1939-1941)
\item 2. La violence paroxystique : génocide, guerre d'anéantissement, société mobilisée
\item 3. L'Afrique et le Cameroun dans la tourmente : enjeux, contributions, mutations
\item 4. 1945 : Un monde nouveau - Conséquences et naissance de l'ONU
\item 5. Héritages et mémoires : Décolonisation, Droits de l'Homme, Construction européenne
\item Activité d'intégration
\item Méthodes et exercices résolus
\item Exercices
\item Corrigés des exercices
\item Fiche bilan
\end{enumerate}
\end{multicols}
\end{sommaire}

\begin{objectifs}
\begin{itemize}
\item À la fin de cette leçon, je sais expliquer les causes structurelles et conjoncturelles de la Seconde Guerre mondiale et identifier les phases majeures du conflit (1939-1945).
\item Je sais caractériser la nature totalitaire et génocidaire du IIIe Reich (Shoah, guerre d'anéantissement) et la mobilisation totale des sociétés (économie, propagande, résistance).
\item Je sais analyser les conséquences démographiques, territoriales, politiques et morales du conflit et expliquer la création de l'ONU et le nouvel ordre bipolaire.
\item Je sais exploiter des documents statistiques et cartographiques pour montrer l'apport militaire, économique et stratégique de l'Afrique et du Cameroun à la victoire alliée.
\item Je sais rédiger une réponse argumentée en mobilisant des connaissances précises (dates, chiffres, acteurs) et en justifiant ma démarche historique (critique de sources, mise en perspective).
\item Je sais réaliser un croquis commenté ou un diagramme à bandes pour représenter la contribution camerounaise (effectifs, productions, ralliement).
\item Je sais établir un lien entre l'expérience de la guerre (vétérans, idées des Droits de l'Homme, affaiblissement colonial) et l'émergence des nationalismes africains d'après-guerre.
\end{itemize}
\end{objectifs}

\begin{prerequis}
\begin{itemize}
\item La Première Guerre mondiale (causes, issue, traité de Versailles) et la SDN (échec) - Programme de 4ème / 1ère.
\item La montée des totalitarismes dans l'entre-deux-guerres (fascisme italien, nazisme, stalinisme) - Programme de 1ère.
\item La politique d'apaisement (Munich 1938) et le pacte germano-soviétique (1939).
\item Notions de base de cartographie (légende, échelle, projection) et de lecture de graphiques statistiques.
\item Contexte colonial camerounais sous mandat français (administration, économie de plantation, travail forcé) avant 1939.
\end{itemize}
\end{prerequis}

\newpage

\coursec{Genèse et extension d'un conflit mondial (1939-1941)}

\noindent\textbf{Situation d'accroche.} En août 1940, alors que la France métropolitaine vient de signer l'armistice et que le régime de Vichy collabore avec l'occupant nazi, un événement majeur se produit en Afrique équatoriale française (AEF). Le gouverneur du Tchad, Félix Éboué, rallie la France libre du général de Gaulle. Quelques jours plus tard, le 27 août 1940, le gouverneur du Cameroun, Richard Brunot, proclame à son tour le ralliement de la colonie. Douala devient, pour un temps, la « capitale de la France combattante » en Afrique. Des milliers de Camerounais, souvent recrutés de force ou par pression administrative, vont bientôt embarquer pour des théâtres d'opérations lointains : Birmanie, Italie, Provence. Comment un conflit né en Europe a-t-il pu embraser le monde entier jusqu'aux rives du Wouri ? Pourquoi l'ordre issu de Versailles a-t-il cédé la place à une guerre d'anéantissement sans précédent ?

\noindent\textbf{Problématique.} \emph{Comment l'agression hitlérienne, facilitée par l'échec de l'ordre de Versailles, a-t-elle transformé un conflit européen en une guerre mondiale totale en moins de deux ans (1939-1941) ?}

\courssub{L'échec de l'ordre de Versailles et la stratégie hitlérienne}

\noindent Le traité de Versailles (1919) visait à empêcher toute résurgence de la puissance allemande : démilitarisation de la Rhénanie, limitation de l'armée à 100 000 hommes, pertes territoriales, réparations financières écrasantes. Mais cet ordre repose sur une illusion : celle d'une Allemagne durablement affaiblie et d'une Société des Nations (SDN) capable de faire respecter le droit international. Dès 1933, Adolf Hitler, arrivé légalement au pouvoir, engage une politique de rupture systématique, fondée sur deux piliers idéologiques exposés dans \emph{Mein Kampf} (1925) : le \cle{Lebensraum} (espace vital) à conquérir à l'Est aux dépens des « sous-hommes » slaves, et le rassemblement de tous les \cle{Volksdeutsche} (Allemands ethniques) dans un \cle{Großdeutschland} (Grand Reich).

\begin{methode}[Analyser les étapes de la rupture de l'ordre de Versailles]
\textbf{Étape 1 :} Identifier la clause du traité violée (militaire, territoriale, diplomatique).
\textbf{Étape 2 :} Dater l'action hitlérienne et nommer l'opération ou l'événement.
\textbf{Étape 3 :} Analyser la réaction des puissances démocratiques (France, Royaume-Uni) : fermeté, protestation diplomatique, ou \cle{politique d'apaisement} (concessions pour éviter la guerre).
\textbf{Étape 4 :} Évaluer les conséquences : renforcement de la position stratégique d'Hitler, perte de crédibilité de la SDN.
\end{methode}

\begin{center}
\begin{tabularx}{\linewidth}{|l|X|X|X|}
\hline
\rowcolor{popblueL}
\textbf{Date} & \textbf{Action hitlérienne (Violation de Versailles)} & \textbf{Réaction France / UK} & \textbf{Conséquence stratégique} \\
\hline
1935 & Réarmement officiel (service militaire obligatoire), création de la Luftwaffe & Protestations verbales, conférence de Stresa (vaine) & Allemagne récupère une puissance militaire moderne. \\
\hline
Mars 1936 & Remilitarisation de la Rhénanie (zone démilitarisée) & France : refus d'intervenir sans appui britannique ; UK : « ils rentrent dans leur jardin » & Barrière naturelle perdue ; Belgique et France exposées ; prestige d'Hitler maximal. \\
\hline
Mars 1938 & \cle{Anschluss} (annexion de l'Autriche) & Apaisement : pas d'opposition militaire & Reich élargi de 7 millions d'habitants, ressources minières, flanc sud sécurisé. \\
\hline
Sept. 1938 & \cle{Accords de Munich} : cession des Sudètes (Tchécoslovaquie) & \textbf{Apaisement total} (Daladier, Chamberlain) : « Paix pour notre temps » & Tchécoslovaquie démembrée ; URSS exclue des négociations ; Hitler méprise la faiblesse adverse. \\
\hline
Mars 1939 & Occupation du reste de la Tchécoslovaquie (protectorat Bohême-Moravie) & Garanties verbales à la Pologne et Roumanie (tardives) & Preuve que Hitler ne veut pas seulement réviser Versailles mais dominer l'Europe. \\
\hline
Août 1939 & \cle{Pacte germano-soviétique} (Ribbentrop-Molotov) : clause secrète de partage de la Pologne & Sidération : l'idéologie anti-communiste d'Hitler s'efface devant l'opportunisme. & Hitler évite la guerre sur deux fronts ; feu vert pour l'invasion de la Pologne. \\
\hline
\end{tabularx}
\end{center}
\legende{Tableau 1 : La stratégie de rupture d'Hitler (1935-1939) et l'échec de l'apaisement.}

\begin{attention}[Piège classique : Confondre Apaisement et Collaboration]
L'\cle{apaisement} (1936-1938) est une politique diplomatique des démocraties (France, UK) consistant à céder aux exigences d'Hitler pour préserver la paix. Elle relève de la naïveté, de la peur de la guerre et de l'anticommunisme. La \cle{collaboration} (1940-1944) est la politique de l'État français de Vichy \emph{après} la défaite, consistant à coopérer activement avec l'occupant nazi. Ne les confondez jamais dans une copie.
\end{attention}

\courssub{Le déclenchement : agression de la Pologne et « drôle de guerre »}

\noindent Le 1er septembre 1939, à 4h45, le cuirassé \emph{Schleswig-Holstein} ouvre le feu sur le dépôt de munitions de Westerplatte (Dantzig) : c'est le début de l'invasion de la Pologne. Hitler lance la \cle{Blitzkrieg} (« guerre éclair »), une doctrine combinant concentration massive de chars (\cle{Panzerdivisionen}), appui aérien rapproché (\cle{Stuka} Ju 87) et liaisons radio pour percer le front, encercler l'ennemi et désorganiser son commandement.

\begin{definition}[Blitzkrieg (Guerre éclair)]
Doctrine militaire allemande reposant sur la \textbf{manœuvre} plutôt que sur l'usure. Elle combine : (1) une percée concentrée sur un point étroit (\emph{Schwerpunkt}) par des divisions blindées et motorisées ; (2) un appui aérien tactique (bombardement en piqué) agissant comme de l'artillerie volante ; (3) des transmissions radio généralisées pour une coordination en temps réel. Objectif : paralyser le centre de décision adverse en quelques jours, éviter la guerre de position.
\end{definition}

\noindent La Pologne, mal équipée, mal commandée, attaquée à l'ouest par l'Allemagne et à l'est par l'URSS (17 sept., clause secrète du pacte), s'effondre en trois semaines. La France et le Royaume-Uni déclarent la guerre à l'Allemagne le 3 septembre, mais n'engagent aucune offensive majeure sur le front occidental. C'est la \cle{drôle de guerre} (sept. 1939 - mai 1940) : huit mois d'attente, de propagande et de réarmement accéléré, pendant lesquels Hitler achève la conquête de la Pologne et se retourne vers le Nord (Danemark, Norvège, avril 1940) pour sécuriser ses approvisionnements en fer suédois.

\begin{popfigure}
\begin{tikzpicture}[scale=0.85, every node/.style={font=\small}]
% Fond Europe simplifié
\draw[popline, thick] (0,0) rectangle (14,10);
% Mers
\fill[popblueL!30] (0,0) rectangle (14,2.5); % Mer du Nord / Baltique
\fill[popblueL!30] (0,7.5) rectangle (14,10); % Méditerranée (partiel)
% Pays gros blocs
\fill[poporangeL!50] (2,2.5) rectangle (6,7.5); % Allemagne + Europe centrale
\fill[popgreenL!50] (6,2.5) rectangle (9,5.5); % Pologne
\fill[poppurpleL!50] (9,2.5) rectangle (12,7.5); % URSS (partie ouest)
\fill[popblueL!50] (0.5,3.5) rectangle (2,6.5); % France
\fill[poppinkL!50] (0.5,1.5) rectangle (2,3.5); % UK (au nord de la France)

% Légendes pays
\node[align=center] at (4,5) {\textbf{REICH}\\(Allemagne)};
\node[align=center] at (7.5,4) {\textbf{POLOGNE}};
\node[align=center] at (10.5,5) {\textbf{URSS}};
\node[align=center] at (1.25,5) {\textbf{FRANCE}};
\node[align=center] at (1.25,2.5) {\textbf{UK}};

% Flèches Blitzkrieg Pologne (Sept 1939)
\draw[poporange, very thick, -{Stealth[length=3mm]}] (4,5.5) -- (6.5,4.5) node[midway, above, sloped, fill=white, inner sep=1pt] {1. Sept 1939};
\draw[poporange, very thick, -{Stealth[length=3mm]}] (10,4) -- (8,4) node[midway, above, sloped, fill=white, inner sep=1pt] {17 sept};

% Flèches Ouest (Mai 1940)
\draw[poporange, very thick, -{Stealth[length=3mm]}, dashed] (4,5) .. controls (3,6) and (2,6) .. (1.5,5.5) node[midway, left, fill=white, inner sep=1pt] {Mai 1940};
\draw[poporange, very thick, -{Stealth[length=3mm]}, dashed] (4,4.5) .. controls (3,3.5) and (2,3.5) .. (1.5,4) node[midway, left, fill=white, inner sep=1pt] {Plan Sichelschnitt};

% Ligne Maginot
\draw[popblue, thick, decorate, decoration={snake, amplitude=1pt, segment length=4pt}] (1.5,3.5) -- (1.5,6.5) node[midway, left=2mm] {Ligne Maginot};

% Zone Vichy / Libre (Juin 1940)
\draw[popred, thick, dashed] (1.25,5) -- (3.5,5) -- (3.5,3.5) -- (1.25,3.5) -- cycle;
\node[popred, align=center, font=\tiny] at (2.3,4.2) {Zone Occupée};
\draw[popgreen, thick, dashed] (1.25,3.5) -- (0.5,3.5) -- (0.5,6.5) -- (1.25,6.5) -- cycle;
\node[popgreen, align=center, font=\tiny] at (0.8,5) {Zone Libre\\(Vichy)};

% Légende intégrée (dans la mer du Nord)
\begin{scope}[xshift=0.5cm, yshift=0.5cm]
\draw[poporange, very thick, -{Stealth[length=2mm]}] (0,0) -- (1.5,0) node[right] {Attaque allemande 1939-40};
\draw[poporange, very thick, dashed, -{Stealth[length=2mm]}] (0,-0.6) -- (1.5,-0.6) node[right] {Plan Jaune / Sichelschnitt 1940};
\draw[popred, thick, dashed] (0,-1.2) rectangle (0.5,-0.7) node[right] {Zone occupée};
\draw[popgreen, thick, dashed] (0,-1.8) rectangle (0.5,-1.3) node[right] {Zone libre (Vichy)};
\draw[popblue, thick, decorate, decoration={snake, amplitude=1pt, segment length=4pt}] (0,-2.4) -- (0.5,-2.4) node[right] {Ligne Maginot};
\end{scope}
\end{tikzpicture}
\legende{Carte schématique : L'expansion allemande en Europe (1939-1940). Flèches pleines : campagne de Pologne (sept. 1939) et offensive à l'Ouest (mai 1940). Flèche est : invasion soviétique (17 sept. 1939). Zones hachurées : découpage de la France après l'armistice (juin 1940).}
\end{popfigure}

\courssub{L'effondrement français (mai-juin 1940) : Armistice, Vichy et France Libre}

\noindent Le 10 mai 1940, l'offensive allemande (\cle{Plan Sichelschnitt}) contourne la ligne Maginot par les Ardennes (jugées infranchissables par l'état-major français). La percée de Sedan (13-15 mai) ouvre la route de la Manche. L'armée française, mal commandée, mal équipée en chars et radios, sans réserve stratégique, est coupée en deux. Dunkerque (opération Dynamo) permet l'évacuation de 338 000 soldats (dont 120 000 Français), mais l'essentiel du matériel est perdu.

\noindent Le 10 juin, le gouvernement français fuit Paris pour Bordeaux. Le maréchal Pétain, héros de Verdun, devient président du Conseil le 16 juin et demande l'armistice. Signé le \textbf{22 juin 1940} à Rethondes (wagon de 1918), il entre en vigueur le 25 juin.

\begin{attention}[Distinction cruciale : Armistice vs Capitulation]
\begin{itemize}
\item \textbf{Armistice (22 juin 1940) :} Accord militaire de \emph{cessez-le-feu}. La France n'est pas juridiquement vaincue (pas de traité de paix). Elle conserve sa flotte, son empire colonial, une armée d'armistice (100 000 hommes) et l'administration de la « zone libre ». C'est un piège : Hitler évite de gérer l'empire et la marine française.
\item \textbf{Capitulation (7-8 mai 1945) :} Acte juridique de \emph{reddition sans condition} (acte de Berlin). L'État allemand cesse d'exister en tant que sujet de droit international.
\end{itemize}
\emph{Au Bac : Si le sujet porte sur 1940, écrivez « armistice ». Si le sujet porte sur 1945, écrivez « capitulation ».}
\end{attention}

\noindent La France est coupée en deux : une \cle{zone occupée} (Nord, Ouest, côte atlantique) sous administration militaire allemande, et une \cle{zone libre} (Sud-Est, empire) administrée par le nouveau régime de \cle{Vichy} (capitale : Vichy). Le 10 juillet 1940, l'Assemblée nationale vote les pleins pouvoirs à Pétain (569 voix pour, 80 contre, 17 abstentions). La IIIe République meurt ; l'\cle{État français} naît, avec sa devise « Travail, Famille, Patrie » et sa \cle{Révolution nationale} (autoritarisme, corporatisme, antisémitisme d'État : Statut des Juifs, oct. 1940).

\noindent Face à la soumission, le général de Gaulle, sous-secrétaire d'État à la Guerre, refuse l'armistice. De Londres, le \textbf{18 juin 1940}, il lance son \cle{Appel du 18 juin} à la BBC : « La France a perdu une bataille, mais la France n'a pas perdu la guerre. » C'est l'acte de naissance de la \cle{France Libre} (puis France Combattante). Elle est d'abord une minorité : quelques milliers de volontaires, des territoires ralliés (Nouvelle-Calédonie, Polynésie, puis AEF grâce à Éboué et Brunot). La légitimité se construira dans la durée, par la lutte et la reconnaissance alliée.

\begin{savaistu}[Le Cameroun, premier territoire de l'AEF rallié]
\emph{Août 1940.} Alors que le Gabon reste fidèle à Vichy jusqu'en novembre 1940 (libéré par la force par Leclerc), le \textbf{Cameroun français} bascule dès le \textbf{27 août 1940}. Le gouverneur \textbf{Richard Brunot}, pressé par les milieux d'affaires (planteurs, commerçants) hostiles à Vichy et par l'opinion locale, proclame le ralliement à Douala. C'est une décision \textbf{stratégique majeure} : le Cameroun offre à de Gaulle un territoire vaste, des revenus douaniers (port de Douala), des ressources agricoles (caoutchouc, cacao, bananes) et une base arrière pour la conquête du Gabon. Des centaines de Camerounais s'engagent dans les Forces françaises libres (FFL) dès 1940, bien avant le recrutement de masse de 1943-44.
\end{savaistu}

\courssub{L'extension mondiale : de la bataille d'Angleterre à Pearl Harbor (1940-1941)}

\noindent Hitler, maître de l'Europe continentale (sauf UK, URSS, Suisse, Suède, Espagne, Portugal), tente de forcer le Royaume-Uni à la paix.

\begin{itemize}
\item \textbf{Bataille d'Angleterre (juillet-oct. 1940) :} La Luftwaffe (Göring) tente de détruire la RAF (Royal Air Force) pour préparer l'opération \cle{Seelöwe} (débarquement). Grâce au radar (\cle{Chain Home}), au décryptage \cle{Enigma} (Ultra) et au courage des pilotes (dont polonais, tchèques, néo-zélandais), la RAF résiste. Hitler renonce au débarquement (oct. 1940) et se tourne vers l'Est. Premier échec stratégique majeur du IIIe Reich.
\item \textbf{Opération \cle{Barbarossa} (22 juin 1941) :} Rupture du pacte germano-soviétique. 3,8 millions de soldats, 3 groupes d'armées (Nord/Leningrad, Centre/Moscou, Sud/Kiev). Objectif : victoire éclair avant l'hiver (guerre de mouvement). Idéologie : \cle{guerre d'anéantissement} (Vernichtungskrieg) contre le « judéo-bolchévisme » : commissaires politiques fusillés (ordre des commissaires), Juifs massacrés par les \cle{Einsatzgruppen} derrière le front, prisonniers de guerre soviétiques laissés mourir de faim (3,3 millions de morts). L'URSS résiste (défense de Moscou, déc. 1941), la guerre s'enlise.
\item \textbf{Pearl Harbor (7 déc. 1941) :} Le Japon, en guerre en Chine depuis 1937, sous embargo pétrolier US, frappe la flotte américaine à Hawaï. Hitler, par solidarité d'Axe (Pacte tripartite sept. 1940), déclare la guerre aux USA (11 déc. 1941). Erreur fatale : le conflit devient \textbf{véritablement mondial}, opposant l'Axe (Allemagne, Italie, Japon, satellites) à la \cle{Grande Alliance} (UK, USA, URSS) : les « \cle{Grand Trois} » (Churchill, Roosevelt, Staline).
\end{itemize}

\begin{definition}[Guerre totale]
Conflit dans lequel les belligérants mobilisent \textbf{l'intégralité de leurs ressources} : humaines (conscription massive, travail forcé, femmes dans l'industrie), économiques (réorientation totale de l'appareil productif vers l'effort de guerre, planification d'État), scientifiques (projet Manhattan, radar, fusées V2, pénicilline), morales (propagande, censure, haine de l'ennemi). \textbf{Effacement de la distinction front/arrière} : les civils deviennent cibles (bombardements stratégiques, blocus, génocides). La guerre devient le mode de fonctionnement normal de la société.
\end{definition}

\begin{popfigure}
\begin{tikzpicture}
\begin{axis}[
    ybar,
    bar width=12pt,
    width=\linewidth,
    height=8cm,
    enlarge x limits=0.25,
    ylabel={Production / PIB (Indices relatifs 1937=100)},
    symbolic x coords={Acier, Charbon, PIB, Aluminium, Camions/Avions},
    xtick=data,
    nodes near coords,
    nodes near coords align={vertical},
    every node near coord/.append style={font=\tiny, rotate=90, anchor=west, yshift=2pt},
    legend style={at={(0.5,-0.2)}, anchor=north, legend columns=-1, font=\small},
    ymin=0, ymax=450,
    ymajorgrids=true,
    grid style=dashed,
    axis x line*=bottom,
    axis y line*=left,
    tick label style={font=\small},
    legend cell align={left}
]
\addplot[popblue!80!popdark, fill=popblue!60] coordinates {(Acier, 380) (Charbon, 250) (PIB, 220) (Aluminium, 400) (Camions/Avions, 350)};
\addplot[poporange!80!popdark, fill=poporange!60] coordinates {(Acier, 120) (Charbon, 110) (PIB, 130) (Aluminium, 100) (Camions/Avions, 110)};
\legend{\textbf{Alliés (UK+USA+URSS 1941)}, \textbf{Axe (Allemagne+Italie+Japon 1941)}}
\end{axis}
\end{tikzpicture}
\legende{Diagramme en barres : Disparité économique et industrielle entre les Alliés (Grande Alliance) et l'Axe en 1941 (indices relatifs, base 100 = niveau 1937). L'entrée en guerre des USA et la tenue de l'URSS créent un déséquilibre de production insurmontable pour l'Axe. Sources : Harrison (1998), \emph{The Economics of World War II}.}
\end{popfigure}

\begin{methode}[Analyser une carte de progression militaire (type Bac)]
\textbf{1. Lire le titre, la légende, l'échelle et l'orientation.} (Quelle zone ? Quelle période ?).
\textbf{2. Identifier les acteurs :} Couleurs / symboles pour chaque camp (Axe vs Alliés).
\textbf{3. Repérer les flèches d'offensive :} Direction, épaisseur (importance des forces), dates associées.
\textbf{4. Localiser les lignes de front :} Ligne continue (front stable), pointillée (front mouvant), hachures (zones occupées).
\textbf{5. Noter les objectifs stratégiques :} Villes encerclées, ressources visées (pétrole, ports), capitales.
\textbf{6. Synthétiser :} Décrire la dynamique (éclair, enlisement, recul) et le rapport de forces visible.
\end{methode}

\courssub{Le basculement : constitution de la « Grande Alliance » (1941-1942)}

\noindent Fin 1941, le monde est divisé en deux blocs antagonistes.

\begin{center}
\begin{tabularx}{\linewidth}{|l|X|X|}
\hline
\rowcolor{popblueL}
\textbf{Dimension} & \textbf{L'Axe (Rome-Berlin-Tokyo)} & \textbf{La Grande Alliance (Grand Trois + « Nations Unies »)} \\
\hline
\textbf{Idéologie} & Fascisme, Nazisme, Militarisme japonais : hiérarchie raciale, expansionnisme, anti-démocratie. & Démocraties libérales (UK, USA) + Communisme (URSS) : alliance de circonstance contre un ennemi commun. \\
\hline
\textbf{Objectifs de guerre} & Hégémonie continentale (Lebensraum, Sphère de coprospérité asiatique), nouvel ordre raciste. & \cle{Déclaration des Nations Unies} (1er janv. 1942, 26 États) : « Victoire complète », pas de paix séparée, principes de la Charte atlantique (droit des peuples, libre-échange, désarmement). \\
\hline
\textbf{Atouts} & Position centrale (lignes intérieures), préparation longue, supériorité tactique initiale. & \textbf{Supériorité économique écrasante} (PIB, production acier, pétrole, main-d'œuvre), maîtrise des mers (US Navy + Royal Navy), potentiel scientifique. \\
\hline
\textbf{Faiblesses} & Ressources limitées (pétrole, caoutchouc, métaux rares), lignes de communication surétendues, coordination médiocre (pas d'état-major commun). & Méfiance idéologique profonde (Staline vs Churchill/Roosevelt), coordination difficile (conférences), opinion publique US isolationniste avant Pearl Harbor. \\
\hline
\end{tabularx}
\end{center}
\legende{Tableau 2 : Les deux blocs antagonistes à l'entrée en guerre des USA (déc. 1941).}

\noindent La \cle{Conférence d'Arcadie} (Washington, déc. 1941 - janv. 1942) scelle la stratégie \cle{"Germany First"} (Allemagne d'abord) : priorité au théâtre européen, guerre défensive dans le Pacifique. Staline obtient la promesse d'un \cle{Second Front} en Europe (débarquement) pour 1942 (reporté à 1944). La machine de guerre alliée se met en marche : l'économie US bascule en « \cle{Arsenal of Democracy} » (Roosevelt), l'URSS déplace ses usines vers l'Oural, le Royaume-Uni tient le choc grâce au \cle{Prêt-Bail} (Lend-Lease Act, mars 1941).

\begin{exoresolu}[Analyse de document : L'Appel du 18 juin 1940 (extrait)]
\textbf{Document :} « Les chefs qui, depuis de nombreuses années, sont à la tête des armées françaises, ont formé un gouvernement. Ce gouvernement, alléguant la défaite de nos armées, s'est mis en rapport avec l'ennemi pour cesser le combat. [...] Quoi qu'il arrive, la flamme de la résistance française ne doit pas s'éteindre et ne s'éteindra pas. »

\textbf{Question :} En quoi cet extrait marque-t-il la naissance d'une nouvelle légitimité politique face à Vichy ?

\textbf{Solution détaillée :}
\begin{enumerate}
\item \textbf{Contexte :} 18 juin 1940, Londres. De Gaulle, général de brigade, seul membre du gouvernement Reynaud à refuser l'armistice. La France est en déroute militaire.
\item \textbf{Analyse du texte :}
\begin{itemize}
\item \emph{« Chefs [...] ont formé un gouvernement [...] alléguant la défaite [...] cesser le combat » :} Délégitimation du gouvernement Pétain. Le mot « alléguant » suggère un prétexte ; « cesser le combat » oppose la soumission à la poursuite de la guerre.
\item \emph{« La flamme de la résistance française ne doit pas s'éteindre » :} Revendication de la continuité de l'État français (la « flamme »), distinction entre l'État (qui continue) et le gouvernement (qui a trahi). Appel à la résistance comme devoir moral.
\end{itemize}
\item \textbf{Portée historique :} Acte fondateur de la France Libre. Légitimité initiale : nulle (pas de territoire, pas d'armée, peu de soutiens). Légitimité future : construite par l'action (ralliements africains, Bir Hakeim, reconnaissance alliés 1942-43). Le texte oppose d'emblée \textbf{légalisme constitutionnel} (Pétain vote par le Parlement) et \textbf{légitimité de la résistance} (continuité de la guerre).
\end{enumerate}
\end{exoresolu}

\begin{aretenir}[Synthèse : 1939-1941, le basculement dans la guerre mondiale]
\begin{itemize}
\item \textbf{Genèse :} Échec de Versailles + stratégie hitlérienne (Lebensraum, réarmement, apaisement) $\rightarrow$ Guerre inévitable.
\item \textbf{Déclenchement :} Agression Pologne (1er sept. 1939) + Pacte germano-soviétique $\rightarrow$ Blitzkrieg $\rightarrow$ Drôle de guerre.
\item \textbf{Effondrement France (juin 1940) :} Armistice (pas capitulation) $\rightarrow$ Vichy (collaboration) vs France Libre (résistance, Appel 18 juin). \textbf{Cameroun rallié août 1940 (Brunot).}
\item \textbf{Mondialisation (1940-1941) :} Échec Bataille d'Angleterre $\rightarrow$ Barbarossa (juin 1941, guerre d'anéantissement) $\rightarrow$ Pearl Harbor (déc. 1941, entrée US).
\item \textbf{Nouvel équilibre :} Grande Alliance (UK, USA, URSS) = supériorité économique/démographique décisive. Guerre totale : mobilisation intégrale, effacement front/arrière.
\end{itemize}
\end{aretenir}

\begin{exercice}[Entraînement Bac - Sujet type]
\textbf{Sujet :} « La guerre éclate en Europe en 1939, mais elle ne devient véritablement mondiale qu'en 1941. » À l'aide de vos connaissances et de la carte ci-jointe (imaginaire), expliquez ce processus d'extension du conflit. Vous structurerez votre réponse en trois parties : 1. L'agression hitlérienne et l'effondrement de l'ordre européen (1939-1940) ; 2. L'élargissement aux théâtres méditerranéen, atlantique et oriental (1940-1941) ; 3. Le basculement décisif de décembre 1941.
\end{exercice}

\begin{corrige}[Exercice n°1 - Éléments de réponse]
\textbf{1. L'agression hitlérienne et l'effondrement de l'ordre européen (1939-1940) :} Rupture de Versailles (réarmement, Rhénanie, Anschluss, Munich). Agression Pologne (1er sept. 1939) déclenche guerre franco-britannique. Blitzkrieg : chute Pologne (partage germano-soviétique), puis Danemark/Norvège, puis France (mai-juin 1940). Armistice 22 juin 1940 : Europe continentale dominée par l'Axe (Allemagne + Italie). Vichy collabore. France Libre (De Gaulle, 18 juin) continue combat depuis Londres et territoires ralliés (Cameroun, AEF).
\textbf{2. L'élargissement (1940-1941) :} Échec invasion UK (Bataille d'Angleterre, été 1940). Guerre méditerranéenne (Italie en Grèce, Afrique du Nord, Balkans). \textbf{Barbarossa (22 juin 1941) :} Rupture pacte, invasion URSS, guerre d'anéantissement (Einsatzgruppen, commissaires). Enlisement devant Moscou (hiver). Japon : expansion Asie (Indochine, Chine) $\rightarrow$ tensions US (embargo pétrole).
\textbf{3. Basculement décembre 1941 :} \textbf{7 déc : Pearl Harbor} (Japon frappe US). \textbf{11 déc : Allemagne/Italie déclarent guerre aux US.} Naissance \textbf{Grande Alliance} (Roosevelt, Churchill, Staline). Déclaration Nations Unies (1er janv 1942). Stratégie "Germany First". Supériorité économique alliée scelle issue longue.
\end{corrige}

\coursec{La violence paroxystique : génocide, guerre d'anéantissement, société mobilisée}

Dans la section précédente, nous avons vu comment une guerre européenne déclenchée en septembre 1939 s'est transformée, entre juin 1941 (invasion de l'URSS) et décembre 1941 (Pearl Harbor), en un conflit véritablement mondial. Mais cette guerre n'est pas seulement plus \emph{étendue} que les précédentes : elle est aussi d'une \emph{nature} nouvelle. Pour la première fois dans l'histoire, un État organise industriellement l'extermination de populations entières, et des sociétés tout entières sont mobilisées jusqu'à l'épuisement. C'est cette \cle{violence paroxystique} — c'est-à-dire portée à son degré extrême — que nous étudions maintenant, avant d'en mesurer les conséquences en 1945.

\courssub{La Shoah : de la persécution au génocide}

\textbf{De l'exclusion à la persécution (1933-1939).} Dès son arrivée au pouvoir en 1933, le régime nazi mène une politique antisémite par étapes. L'intuition à retenir est simple : les nazis commencent par \emph{sortir les Juifs de la société allemande} avant de songer à les exterminer physiquement. Les \cle{lois de Nuremberg} (15 septembre 1935) privent les Juifs de la citoyenneté allemande et interdisent les mariages entre Juifs et « Aryens » : le Juif devient juridiquement un étranger dans son propre pays. Dans la nuit du 9 au 10 novembre 1938, la \cle{Nuit de Cristal} (\emph{Kristallnacht}) marque le passage à la violence de masse : des centaines de synagogues sont incendiées, des milliers de commerces juifs saccagés, une centaine de Juifs assassinés et 30 000 déportés vers des camps. L'objectif nazi est alors d'\emph{expulser} les Juifs d'Allemagne.

\textbf{De la persécution à la « Solution finale » (1941-1942).} La guerre change tout. L'invasion de la Pologne (1939) place trois millions de Juifs sous domination nazie ; ils sont enfermés dans des ghettos (Varsovie, Lodz). Avec l'invasion de l'URSS (juin 1941), des unités mobiles de tuerie, les \cle{Einsatzgruppen}, massacrent environ un million de Juifs par fusillades (Babi Yar, près de Kiev : 33 771 Juifs tués les 29-30 septembre 1941). Enfin, le 20 janvier 1942, la \cle{conférence de Wannsee}, réunie près de Berlin sous la direction de Reinhard Heydrich, coordonne les administrations pour organiser la « \cle{Solution finale de la question juive} » : la déportation de tous les Juifs d'Europe vers des centres de mise à mort. C'est le passage d'une violence de guerre à un \emph{génocide planifié, bureaucratique et industriel}.

\begin{definition}[La Shoah et le crime contre l'humanité]
La \cle{Shoah} (mot hébreu signifiant « catastrophe ») désigne l'extermination systématique d'environ six millions de Juifs d'Europe par l'Allemagne nazie et ses collaborateurs entre 1941 et 1945. Elle constitue un \cle{génocide} : la destruction intentionnelle d'un groupe humain en tant que tel.
Le \cle{crime contre l'humanité}, défini par le Statut du tribunal de Nuremberg (1945), englobe le meurtre, l'extermination, la réduction en esclavage, la déportation et toute persécution commise contre des populations civiles pour des motifs politiques, raciaux ou religieux. Il est \emph{imprescriptible} : ses auteurs peuvent être jugés à tout moment.
\end{definition}

\begin{popfigure}
\begin{center}
\begin{tikzpicture}[font=\small, node distance=0.9cm]
\node[draw, rounded corners, fill=popblueL, text width=3.4cm, align=center] (nuremberg) {\textbf{1935} \\ Lois de Nuremberg : exclusion juridique};
\node[draw, rounded corners, fill=popblueL, text width=3.4cm, align=center, right=of nuremberg] (cristal) {\textbf{1938} \\ Nuit de Cristal : violence de masse};
\node[draw, rounded corners, fill=poporangeL, text width=3.4cm, align=center, right=of cristal] (ghettos) {\textbf{1939-1941} \\ Ghettos de Pologne};
\node[draw, rounded corners, fill=poporangeL, text width=3.4cm, align=center, below=of ghettos] (einsatz) {\textbf{1941} \\ Einsatzgruppen : fusillades de masse (Babi Yar)};
\node[draw, rounded corners, fill=poppinkL, text width=3.4cm, align=center, left=of einsatz] (wannsee) {\textbf{20 janv. 1942} \\ Conférence de Wannsee : coordination de la « Solution finale »};
\node[draw, rounded corners, fill=poppink!60, text width=3.4cm, align=center, left=of wannsee] (camps) {\textbf{1942-1944} \\ Centres de mise à mort : Auschwitz-Birkenau, Treblinka...};
\draw[-{Stealth}, thick, popdark] (nuremberg) -- (cristal);
\draw[-{Stealth}, thick, popdark] (cristal) -- (ghettos);
\draw[-{Stealth}, thick, popdark] (ghettos) -- (einsatz);
\draw[-{Stealth}, thick, popdark] (einsatz) -- (wannsee);
\draw[-{Stealth}, thick, popdark] (wannsee) -- (camps);
\end{tikzpicture}
\end{center}
\legende{Schéma conceptuel : la mise en place progressive de la « Solution finale » (1935-1944). On observe une radicalisation par étapes : exclusion juridique, violence, enfermement, massacres de masse, puis extermination industrielle coordonnée à Wannsee.}
\end{popfigure}

\textbf{Le fonctionnement du système concentrationnaire et exterminateur.} Les Juifs d'Europe sont convoyés par trains vers les camps. À l'arrivée, une « sélection » envoie les valides au travail forcé et les autres — enfants, vieillards, mères avec nourrissons — directement aux chambres à gaz. Auschwitz-Birkenau, à la fois camp de concentration et centre de mise à mort, devient le symbole de ce système : environ un million de Juifs y sont assassinés.

\begin{popfigure}
\begin{center}
\begin{tikzpicture}
\begin{axis}[
  width=0.85\linewidth, height=6cm,
  ybar, bar width=8pt,
  xlabel={Période}, ylabel={Déportés arrivant à Auschwitz-Birkenau (milliers)},
  symbolic x coords={1942 (moy.), Printemps 1943, Été 1943, Mars-avr. 1944, Mai-juil. 1944},
  xtick=data, x tick label style={rotate=20, anchor=east, font=\footnotesize},
  ymin=0, ymax=500,
  nodes near coords, nodes near coords style={font=\footnotesize},
  every axis plot/.append style={fill=popblue, draw=popdark}
]
\addplot coordinates {(1942 (moy.),15) (Printemps 1943,10) (Été 1943,12) (Mars-avr. 1944,20) (Mai-juil. 1944,140)};
\end{axis}
\end{tikzpicture}
\end{center}
\legende{Évolution mensuelle moyenne des convois de déportés vers Auschwitz-Birkenau (en milliers de personnes, ordres de grandeur d'après le \emph{Kalendarium} d'Auschwitz de Danuta Czech). Le pic de mai-juillet 1944 correspond à la déportation massive des Juifs de Hongrie : environ 430 000 personnes en moins de deux mois, dont la majorité est gazée à l'arrivée.}
\end{popfigure}

\begin{exemplebox}[L'ampleur des pertes : des chiffres à connaître]
\begin{itemize}
\item \textbf{6 millions} de Juifs assassinés, dont environ \textbf{1,5 million d'enfants} : c'est la Shoah.
\item \textbf{27 millions} de morts soviétiques (militaires et civils) : l'URSS paie le plus lourd tribut du conflit.
\item \textbf{50 à 60 millions} de morts au total dans le monde, dont une majorité de \emph{civils} — fait nouveau par rapport à 1914-1918.
\end{itemize}
Ces chiffres doivent être cités avec précision dans une copie de Probatoire ou de Baccalauréat technique : ils fondent la qualification de « guerre totale » et de « guerre d'anéantissement ».
\end{exemplebox}

\begin{attention}[Ne pas confondre camps de concentration et centres de mise à mort]
Erreur fréquente dans les copies : appeler tous les camps nazis « camps de concentration ». Il faut distinguer :
\begin{itemize}
\item les \cle{camps de concentration} (Dachau, Buchenwald, camps de travail) : on y enferme et on y épuise par le travail, la faim et la violence ; la mort y est massive mais n'est pas la fonction première ;
\item les \cle{centres de mise à mort} (ou camps d'extermination) : \textbf{Treblinka, Sobibor, Belzec, Chelmno, Auschwitz-Birkenau, Majdanek}, tous situés en Pologne occupée. Leur fonction unique est de tuer à l'arrivée, par gazage. Treblinka, par exemple, n'a presque pas de détenus permanents : les convois y sont exterminés en quelques heures.
\end{itemize}
Employer le vocabulaire exact est une exigence de rigueur historique et de respect des victimes.
\end{attention}

\courssub{Les autres crimes de masse du nazisme}

La Shoah est au cœur du système criminel nazi, mais elle ne doit pas masquer les autres victimes d'une idéologie qui hiérarchise les vies humaines.

\begin{itemize}
\item \textbf{Les Tziganes (Roms et Sintis)} : environ 250 000 à 500 000 sont assassinés, victimes du même racisme biologique que les Juifs (le \emph{Porajmos}, « dévoration » en romani).
\item \textbf{Les handicapés} : le programme \cle{Aktion T4} (1939-1941) organise l'« euthanasie » d'environ 200 000 personnes handicapées physiques ou mentales, gazées dans des centres spécialisés en Allemagne même. C'est là que sont expérimentées les techniques de gazage ensuite utilisées dans les centres de mise à mort — preuve de la continuité du système.
\item \textbf{Les opposants politiques} : communistes, socialistes, syndicalistes, résistants de tous les pays occupés sont internés, déportés, exécutés.
\item \textbf{Les prisonniers de guerre soviétiques} : sur 5,7 millions de captifs, environ 3,3 millions meurent, la plupart d'une \emph{famine délibérée} organisée par la Wehrmacht, qui refuse de les nourrir conformément à l'idéologie raciale.
\end{itemize}

\courssub{La guerre d'anéantissement à l'Est}

Le front germano-soviétique (1941-1944) n'est pas une guerre « classique » : c'est une \cle{guerre d'anéantissement}, à la fois idéologique et raciale. Hitler veut détruire le « judéo-bolchevisme » et conquérir un \emph{Lebensraum} (« espace vital ») pour la race allemande, en réduisant les Slaves en esclavage.

Cette volonté se traduit par des ordres criminels écrits \emph{avant} l'invasion :
\begin{itemize}
\item l'\cle{ordre des Commissaires} (juin 1941) : tout commissaire politique de l'Armée rouge capturé doit être fusillé immédiatement, en violation du droit de la guerre ;
\item les directives de l'opération \cle{Barbarossa} dispensent les soldats allemands de toute poursuite pour crimes contre les civils soviétiques ;
\item le « plan de la faim » prévoit de prélever les récoltes ukrainiennes pour l'Allemagne, condamnant des millions de civils à la disette (Leningrad assiégée : près d'un million de morts de faim).
\end{itemize}

\begin{aretenir}[L'essentiel sur la violence de guerre]
La guerre à l'Est est une guerre d'anéantissement : elle vise non seulement à vaincre une armée, mais à détruire un État, une idéologie et des populations jugées « inférieures ». Génocide des Juifs et des Tziganes, famine organisée des prisonniers et des civils, ordres criminels : la violence nazie est \emph{planifiée} et repose sur une idéologie raciale. C'est ce caractère qui fondera, en 1945, la notion juridique nouvelle de crime contre l'humanité.
\end{aretenir}

\courssub{La mobilisation totale des sociétés}

Face à un conflit mondial de longue durée, tous les belligérants passent à l'\cle{économie de guerre} : l'État dirige toute la production vers l'effort militaire.

\begin{itemize}
\item \textbf{Allemagne} : Albert Speer, ministre de l'Armement à partir de 1942, organise la production de guerre en exploitant le travail forcé de millions de déportés et de prisonniers.
\item \textbf{URSS} : le \cle{Gosplan} (commission de planification) déplace plus de 1 500 usines à l'est de l'Oural, hors de portée des bombardiers allemands, et y fait travailler des millions d'ouvriers, dont beaucoup de femmes.
\item \textbf{États-Unis} : le \cle{War Production Board} convertit l'industrie civile (automobile, réfrigérateurs) à la production militaire ; l'Amérique devient « l'arsenal des démocraties ».
\end{itemize}

Les \textbf{femmes} remplacent massivement les hommes partis au front : ouvrières (les « Rosie the Riveter » américaines), infirmières, auxiliaires, paysannes. La \textbf{propagande} devient une arme : Goebbels, ministre nazi de la Propagande, proclame en février 1943 la « guerre totale » au Sportpalast de Berlin ; aux États-Unis, l'\emph{Office of War Information} produit affiches et films pour soutenir le moral.

Enfin, des \textbf{résistances} s'organisent partout dans l'Europe occupée.

\begin{definition}[Résistance, collaboration, collaborationnisme]
La \cle{Résistance} désigne l'ensemble des luttes contre l'occupant et les régimes collaborateurs. On distingue :
\begin{itemize}
\item la \emph{Résistance intérieure} : réseaux de renseignement, sabotages, journaux clandestins, maquis (groupes armés cachés) ;
\item la \emph{Résistance extérieure} : combattants ayant rejoint les Alliés, comme les \cle{Forces françaises libres} (FFL) du général de Gaulle, créées à Londres en juin 1940.
\end{itemize}
À ne pas confondre : la \emph{collaboration} est une coopération d'État à État (le régime de Vichy collabore avec l'Allemagne) ; le \emph{collaborationnisme} est un engagement idéologique actif en faveur de l'ennemi (miliciens, dénonciateurs).
\end{definition}

\begin{methode}[Critiquer une source iconographique]
Face à une photographie d'époque (exercice fréquent au Baccalauréat technique), suivre quatre étapes :
\begin{enumerate}
\item \textbf{Identifier} : nature du document (photo de propagande officielle ou photo clandestine ?), auteur, date, lieu, commanditaire.
\item \textbf{Décrire} : cadrage (qu'est-ce qui est montré, qu'est-ce qui est exclu du champ ?), angle de vue, mise en scène des personnages.
\item \textbf{Interpréter} : intention de l'auteur (exalter, rassurer, terroriser, témoigner ?), contexte de diffusion (presse contrôlée, tracts clandestins ?).
\item \textbf{Confronter} : croiser avec d'autres sources. Une photo de propagande de Goebbels montrant des soldats souriants cache la réalité du front ; une photo clandestine d'un camp, prise au péril du photographe (comme celles du \emph{Sonderkommando} d'Auschwitz), a une valeur de témoignage exceptionnelle mais un cadrage contraint.
\end{enumerate}
Conclusion obligatoire : une image n'est jamais neutre ; elle est une construction qu'il faut démonter.
\end{methode}

\courssub{Les tournants militaires : le recul de l'Axe (1942-1944)}

La violence de l'Axe ne lui assure pas la victoire. Entre fin 1942 et 1944, une série de batailles inverse le cours de la guerre :

\begin{itemize}
\item \textbf{Midway} (juin 1942) : la marine américaine détruit quatre porte-avions japonais ; le Japon perd l'initiative dans le Pacifique.
\item \textbf{El Alamein} (octobre-novembre 1942) : les Britanniques de Montgomery arrêtent Rommel en Égypte ; l'Afrique du Nord échappe progressivement à l'Axe.
\item \textbf{Stalingrad} (17 juillet 1942 - 2 février 1943) : tournant majeur de la guerre. La 6\textsuperscript{e} armée allemande du maréchal Paulus, encerclée dans les ruines de la ville, capitule : environ 250 000 soldats de l'Axe sont tués ou faits prisonniers. L'Armée rouge entame sa longue contre-offensive vers Berlin.
\item \textbf{Les débarquements alliés} : Sicile (juillet 1943, qui provoque la chute de Mussolini), \textbf{Normandie} (6 juin 1944, le « Jour J », plus grande opération amphibie de l'histoire) et \textbf{Provence} (15 août 1944). La France est libérée ; l'Allemagne se retrouve prise en tenaille entre l'Ouest et l'Est.
\end{itemize}

\begin{savaistu}[Des Camerounais dans la libération de la France]
Des tirailleurs camerounais, engagés dans les forces de la France libre dès 1940, ont participé à la libération de l'Europe. Le \cle{Régiment de Marche du Cameroun} a combattu en Italie — notamment lors de la terrible bataille de Monte Cassino (1944) — puis a débarqué en \textbf{Provence en août 1944} avant de remonter la vallée du Rhône. Leur contribution, longtemps oubliée des manuels, rappelle que la libération de la France fut aussi l'œuvre des soldats africains : c'est le fil conducteur du dossier sur le rôle du Cameroun dans la guerre, traité en travaux dirigés.
\end{savaistu}

\begin{exoresolu}[Exercice d'application : analyser un tableau de données]
\textbf{Énoncé.} Voici les pertes humaines approximatives de quelques pays (en millions) : URSS : 27 ; Chine : 15 ; Allemagne : 7 ; Pologne : 6 ; Japon : 3 ; France : 0,6 ; États-Unis : 0,4. Total mondial : 55 millions environ.
1. Calculez la part de l'URSS dans les pertes mondiales.
2. Pourquoi peut-on parler de « guerre d'anéantissement » sur le front de l'Est ?
\tcblower
\textbf{Solution détaillée.}
1. Part de l'URSS : $\dfrac{27}{55} \times 100 \approx 49\,\%$. L'URSS supporte à elle seule près de la moitié des pertes humaines mondiales du conflit.
2. Trois arguments, appuyés sur le cours : (a) le caractère idéologique et racial du conflit (destruction du « judéo-bolchevisme », \emph{Lebensraum}) ; (b) les ordres criminels (ordre des Commissaires, directives Barbarossa) et la famine délibérée des 3,3 millions de prisonniers soviétiques ; (c) la politique d'extermination menée sur les territoires occupés (Einsatzgruppen, centres de mise à mort en Pologne). La disproportion des pertes soviétiques s'explique ainsi par une guerre qui visait les populations autant que les armées.
\end{exoresolu}

\begin{exercice}[Pour s'entraîner]
1. Définis en deux phrases le crime contre l'humanité et explique pourquoi cette notion est créée en 1945.
2. Construis une frise chronologique plaçant : lois de Nuremberg, Nuit de Cristal, Wannsee, Stalingrad, débarquement de Normandie.
3. À partir de la méthode vue en classe, rédige en dix lignes l'analyse critique d'une affiche de propagande de ton choix (nazie ou alliée) : nature, cadrage, intention, diffusion.
\end{exercice}

\begin{corrige}[Exercice]
1. Le crime contre l'humanité désigne, selon le Statut de Nuremberg (1945), le meurtre, l'extermination, la déportation, la réduction en esclavage et les persécutions commises contre des populations civiles pour des motifs politiques, raciaux ou religieux. Il est créé en 1945 parce que les crimes nazis (Shoah, famine des prisonniers soviétiques) ne rentraient dans aucune catégorie juridique existante : il fallait un outil nouveau pour juger des crimes d'État, commis parfois contre les propres citoyens de l'État criminel.
2. 1935 : lois de Nuremberg ; 1938 : Nuit de Cristal ; 1942 : conférence de Wannsee ; 1943 : fin de la bataille de Stalingrad ; 6 juin 1944 : débarquement de Normandie. La frise doit montrer la radicalisation (1935-1942) puis le recul militaire de l'Axe (1943-1944).
3. Attendu : identification (auteur, date, commanditaire), description du cadrage et des symboles, intention (mobiliser, exalter, désigner un ennemi), mode de diffusion, et mise en perspective critique (ce que l'affiche cache). La copie doit conclure sur la non-neutralité de l'image.
\end{corrige}

Ainsi, entre 1941 et 1944, la guerre atteint un degré de violence inédit : génocide industriel, guerre d'anéantissement à l'Est, mobilisation totale des sociétés. Mais cette violence extrême appelle une réponse : dès 1945, les vainqueurs devront construire un ordre nouveau pour qu'elle ne se reproduise pas — c'est l'objet de la section 4, consacrée aux conséquences du conflit et à la naissance de l'ONU, après que la section 3 aura montré comment l'Afrique et le Cameroun ont été entraînés dans cette tourmente mondiale.

\coursec{L'Afrique et le Cameroun dans la tourmente : enjeux, contributions, mutations}

La violence paroxystique qui caractérise le conflit en Europe — génocide, guerre d'anéantissement, mobilisation totale des sociétés — trouve un écho particulier sur le continent africain. Loin d'être un simple théâtre d'opérations périphérique, l'Afrique devient, dès 1940, le \cle{territoire de résistance} de la France libre et le \cle{réservoir stratégique} indispensable à l'effort de guerre allié. Le Cameroun, par sa position géographique (golfe de Guinée), sa production de matières premières et son engagement humain massif, incarne cette bascule : d'une colonie soumise à l'exploitation, il devient un acteur de la victoire, dont les sacrifices accélèrent la prise de conscience politique et posent les jalons de la décolonisation.

\courssub{L'Afrique, enjeu stratégique planétaire : matières, bases et main-d'œuvre}

Dès la défaite de juin 1940 et l'armistice, l'Empire colonial français se fracture entre partisans de Vichy et partisans de la France libre. L'Afrique-Équatoriale française (AEF) et le Cameroun, ralliés à de Gaulle dès août 1940, deviennent les \cle{premiers territoires de l'Empire} à se placer aux côtés des Alliés, bien avant le ralliement de l'Afrique du Nord (novembre 1942). Cette position confère à l'Afrique trois fonctions vitales pour l'effort de guerre allié.

\begin{definition}[Réservoir stratégique et matières premières]
L'Afrique fournit des \cle{matières premières} indispensables à l'industrie de guerre alliée : le \cle{caoutchouc} (pneus, chambres à air pour avions et véhicules), l'\cle{étain} (soudure), l'\cle{uranium} du Congo belge (utilisé dans le projet Manhattan), la \cle{bauxite} (aluminium pour l'aéronautique), l'huile de palme et le cacao. Après la perte de l'Asie du Sud-Est (occupée par le Japon en 1942), premier producteur mondial de caoutchouc, l'Afrique de l'Ouest et le Cameroun deviennent des fournisseurs essentiels de caoutchouc naturel pour les Alliés.
\end{definition}

\begin{definition}[Prestations (travail forcé légalisé)]
Les \cle{prestations} désignent le travail forcé légalisé en AEF et en AOF pour des travaux d'intérêt public (routes, ports, plantations, évacuation des récoltes). Instaurées dans l'Empire français au début du XX\textsuperscript{e} siècle, elles sont \emph{intensifiées} pendant la guerre : durée prolongée, quotas accrus, répression des récalcitrants. Elles constituent la main-d'œuvre invisible qui permet l'extraction et l'évacuation des richesses.
\end{definition}

L'Afrique est aussi une \cle{plateforme logistique}. La \cle{route Sud-Atlantique} (South Atlantic Air Ferry Route) relie le Brésil (Natal) à l'Afrique de l'Ouest puis au Moyen-Orient. Les bases aériennes de \cle{Dakar}, \cle{Pointe-Noire}, \cle{Brazzaville} et \cle{Douala} (Cameroun) servent d'escales aux avions alliés qui ravitaillent le front du Moyen-Orient.

\begin{savaistu}[La base aérienne de Douala, carrefour stratégique]
L'aérodrome de Douala sert, pendant la guerre, d'escale sur la \emph{South Atlantic Air Ferry Route} : les avions alliés y font escale pour ravitailler en carburant avant de poursuivre vers l'Afrique centrale, le Soudan (Khartoum) puis l'Égypte (Le Caire). Cette infrastructure, construite et entretenue par la main-d'œuvre camerounaise (prestations), laisse en héritage une piste capable d'accueillir les gros porteurs, fondement de l'aviation civile camerounaise d'après-guerre.
\end{savaistu}

Enfin, l'Afrique est un \cle{réservoir de main-d'œuvre militaire}. Des dizaines de milliers de tirailleurs d'Afrique subsaharienne sont engagés dans les Forces françaises libres (FFL) puis dans l'armée de la Libération ; le Cameroun fournit un contingent important par rapport à sa population.

\courssub{Le ralliement de l'AEF et du Cameroun (août 1940) : la rupture avec Vichy}

Le 26 août 1940, le gouverneur du Tchad, \cle{Félix Éboué} (premier gouverneur noir d'une colonie française), proclame le ralliement du Tchad au général de Gaulle ; l'ensemble de l'AEF suit dans les jours suivants. Le 27 août, le gouverneur du Cameroun, \cle{Richard Brunot}, rallie à son tour le territoire à la France libre. C'est un acte politique majeur : la France libre dispose désormais d'un territoire, d'une population et de ressources pour exister physiquement.

\begin{exemplebox}[Dossier Programme : Le rôle du Cameroun dans la Deuxième Guerre mondiale]
\textbf{Document 1 : Déclaration de ralliement du Cameroun (27 août 1940).} Le Gouverneur du Cameroun décide que le territoire continue la guerre aux côtés de la Grande-Bretagne et des Alliés.
\textbf{Document 2 : Lettre d'un tirailleur camerounais (1944).} « Nous nous battons pour la France afin qu'à notre retour, nous soyons traités en hommes libres. »
\textbf{Document 3 : Statistiques de production du caoutchouc.} La production camerounaise est multipliée par plusieurs fois entre 1939 et 1943, grâce à l'intensification des prestations obligatoires dans les plantations du Sud.
\textbf{Document 4 : Affiche de propagande « L'Afrique aide la France » (1942).} Elle représente un tirailleur souriant offrant du caoutchouc à une Marianne libératrice. Elle masque la contrainte réelle.
\textbf{Exploitation pédagogique :} Croiser les sources officielles (déclaration, statistiques) et le vécu (lettre) pour montrer l'écart entre le discours volontariste et la réalité contrainte.
\end{exemplebox}

Ce ralliement n'est pas sans heurts. Le Gabon, resté fidèle à Vichy, est conquis militairement en novembre 1940 (bataille de Libreville) par les forces de la France libre, parmi lesquelles des unités formées au Cameroun. Le blocus et les combats perturbent l'économie camerounaise (pénurie de tissus, médicaments, carburant) jusqu'en 1942.

\courssub{L'effort de guerre camerounais : hommes, matières, contrainte}

L'engagement camerounais repose sur une triple mobilisation : humaine, économique, contrainte.

\begin{exemplebox}[Exemple chiffré : L'apport camerounais]
\begin{itemize}
    \item \textbf{Hommes :} Plusieurs milliers de Camerounais sous les drapeaux (régiments de tirailleurs, bataillons de marche, services auxiliaires). Des unités camerounaises combattent au Gabon (1940), puis aux côtés des Alliés en Afrique du Nord, en Italie et en France (1944-1945).
    \item \textbf{Caoutchouc :} La production est fortement multipliée entre 1939 et 1943 pour compenser la perte des plantations d'Asie du Sud-Est occupées par le Japon.
    \item \textbf{Produits vivriers :} Réquisitions massives de manioc, maïs, bananes, huile de palme pour nourrir les troupes et les populations urbaines (Douala, Yaoundé).
    \item \textbf{Bois :} Exploitation intensive des forêts du Sud pour le charbon de bois (gazogènes) et les traverses de chemin de fer.
\end{itemize}
\end{exemplebox}

\begin{attention}[Le piège du « volontariat » : contrainte administrative et chefferie]
Ne pas surestimer l'enthousiasme spontané. Le recrutement repose sur une \cle{pression administrative} intense : les chefs de canton et de village reçoivent des \cle{quotas d'hommes} à fournir. Le refus est sanctionné (amendes, destitution, prison). La prime d'engagement et la solde attirent des volontaires, mais beaucoup répondent à l'ordre colonial. L'historien doit distinguer \emph{engagement volontaire} (choix individuel) et \emph{mobilisation contrainte} (pression collective via l'administration coloniale).
\end{attention}

L'intensification des \cle{prestations} (travail forcé) est la clé de voûte de cet effort. Pour évacuer le caoutchouc vers Douala, il faut des routes, des ports, des voies ferrées. Les populations rurales (surtout dans le Sud forestier et le Nord) sont réquisitionnées pour l'entretien de la voie ferrée Douala-Yaoundé, du port de Douala et des pistes d'aviation. Cette surexploitation provoque des fuites de population vers la brousse, une mortalité accrue et des disettes locales masquées par les exportations records.

\begin{popfigure}
\begin{tikzpicture}[scale=0.75, every node/.style={font=\small}]
    % Contour très simplifié du Cameroun (polygone approximatif)
    \draw[popdark, thick, fill=popcream]
        (0,0) -- (8,0) -- (9,2) -- (8.5,5) -- (7,7) -- (4,7.5) -- (1,6) -- (0,3) -- cycle;

    % Villes principales
    \node[circle, fill=popblue, inner sep=2pt, label=below:\textbf{Douala}] (Douala) at (1.5, 1.2) {};
    \node[circle, fill=popblue, inner sep=2pt, label=above:\textbf{Yaoundé}] (Yaounde) at (4.5, 3.5) {};
    \node[circle, fill=popblue, inner sep=2pt, label=above:\textbf{Garoua}] (Garoua) at (7.5, 5.5) {};
    \node[circle, fill=popblue, inner sep=2pt, label=left:\textbf{Nkongsamba}] (Nkongsamba) at (2.5, 2.5) {};
    \node[circle, fill=popblue, inner sep=2pt, label=right:\textbf{Bertoua}] (Bertoua) at (6.5, 4.5) {};

    % Base aérienne / Port (Douala)
    \node[star, star points=5, star point ratio=2.25, fill=popgold, draw=poporange, thick, inner sep=3pt] at (1.5, 1.2) {};

    % Zones de recrutement (surfacique - hachuré)
    \pattern[pattern=north east lines, pattern color=poppink!50] (2,2) rectangle (5,5);
    \pattern[pattern=north west lines, pattern color=poppink!50] (6,4) rectangle (8.5,6.5);
    \node[poppink, font=\tiny\bfseries] at (3.5, 3.5) {Zone recrutement intense};
    \node[poppink, font=\tiny\bfseries, rotate=-20] at (7.2, 5.2) {Zone Nord};

    % Axes d'évacuation matières premières
    \draw[popgreen, very thick, ->, >=Stealth] (Bertoua) -- (Nkongsamba) -- (Douala) node[midway, above, sloped, font=\tiny] {Bois, Caoutchouc};
    \draw[popgreen, very thick, ->, >=Stealth] (Yaounde) -- (Douala) node[midway, above, sloped, font=\tiny] {Cacao, Vivres};
    \draw[popgreen, very thick, ->, >=Stealth] (Garoua) -- (Yaounde) node[midway, above, sloped, font=\tiny] {Coton, Bétail};

    % Route aérienne Sud-Atlantique
    \draw[poppurple, thick, dashed, ->, >=Stealth] (-1, 1) -- (Douala) node[midway, above, font=\tiny] {Route aérienne Sud-Atlantique};

    % Flèche ralliement (depuis l'AEF/Est)
    \draw[popblue, thick, ->, >=Stealth] (8.5, 3.5) -- (7, 3.5) node[midway, above, font=\tiny] {Ralliement AEF (août 1940)};

    % Légende intégrée
    \begin{scope}[xshift=10cm, yshift=2cm]
        \draw[popdark, thick, fill=popcream] (0,0) rectangle (5,4.5);
        \node[align=left, font=\small\bfseries] at (2.5, 4.2) {Légende};
        \node[circle, fill=popblue, inner sep=2pt] at (0.5, 3.5) {}; \node[anchor=west, font=\small] at (0.8, 3.5) {Villes principales};
        \node[star, star points=5, star point ratio=2.25, fill=popgold, draw=poporange, thick, inner sep=3pt] at (0.5, 3.0) {}; \node[anchor=west, font=\small] at (0.8, 3.0) {Base aérienne / Port stratégique};
        \draw[popgreen, very thick, ->, >=Stealth] (0.3, 2.5) -- (1.5, 2.5); \node[anchor=west, font=\small] at (1.6, 2.5) {Axes d'évacuation (rail/route)};
        \draw[poppurple, thick, dashed, ->, >=Stealth] (0.3, 2.0) -- (1.5, 2.0); \node[anchor=west, font=\small] at (1.6, 2.0) {Route aérienne Sud-Atlantique};
        \draw[popblue, thick, ->, >=Stealth] (0.3, 1.5) -- (1.5, 1.5); \node[anchor=west, font=\small] at (1.6, 1.5) {Flux de ralliement / Renforts};
        \fill[pattern=north east lines, pattern color=poppink!50] (0.2, 1.0) rectangle (1.0, 0.5); \node[anchor=west, font=\small] at (1.1, 0.75) {Zones de recrutement / prestations};
    \end{scope}

    % Flèche Nord
    \node[draw=popdark, thick, circle, inner sep=2pt, label=90:\textbf{N}] at (9.5, 6.5) {};
\end{tikzpicture}
\legende{Croquis commenté : le Cameroun dans la Seconde Guerre mondiale. Le territoire (contour simplifié) montre la centralité de Douala (base aérienne et port), les axes d'évacuation des matières premières (vert) vers la côte, les zones de recrutement et de prestations intenses (hachuré) et la route aérienne stratégique (violet). Croquis schématique, sans précision cartographique.}
\end{popfigure}

\courssub{L'économie de guerre : le pic des exportations (1942-1943)}

L'effort de guerre se lit dans les statistiques douanières. Le blocus et la rupture avec la métropole (1940-1942) effondrent d'abord les exportations. Le ralliement à la France libre et l'entrée en guerre des États-Unis (décembre 1941) ouvrent le marché allié (États-Unis, Grande-Bretagne, Afrique du Nord). Les prix sont contrôlés, mais les volumes explosent, en particulier pour le caoutchouc.

\begin{popfigure}
\begin{tikzpicture}
\begin{axis}[
    width=\linewidth,
    height=10cm,
    title={Évolution indicative des exportations camerounaises (1938-1945)},
    xlabel={Année},
    ylabel={Volume (en milliers de tonnes)},
    xmin=1937.5, xmax=1945.5,
    ymin=0, ymax=160,
    xtick={1938,1939,1940,1941,1942,1943,1944,1945},
    xticklabel style={rotate=45, anchor=east, font=\small},
    ytick={0,20,40,60,80,100,120,140,160},
    legend pos=north west,
    legend style={font=\small, fill=white, draw=popline},
    grid=major, grid style={popline},
    popcurve/.style={thick, mark=*, mark size=2.5},
]
\addplot[popblue, popcurve] coordinates {
    (1938, 1.5) (1939, 2) (1940, 1) (1941, 3) (1942, 8) (1943, 14) (1944, 10) (1945, 6)
};
\addlegendentry{Caoutchouc (ordre de grandeur)}

\addplot[popgreen, popcurve] coordinates {
    (1938, 85) (1939, 80) (1940, 20) (1941, 30) (1942, 55) (1943, 70) (1944, 75) (1945, 80)
};
\addlegendentry{Cacao (ordre de grandeur)}

\addplot[poporange, popcurve] coordinates {
    (1938, 45) (1939, 50) (1940, 40) (1941, 55) (1942, 70) (1943, 85) (1944, 80) (1945, 60)
};
\addlegendentry{Bois (ordre de grandeur)}

\node[popblue, anchor=south, font=\small\bfseries] at (axis cs:1943,14) {Pic 1943};
\draw[popblue, ->, thick] (axis cs:1943.5,13) -- (axis cs:1943.1,14.1);
\end{axis}
\end{tikzpicture}
\legende{Graphique : l'explosion du caoutchouc (bleu) masque la chute temporaire du cacao (vert) due au blocus (1940-1941). Le bois (orange) progresse pour les gazogènes. Valeurs indicatives reconstituées à partir des tendances des statistiques douanières camerounaises, 1938-1945.}
\end{popfigure}

\courssub{Conséquences sociales : urbanisation, inflation, naissance du syndicalisme}

La guerre transforme en profondeur la société camerounaise.

\begin{propriete}[Mutations sociales accélérées (1940-1945)]
\begin{itemize}
    \item \textbf{Urbanisation accélérée :} Douala et Yaoundé voient leur population croître fortement (afflux de recrues, fonctionnaires, travailleurs, populations fuyant les réquisitions rurales). Apparaissent des quartiers spontanés, une crise du logement et de l'insalubrité.
    \item \textbf{Inflation et pénuries :} Blocus (1940-1942) puis inflation liée à l'économie de guerre ; le pouvoir d'achat s'effondre. Développement du marché noir et de la corruption.
    \item \textbf{Diffusion d'idées nouvelles :} Contact avec les soldats alliés, diffusion de la \cle{Charte de l'Atlantique} (1941 : droit des peuples à disposer d'eux-mêmes), propagande antifasciste : prise de conscience des contradictions coloniales (lutte pour la liberté en Europe, déni de liberté en Afrique).
    \item \textbf{Montée des revendications :} Grèves de 1945 (Douala, Edéa, Yaoundé) : revendications salariales, fin du travail forcé, droit syndical. Création de syndicats camerounais (1944) puis, dans ce terreau, des premiers partis nationalistes après 1945.
\end{itemize}
\end{propriete}

\begin{exoresolu}[Analyse d'un document : la grève des dockers de Douala (1945)]
\textbf{Énoncé :} À partir de l'extrait ci-dessous, expliquez comment la guerre a fait émerger une conscience de classe et politique chez les travailleurs camerounais.
\emph{« Les dockers de Douala cessent le travail. Ils réclament l'alignement des salaires sur le coût de la vie, la suppression du travail forcé (prestations), la libération des prisonniers. L'administration répond par l'arrestation des meneurs et l'envoi de troupes. La grève dure plusieurs jours. »} (Rapport administratif, 1945-1946, adapté).

\tcblower

\textbf{Solution détaillée :}
\begin{enumerate}
    \item \textbf{Contexte de guerre :} L'inflation est une conséquence directe de l'économie de guerre (pénuries, émission monétaire). Les dockers, essentiels à l'évacuation du caoutchouc, ont conscience de leur pouvoir de nuisance stratégique.
    \item \textbf{Revendications hybrides :} Socio-économiques (salaires) ET politiques (fin des prestations, prisonniers). Le mot « prestations » renvoie au travail forcé, symbole de la soumission coloniale.
    \item \textbf{Référentiels nouveaux :} La référence au « coût de la vie » montre l'appropriation du vocabulaire syndical métropolitain et des idées nouvelles (Charte de l'Atlantique, Conférence de Brazzaville, 1944).
    \item \textbf{Répression et organisation :} L'envoi de troupes montre la fébrilité coloniale ; l'ampleur du mouvement prouve l'organisation ouvrière naissante.
    \item \textbf{Conclusion :} La guerre a créé les conditions objectives (concentration ouvrière, inflation) et subjectives (idées nouvelles, expérience du monde) de l'entrée en politique des travailleurs camerounais.
\end{enumerate}
\end{exoresolu}

\courssub{La Conférence de Brazzaville (janvier-février 1944) : l'illusion réformiste}

Convoquée par le général de Gaulle et préparée par Félix Éboué (qui meurt en mai 1944), cette conférence rassemble les gouverneurs de l'Empire français. Elle doit définir l'avenir de la colonisation après la victoire.

\begin{aretenir}[Les résolutions de Brazzaville (1944) : rupture ou continuité ?]
\begin{itemize}
    \item \textbf{Rejet catégorique de l'autonomie :} « Les fins de la civilisation [...] excluent toute idée d'autonomie, toute possibilité d'évolution hors du bloc français de l'Empire » (discours de Gaulle). L'indépendance n'est pas à l'ordre du jour.
    \item \textbf{Promesses de réformes majeures :} suppression de l'indigénat (abrogé en 1946), accès à la citoyenneté française (loi Lamine Guèye, 1946), fin officielle du \cle{travail forcé} (loi de 1946), représentation au Parlement français (députés), création de conseils locaux élus.
    \item \textbf{Portée historique :} Premier discours officiel reconnaissant les « devoirs » de la France envers ses sujets. Il légitime les revendications nationalistes qui exigeront l'application stricte de ces promesses. Éboué, artisan du ralliement, meurt avant la fin de la guerre : loyaliste convaincu, il a ouvert une voie que la métropole refermera mal.
\end{itemize}
\end{aretenir}

\begin{methode}[Construire un croquis commenté : « Le Cameroun dans la Seconde Guerre mondiale »]
\textbf{Consigne de TP / Examen (Probatoire) :} à partir d'une carte muette du Cameroun, réaliser un croquis organisé et légendé.
\begin{enumerate}
    \item \textbf{Analyse du sujet :} identifier les 3 dimensions : \emph{Stratégie} (bases, axes), \emph{Économie} (zones de production, flux), \emph{Humain} (recrutement, urbanisation).
    \item \textbf{Choix des symboles (sémiologie graphique) :}
        \begin{itemize}
            \item \textbf{Ponctuels (points) :} villes (Douala, Yaoundé), bases aériennes (étoile), ports (ancre), centres de recrutement.
            \item \textbf{Linéaires (flèches/traits) :} axes ferroviaires et routiers d'évacuation (flèches vertes d'épaisseur proportionnelle au volume), route aérienne (flèche pointillée violette), flux de ralliement (flèche bleue).
            \item \textbf{Surfaciques (zones colorées/hachurées) :} zones de production caoutchouc/bois (hachures vertes), zones de recrutement intensif et de prestations (hachures rouges), zones urbaines en croissance.
        \end{itemize}
    \item \textbf{Construction de la légende :} organisée en 3 items numérotés (I. Stratégie, II. Économie de guerre, III. Société mobilisée). Chaque item regroupe ses symboles. Titre précis, sources, échelle indicative, flèche du Nord.
    \item \textbf{Soin graphique :} propreté, lisibilité (pas de surcharge), couleurs contrastées (4 ou 5 maximum), écriture lisible à la règle.
\end{enumerate}
\textbf{Critères d'évaluation (barème type) :} titre/légende/échelle/Nord (4 pts) + choix des symboles pertinents (4 pts) + exactitude des localisations (4 pts) + organisation de la légende (4 pts) + soin et esthétique (4 pts) = 20 pts.
\end{methode}

\begin{exercice}[Entraînement : croquis et analyse]
\begin{enumerate}
    \item Sur une carte muette du Cameroun fournie, réaliser le croquis « Le Cameroun dans la Seconde Guerre mondiale » en appliquant la méthode ci-dessus.
    \item \textbf{Question d'analyse :} « En quoi le croquis réalisé met-il en évidence la contradiction entre l'apport massif du Cameroun à la victoire alliée et le maintien de l'ordre colonial après 1945 ? » (Rédaction attendue : 20 lignes, structurée en 2 arguments illustrés par le croquis).
\end{enumerate}
\end{exercice}

\begin{corrige}[Exercice n°1 - Éléments de corrigé]
\textbf{1. Croquis :} voir le modèle en figure ci-dessus. Vérifier : Douala (base/port) central, flèches vertes convergentes vers Douala, hachures dans le Sud et le Nord, flèche aérienne Ouest-Est, légende structurée en 3 parties.

\textbf{2. Argumentation type :}
\begin{itemize}
    \item \textbf{Argument 1 (visible sur le croquis — flux et production) :} la convergence de toutes les flèches vertes (caoutchouc, bois, cacao, coton) vers Douala montre l'\emph{extraction unilatérale} des richesses pour l'effort de guerre allié, sans transformation locale, au profit de la métropole et des Alliés.
    \item \textbf{Argument 2 (visible sur le croquis — humain et contrainte) :} la superposition des zones de recrutement (hachures) et des zones de production montre que les \emph{mêmes populations} fournissent les hommes (tirailleurs) ET la main-d'œuvre forcée (prestations) pour évacuer la production. Le sang et la sueur camerounais ont payé la libération de la France.
    \item \textbf{Synthèse :} Brazzaville (1944) promet la fin du travail forcé et la citoyenneté, mais le croquis révèle une économie et une administration restées coloniales (extraction, contrainte). Cette contradiction alimente le nationalisme camerounais (grèves de 1945, premiers partis).
\end{itemize}
\end{corrige}

\coursec{Un monde nouveau - Conséquences et naissance de l'ONU}

Nous avons vu, dans la section précédente, que l'Afrique et le Cameroun ont payé un lourd tribut à la guerre : soldats engagés dans les forces françaises libres, porteurs, réquisitions, et ralliement du Cameroun dès 1940 au général de Gaulle. Mais ce conflit, le plus meurtrier de l'histoire de l'humanité, ne s'achève pas seulement par la capitulation de l'Allemagne (8 mai 1945) et du Japon (2 septembre 1945). Il laisse derrière lui un monde détruit, déplacé, traumatisé — et il fait naître, dans les cendres, de nouvelles institutions et de nouvelles règles destinées à empêcher le retour d'une telle catastrophe. C'est ce \cle{monde nouveau} de 1945 que nous allons maintenant étudier : le bilan du conflit, la justice rendue aux vaincus, le redécoupage de l'Europe, la naissance de l'\cle{ONU}, et enfin les premiers signes d'un nouvel affrontement, la \cle{Guerre froide}.

\courssub{Un bilan humain et matériel sans précédent}

\textbf{Intuition.} Pour mesurer l'ampleur d'une catastrophe, on compte les morts, les ruines et les déplacés. Appliquons cette démarche à 1945 : les chiffres dépassent tout ce que l'histoire avait connu jusqu'alors.

La Deuxième Guerre mondiale a fait environ \cle{60 millions de morts}, soit près de 3\,\% de la population mondiale de l'époque (environ 2,3 milliards d'habitants). La Première Guerre mondiale, avec ses 10 millions de morts, fait figure de conflit « limité » en comparaison. Surtout, la nature des pertes a changé : les deux tiers des victimes sont des \cle{civils} (bombardements, massacres, famines, génocide), alors que la guerre de 1914-1918 tuait d'abord des soldats. L'URSS paie le plus lourd tribut (environ 27 millions de morts), suivie de la Chine, de l'Allemagne et de la Pologne (qui perd près de 20\,\% de sa population, dont 3 millions de Juifs polonais assassinés dans la Shoah).

\begin{exemplebox}[Le coût matériel du conflit]
Le coût matériel total de la guerre est estimé à plus de \textbf{1000 milliards de dollars de 1945}. Des villes entières sont rasées : \cle{Varsovie} détruite à plus de 80\,\% après l'insurrection de 1944, \cle{Berlin} réduite en champs de ruines par les bombardements alliés et la bataille finale, \cle{Tokyo} incendiée par les bombardements américains de mars 1945 (environ 100\,000 morts en une nuit), \cle{Hiroshima} et \cle{Nagasaki} anéanties par les deux bombes atomiques des 6 et 9 août 1945 (environ 200\,000 morts au total, immédiats et différés). En Europe, des millions de logements, de ponts, de voies ferrées et d'usines sont détruits : il faudra une décennie pour reconstruire.
\end{exemplebox}

À cette destruction s'ajoutent des déplacements de populations massifs. Environ \cle{12 millions d'Allemands} sont expulsés entre 1945 et 1950 des territoires où ils vivaient depuis des siècles (Prusse orientale, Silésie, Sudètes, Poméranie), au fur et à mesure que les frontières de la Pologne et de la Tchécoslovaquie sont déplacées vers l'ouest. S'y ajoutent des millions de « personnes déplacées » (rescapés des camps, prisonniers, travailleurs forcés) que les Alliés doivent rapatrier. L'Europe de 1945 est un continent de réfugiés.

\courssub{Juger les vaincus : les procès de Nuremberg et de Tokyo}

\textbf{Intuition.} Après la découverte des camps d'extermination, les Alliés refusent de fusiller sommairement les responsables nazis : ils choisissent de les \emph{juger}, pour établir les faits, fonder la mémoire et créer un droit nouveau. C'est une rupture historique : pour la première fois, des chefs d'État et des hauts dignitaires répondent personnellement de leurs actes devant la justice internationale.

Le \cle{procès de Nuremberg} se tient de \textbf{novembre 1945 à octobre 1946} devant un Tribunal militaire international composé de juges des quatre puissances victorieuses (États-Unis, URSS, Royaume-Uni, France). Vingt-deux dignitaires nazis y comparaissent (Göring, Hess, Ribbentrop...). Douze sont condamnés à mort. Le \cle{procès de Tokyo} (1946-1948) juge de la même manière les criminels de guerre japonais.

\begin{definition}[Crime contre la paix]
Le \cle{crime contre la paix} désigne la \textbf{planification, la préparation, le déclenchement ou la conduite d'une guerre d'agression}, en violation des traités et accords internationaux. C'est l'innovation juridique majeure de Nuremberg : faire la guerre n'est plus un droit des États, mais un crime engageant la responsabilité \emph{individuelle} de ceux qui la décident. Le tribunal définit aussi le \cle{crime de guerre} (violation des lois de la guerre : massacres de prisonniers, déportations) et le \cle{crime contre l'humanité} (extermination, réduction en esclavage et persécution de populations civiles, pour des motifs politiques, raciaux ou religieux).
\end{definition}

Ces procès ont une portée immense : ils affirment qu'obéir aux ordres n'exonère pas de la responsabilité personnelle, et ils fondent le droit international pénal moderne (qui aboutira, bien plus tard, à la Cour pénale internationale). En parallèle, les Alliés mènent en Allemagne une politique de \cle{dénazification} : épuration de l'administration, interdiction du parti nazi, rééducation démocratique. Mais cette dénazification reste \textbf{inachevée} : devant l'ampleur de la tâche et les besoins de la reconstruction, de nombreux anciens nazis de second plan conservent leurs postes, surtout en Allemagne de l'Ouest.

\begin{attention}[Piège fréquent]
Ne pas écrire que le procès de Nuremberg a jugé « l'Allemagne » ou « le nazisme » en général : il a jugé des \emph{individus} (et déclaré criminelles certaines organisations comme la SS). De même, ne pas confondre crime de guerre et crime contre l'humanité : le second vise des exactions contre des \emph{populations civiles}, y compris les propres ressortissants de l'État coupable.
\end{attention}

\courssub{Le redessin des frontières : Yalta et Potsdam}

\textbf{Intuition.} Les vainqueurs se réunissent avant même la fin de la guerre pour organiser l'après-guerre. Deux conférences dominent : \cle{Yalta} (Crimée, février 1945 : Roosevelt, Churchill, Staline) et \cle{Potsdam} (juillet-août 1945 : Truman, Attlee, Staline). Les décisions prises redessinent la carte de l'Europe pour près d'un demi-siècle.

Les grandes décisions sont les suivantes :
\begin{itemize}
\item La \textbf{Pologne est déplacée vers l'ouest} : l'URSS annexe ses territoires orientaux (au-delà de la « ligne Curzon »), et la Pologne est compensée par des territoires allemands à l'ouest (Silésie, Poméranie, sud de la Prusse orientale), jusqu'à la ligne Oder-Neisse. D'où l'expulsion des 12 millions d'Allemands déjà évoquée.
\item L'\textbf{URSS annexe} les pays baltes (Estonie, Lettonie, Lituanie, occupés dès 1940), le nord de la Prusse orientale, la Bessarabie et des territoires finlandais.
\item L'\textbf{Allemagne et Berlin} sont divisées en \textbf{quatre zones d'occupation} (américaine, soviétique, britannique, française). Berlin, enclavée dans la zone soviétique, est elle-même partagée en quatre secteurs. L'Autriche et Vienne connaissent le même régime quadripartite.
\item Les États d'Europe de l'Est « libérés » par l'Armée rouge (Pologne, Hongrie, Roumanie, Bulgarie, Tchécoslovaquie) passent sous influence soviétique : Staline y impose progressivement des régimes communistes, les \cle{démocraties populaires}.
\end{itemize}

\begin{popfigure}
\begin{center}
\begin{tikzpicture}[scale=1.05, font=\small]
% Fond Europe
\fill[popcream] (-6.2,-3.2) rectangle (6.2,3.4);
% URSS
\fill[poporange!45] (2.6,-3.2) -- (2.6,3.4) -- (6.2,3.4) -- (6.2,-3.2) -- cycle;
% Bloc Est
\fill[poporange!20] (0.2,-2.6) -- (0.2,1.2) -- (1.2,2.2) -- (2.6,2.4) -- (2.6,-3.2) -- (1.4,-3.2) -- cycle;
% Zone soviétique Allemagne
\fill[poporange!60] (0.2,0.6) -- (0.2,2.0) -- (-0.8,2.6) -- (-0.8,0.6) -- cycle;
% Zones ouest Allemagne
\fill[popblue!35] (-0.8,0.6) -- (-0.8,2.6) -- (-2.6,2.2) -- (-2.6,0.4) -- cycle;
\fill[popblue!20] (-2.6,0.4) -- (-2.6,-1.0) -- (-0.8,-1.0) -- (-0.8,0.6) -- cycle;
\fill[poppurple!25] (-0.8,-1.0) -- (-0.8,0.6) -- (0.2,0.6) -- (0.2,-1.0) -- cycle;
% Europe de l'Ouest
\fill[popblue!10] (-6.2,-3.2) rectangle (-2.6,3.4);
% Berlin
\fill[popdark] (-0.3,1.6) circle (0.07) node[right=1pt, font=\scriptsize]{Berlin (4 secteurs)};
% Rideau de fer
\draw[popdark, line width=1.6pt, dashed] (0.2,-3.2) -- (0.2,1.2) -- (1.2,2.2) -- (2.6,2.4);
\node[popdark, font=\scriptsize\bfseries, rotate=90] at (-0.15,-2.0) {Rideau de fer (à partir de 1946)};
% Flèche déplacement Pologne
\draw[-{Stealth}, popgreen, line width=1.4pt] (1.9,0.3) -- (0.9,0.3);
\node[popgreen, font=\scriptsize, align=center] at (1.4,-0.35) {Pologne déplacée\\vers l'ouest};
% Labels
\node[font=\scriptsize\bfseries] at (4.4,2.6) {URSS (annexions : pays baltes,};
\node[font=\scriptsize\bfseries] at (4.4,2.2) {Prusse or. nord, Bessarabie...)};
\node[font=\scriptsize\bfseries, poporange!80!black] at (1.4,-2.4) {Démocraties populaires};
\node[font=\scriptsize\bfseries, poporange!80!black] at (1.4,-2.8) {(Pologne, Hongrie, Roumanie...)};
\node[font=\scriptsize\bfseries, popblue!80!black] at (-4.4,2.8) {Europe de l'Ouest};
\node[font=\scriptsize, align=center] at (-1.7,1.7) {Zones\\am./brit.};
\node[font=\scriptsize] at (-1.7,-0.3) {Zone fr.};
\node[font=\scriptsize] at (-0.3,1.1) {Zone sov.};
% Nord
\draw[-{Stealth}, popdark] (-5.8,-2.2) -- (-5.8,-1.4) node[above, font=\scriptsize]{N};
\end{tikzpicture}
\end{center}
\legende{Croquis schématique de l'Europe en 1945 : zones d'occupation en Allemagne (américaine et britannique en bleu foncé, française en bleu clair, soviétique en orange foncé), déplacement de la Pologne vers l'ouest (flèche verte), annexions soviétiques à l'est, et tracé du « rideau de fer » séparant les démocraties populaires de l'Europe occidentale. Carte indicative, sans précision géographique.}
\end{popfigure}

\courssub{La naissance de l'ONU : San Francisco, juin 1945}

\textbf{Intuition.} La Société des Nations (SDN), créée en 1919, a échoué à empêcher la guerre : pas d'armée, règle de l'unanimité paralysante, absence des États-Unis. En 1945, les vainqueurs veulent une organisation \emph{plus forte et plus universelle}. Mais l'idée ne naît pas d'un coup en 1945.

\begin{attention}[L'ONU ne naît pas seulement de la guerre]
L'ONU est le fruit d'une maturation longue : la \textbf{Déclaration des 26 nations} du 1\textsuperscript{er} janvier 1942 (les Alliés s'engagent contre l'Axe au nom des principes de la Charte de l'Atlantique), puis la \textbf{conférence de Dumbarton Oaks} (1944) qui prépare le projet de Charte, avant la conférence finale de San Francisco. Ne pas attribuer la création de l'ONU à la seule conférence de Yalta.
\end{attention}

La \cle{conférence de San Francisco} (avril-juin 1945) réunit 50 délégations et aboutit à la signature de la \cle{Charte des Nations unies} le \textbf{26 juin 1945} (entrée en vigueur le 24 octobre 1945). L'ONU compte alors \cle{51 États fondateurs} ; elle en compte 193 aujourd'hui, ce qui témoigne de son universalité croissante, notamment grâce à la décolonisation. Ses objectifs, fixés par la Charte, sont : le maintien de la \textbf{paix} et de la sécurité internationales, le respect des \textbf{droits de l'Homme} (la Déclaration universelle des droits de l'Homme sera adoptée en 1948), et le \textbf{développement} économique et social.

\begin{definition}[Droit de veto]
Le \cle{droit de veto} (article 27 de la Charte) permet à chacun des \textbf{cinq membres permanents} du Conseil de sécurité — États-Unis, URSS (aujourd'hui Russie), Royaume-Uni, France, Chine (RPC depuis 1971) — de \textbf{bloquer toute résolution} du Conseil par un simple vote négatif. Il garantit que les grandes puissances ne quittent pas l'organisation, mais il peut aussi la paralyser lorsqu'elles s'opposent, comme pendant la Guerre froide.
\end{definition}

\begin{popfigure}
\begin{center}
\begin{tikzpicture}[font=\small, node distance=0.5cm,
box/.style={draw=popblue, fill=popblueL, rounded corners=2pt, align=center, minimum width=3.4cm, minimum height=0.85cm},
secr/.style={draw=poppurple, fill=poppurpleL, rounded corners=2pt, align=center, minimum width=3.4cm, minimum height=0.85cm},
cs/.style={draw=poporange, fill=poporangeL, rounded corners=2pt, align=center, minimum width=4.2cm, minimum height=0.85cm}]
\node[box, fill=popblue!30, minimum width=5cm, align=center] (ag){\textbf{Assemblée générale}\\tous les États membres (1 voix chacun)};
\node[cs, right=1.1cm of ag, align=center] (cse){\textbf{Conseil de sécurité}\\15 membres dont \textbf{5 permanents}};
\node[secr, below=0.9cm of ag, align=center] (sec){\textbf{Secrétariat}\\Secrétaire général};
\node[box, below=0.9cm of cse, align=center] (eco){\textbf{ECOSOC}\\Conseil économique et social};
\node[box, below=0.6cm of sec, align=center] (cij){\textbf{CIJ}\\Cour internationale de Justice (La Haye)};
\node[box, below=0.6cm of eco, align=center] (tut){\textbf{Conseil de tutelle}\\(territoires sous tutelle, inactif depuis 1994)};
\draw[-{Stealth}, popline] (ag) -- (sec);
\draw[-{Stealth}, popline] (cse) -- (eco);
\draw[-{Stealth}, popline] (ag.south west) ++(0.3,-0.05) |- (cij.west);
\draw[-{Stealth}, popline] (cse.south east) ++(-0.3,-0.05) |- (tut.east);
% Veto
\node[draw=poporange, fill=poporange!25, rounded corners=2pt, align=center, font=\scriptsize, right=0.7cm of cse, minimum width=3.6cm] (veto) {\textbf{Droit de veto (P5)}\\USA -- Russie -- Chine\\Royaume-Uni -- France\\1 vote négatif = résolution bloquée};
\draw[-{Stealth}, poporange, line width=1.2pt] (veto.west) -- (cse.east);
\end{tikzpicture}
\end{center}
\legende{Organigramme simplifié des six organes principaux de l'ONU. Le Conseil de sécurité, seul organe pouvant imposer des décisions contraignantes (sanctions, interventions), est dominé par les cinq membres permanents disposant du droit de veto.}
\end{popfigure}

\begin{methode}[Comparer la SDN et l'ONU (démarche en 4 étapes)]
Pour construire une comparaison rigoureuse au Baccalauréat :
\begin{enumerate}
\item \textbf{Universalité} : la SDN compte une soixantaine de membres, sans les États-Unis ; l'ONU part de 51 membres et devient quasi universelle (193 États).
\item \textbf{Pouvoir coercitif} : la SDN exige l'unanimité et ne peut rien imposer ; le Conseil de sécurité de l'ONU vote à la majorité (sous réserve du veto) et peut décider sanctions et opérations militaires.
\item \textbf{Absence d'armée propre} : ni la SDN ni l'ONU ne disposent de forces permanentes ; l'ONU s'appuie sur des contingents fournis par les États (les Casques bleus).
\item \textbf{Rôle du Secrétaire général} : chez l'ONU, il dispose d'une initiative diplomatique reconnue (article 99), bien supérieure à celle de son homologue de la SDN.
\end{enumerate}
Conclure en montrant que l'ONU corrige les défauts de la SDN sans les supprimer tous (le veto peut paralyser, comme l'unanimité autrefois).
\end{methode}

\begin{savaistu}[Le Cameroun et la Charte de l'ONU]
Le Cameroun, bien que non indépendant en 1945, est lié à la Charte de l'ONU par l'intermédiaire de la \textbf{France}, puissance administrante signataire. Placé sous le régime international de la \textbf{tutelle} (succédant au mandat de la SDN), le Cameroun relève du Conseil de tutelle : la France doit rendre compte chaque année à l'ONU de son administration. C'est d'ailleurs devant l'Assemblée générale de l'ONU que se jouera une partie du combat pour l'indépendance du Cameroun (pétitions de l'UPC dès 1948).
\end{savaistu}

\begin{exoresolu}[Exercice : les fondateurs et l'élargissement de l'ONU]
Énoncé. L'ONU compte 51 États fondateurs en 1945 et 193 États membres aujourd'hui. 1) Calculez le taux d'augmentation du nombre de membres. 2) Expliquez en quelques lignes la principale cause historique de cette progression.
\tcblower
Solution. 1) Le taux de variation se calcule ainsi :
\[ t = \frac{193 - 51}{51} \times 100 = \frac{142}{51} \times 100 \approx 278\,\%. \]
Le nombre de membres a donc été presque multiplié par quatre (augmentation d'environ 278\,\%). 2) Cette progression s'explique essentiellement par la \textbf{décolonisation} : à partir des années 1950-1960, les anciennes colonies d'Asie et d'Afrique accèdent à l'indépendance et adhèrent massivement à l'ONU (17 États africains entrent rien qu'en 1960, « année de l'Afrique »). L'ONU devient ainsi le porte-parole du Tiers-Monde, ce qui transforme ses équilibres internes.
\end{exoresolu}

\courssub{De l'alliance à la rupture : les débuts de la Guerre froide}

\textbf{Intuition.} Les Alliés de 1945 n'ont en commun que leur ennemi. Une fois Hitler vaincu, les deux superpuissances que sont devenus les États-Unis et l'URSS découvrent l'incompatibilité de leurs modèles : démocratie libérale et capitalisme contre parti unique et économie collectivisée. En quelques années, l'alliance se transforme en affrontement mondial — sans guerre directe : c'est la \cle{Guerre froide}.

Les étapes de la rupture sont rapides :
\begin{itemize}
\item \textbf{1946} : à Fulton (Missouri), Churchill dénonce dans un discours célèbre un « \cle{rideau de fer} » qui « est descendu à travers le continent », de Stettin à Trieste, séparant l'Europe libre de l'Europe soviétisée.
\item \textbf{1947} : le président Truman énonce la \cle{doctrine Truman} d'« endiguement » (\emph{containment}) du communisme, en aidant les États menacés (Grèce, Turquie). La même année, le \cle{plan Marshall} offre une aide économique massive (environ 13 milliards de dollars) à la reconstruction de l'Europe occidentale — aide refusée par l'URSS, qui y voit un instrument d'influence américaine.
\item \textbf{1947} : en riposte, Jdanov énonce la \cle{doctrine Jdanov} (« deux camps » : impérialiste et anti-impérialiste) et Staline crée le \textbf{Kominform} pour encadrer les partis communistes européens.
\end{itemize}

Le monde se \cle{bipolarise} : à l'ouest, le bloc américain (future OTAN, 1949) ; à l'est, le bloc soviétique (futur pacte de Varsovie, 1955). L'Allemagne en devient le symbole : le blocus de Berlin (1948-1949) puis la création de deux États allemands (RFA et RDA, 1949) concrétisent la coupure. La guerre chaude a laissé place à une guerre froide qui structurera le monde jusqu'en 1991.

\begin{aretenir}[L'essentiel de la section]
\begin{itemize}
\item 1945 : 60 millions de morts (2/3 de civils), plus de 1000 milliards de dollars de destructions, 12 millions d'Allemands expulsés.
\item Nuremberg (1945-1946) et Tokyo inventent les crimes contre la paix, de guerre et contre l'humanité ; la dénazification reste inachevée.
\item Yalta et Potsdam déplacent la Pologne vers l'ouest, divisent l'Allemagne et Berlin en quatre zones, et installent l'Europe de l'Est dans l'orbite soviétique.
\item L'ONU naît à San Francisco (juin 1945, 51 membres fondateurs) ; son Conseil de sécurité est dominé par les cinq membres permanents dotés du veto.
\item Dès 1946-1947 (Fulton, doctrine Truman, plan Marshall, doctrine Jdanov), la rupture Est-Ouest ouvre la Guerre froide et bipolarise le monde.
\end{itemize}
\end{aretenir}

\begin{exercice}[Pour s'entraîner]
1) Définissez le crime contre l'humanité et expliquez en quoi il constitue une innovation juridique. 2) Sur un croquis, placez les quatre zones d'occupation de l'Allemagne et la position de Berlin. 3) Rédigez un paragraphe argumenté répondant à la question : « L'ONU a-t-elle vraiment corrigé les défauts de la SDN ? »
\end{exercice}

\coursec{Héritages et mémoires : Décolonisation, Droits de l'Homme, Construction européenne}

La signature de la Charte de San Francisco en juin 1945 et la mise en place des organes de l'ONU marquent la fin institutionnelle de la guerre, mais elles ouvrent une ère nouvelle dont les secousses redessinent la carte du monde bien au-delà de 1945. L'ordre international issu de Yalta et de Potsdam porte en lui les germes de sa propre transformation : le principe d'égalité souveraine des États (Art. 2.1 de la Charte) entre en contradiction flagrante avec la réalité des empires coloniaux ; le procès de Nuremberg pose les fondations d'une justice pénale internationale que la Guerre froide va geler pendant cinquante ans ; le traumatisme du continent européen pousse la France et l'Allemagne de l'Ouest à inventer une forme inédite d'organisation politique, la construction supranationale. Cette section analyse comment la Seconde Guerre mondiale agit comme un \cle{révélateur} et un \cle{accélérateur} de l'histoire contemporaine, en structurant trois héritages majeurs : la décolonisation, le régime universel des droits de l'homme et la construction européenne, tout en interrogeant les mémoires conflictuelles qui en découlent.

\courssub{L'accélération des décolonisations : l'effondrement de la légitimité impériale}

\begin{popfigure}
\begin{tikzpicture}[scale=0.9, every node/.style={font=\small}]
% Axe principal
\draw[->, thick, popdark] (0,0) -- (15.5,0) node[right] {Années};
% Graduations
\foreach \x/\year in {0/1945, 1.5/1947, 3/1948, 4.5/1950, 6/1955, 7.5/1957, 9/1958, 10.5/1960, 12/1962, 13.5/1965, 15/1975} {
    \draw (\x,0.1) -- (\x,-0.1) node[below=2pt] {\year};
}
% Événements majeurs (au-dessus)
\coordinate (A) at (1.5,1.5); \node[above, align=center] at (A) {\textbf{Inde/Pakistan}\\(1947)}; \draw[popblue, thick] (1.5,0) -- (A);
\coordinate (B) at (3,2.2); \node[above, align=center] at (B) {\textbf{Création d'Israël}\\(1948)}; \draw[popblue, thick] (3,0) -- (B);
\coordinate (C) at (4.5,1.5); \node[above, align=center] at (C) {\textbf{DUDH}\\(1948)}; \draw[poppurple, thick] (4.5,0) -- (C);
\coordinate (D) at (6,2.2); \node[above, align=center] at (D) {\textbf{Conf. Bandung}\\(1955)}; \draw[popgreen, thick] (6,0) -- (D);
\coordinate (E) at (7.5,1.5); \node[above, align=center] at (E) {\textbf{Ghana}\\(1957)}; \draw[poporange, thick] (7.5,0) -- (E);
\coordinate (F) at (9,2.2); \node[above, align=center] at (F) {\textbf{Guinée « Non »}\\(1958)}; \draw[poporange, thick] (9,0) -- (F);
\coordinate (G) at (10.5,1.5); \node[above, align=center, popred] at (G) {\textbf{Cameroun}\\(1er janv 1960)}; \draw[popred, very thick] (10.5,0) -- (G);
\coordinate (H) at (12,2.2); \node[above, align=center] at (H) {\textbf{Algérie}\\(1962)}; \draw[popblue, thick] (12,0) -- (H);
\coordinate (I) at (13.5,1.5); \node[above, align=center] at (I) {\textbf{Fin Portugal}\\(1975)}; \draw[popblue, thick] (13.5,0) -- (I);
% Flux Vétérans -> Syndicats -> Partis (en dessous)
\draw[popgold, very thick, ->] (0.5,-1.2) -- (3,-1.2) node[midway, below, align=center] {\textbf{Vétérans}\\\textbf{africains}};
\draw[popgold, very thick, ->] (3,-1.2) -- (6,-1.2) node[midway, below, align=center] {\textbf{Syndicats}\\\textbf{(CGT, UGTAN)}};
\draw[popgold, very thick, ->] (6,-1.2) -- (10.5,-1.2) node[midway, below, align=center] {\textbf{Partis nationaux}\\\textbf{(RDA, UPC, CPP)}};
\draw[popgold, very thick, ->] (10.5,-1.2) -- (13,-1.2) node[midway, below, align=center] {\textbf{Indépendances}};
% Annotation Charte de l'Atlantique
\node[align=center, text width=4cm, fill=popblueL, rounded corners, draw=popblue] at (0.75, -3) {\textbf{Charte de l'Atlantique (1941)}\\\textit{Art. 3 : Droit des peuples à choisir leur forme de gouvernement}};
\draw[popblue, dashed] (0.75,-2.5) -- (0.75,-0.3);
% Annotation Chap XI-XII ONU
\node[align=center, text width=4cm, fill=poppurpleL, rounded corners, draw=poppurple] at (5.25, -3) {\textbf{Charte ONU Chap. XI-XII}\\\textit{Principes de tutelle et d'autogouvernement}};
\draw[poppurple, dashed] (5.25,-2.5) -- (5.25,-0.3);
\end{tikzpicture}
\legende{Frise chronologique : L'accélération des décolonisations (1945-1975) et le rôle moteur des vétérans africains. En rouge, l'indépendance du Cameroun (1er janvier 1960), premier État d'Afrique subsaharienne francophone à accéder à la souveraineté internationale.}
\end{popfigure}

La Seconde Guerre mondiale a agi comme un \cle{choc structurel} sur le système colonial. Trois mécanismes conjugués expliquent l'effondrement rapide des empires entre 1945 et 1960.

\begin{definition}[Décolonisation]
Processus historique (1945-1975) par lequel les territoires colonisés accèdent à la souveraineté internationale reconnue. Il ne se réduit pas au transfert formel des pouvoirs : il implique une réappropriation de l'État, de l'économie et de l'histoire par les peuples anciennement dominés, souvent au prix de luttes armées ou de négociations sous contrainte.
\end{definition}

\paragraph{1. L'affaiblissement matériel et moral des métropoles.} La France, les Pays-Bas, la Belgique et le Royaume-Uni sortent exsangues du conflit. Leur prestige est brisé par les défaites de 1940 (France, Pays-Bas) ou la chute de Singapour (1942) qui a pulvérisé le mythe de l'invincibilité blanche face au Japon. La métropole n'est plus le centre du monde ; elle devient une périphérie dépendante de l'aide américaine (Plan Marshall). Maintenir l'ordre colonial coûte trop cher à des économies en reconstruction.

\paragraph{2. Le discrédit idéologique du colonialisme.} La Charte de l'Atlantique (août 1941), signée par Churchill et Roosevelt, proclame le « droit de tous les peuples à choisir la forme de gouvernement sous laquelle ils veulent vivre » (Art. 3). Bien que Churchill ait tenté de l'exclure de l'Empire britannique, le texte devient la référence morale des nationalistes (Hô Chi Minh, Nkrumah, Senghor). La Charte de l'ONU (Chap. XI « Déclaration relative aux territoires non autonomes » et Chap. XII « Système de tutelle ») institutionnalise l'obligation de mener les territoires vers l'autogouvernement. Le colonialisme devient une anomalie juridique dans l'ordre international nouveau.

\paragraph{3. Le rôle décisif des vétérans africains.} Près de \textbf{800 000 Africains} ont combattu sous drapeaux français, britanniques, belges. Au Cameroun, on estime à \textbf{15 000 - 20 000} le nombre de tirailleurs engagés (dont la majorité dans le Corps expéditionnaire français en Italie, en Provence, en Allemagne). Leur retour au pays opère une \cle{révolution des mentalités} :
\begin{itemize}
    \item \textbf{Expérience de l'égalité :} Au front, le tirailleur mange, dort, se bat et meurt aux côtés du métropolitain. La hiérarchie raciale s'efface face à la mort commune.
    \item \textbf{Acquisition de compétences :} Maniement d'armes modernes, logistique, secourisme, alphabétisation militaire.
    \item \textbf{Revendications citoyennes :} La promesse d'égalité (Conférence de Brazzaville 1944, Loi Lamine Guèye 1946) se heurte à la réalité : pensions inégales (cristallisées en 1959, gelées jusqu'aux lois de 2007 et 2010 pour la « cristallisation » puis la revalorisation), refus du droit de vote plénier, répression des grèves (grève des cheminots au Sénégal 1947, insurrection malgache 1947, répression au Cameroun dès 1955).
\end{itemize}

\begin{exemplebox}[Exemple chiffré : Le sort des vétérans camerounais]
Sur les \num{15000} à \num{20000} Camerounais mobilisés (chiffres variables selon sources : archives militaires vs mémoire orale), beaucoup sont démobilisés sans qualification civile reconnue. La \cle{pension militaire} versée aux anciens combattants africains est fixée en 1959 à un taux \textbf{inférieur de 50 à 75 \%} à celle des métropolitains pour un même grade. Ce n'est qu'en \textbf{2007} (loi sur la cristallisation) puis \textbf{2010} (revalorisation) que l'égalité formelle est rétablie, sans effet rétroactif complet. Cette injustice alimente le ressentiment nationaliste : l'UPC (Union des Populations du Cameroun) recrute massivement parmi ces vétérans « oubliés » (ex: \textbf{Ernest Ouandié}, ancien tirailleur).
\end{exemplebox}

\paragraph{La dynamique « Vétérans $\to$ Syndicats $\to$ Partis ».} La frise ci-dessus illustre ce continuum. Les vétérans forment le noyau dur des premiers syndicats autonomes (CGTA, UGTAN) qui structurent la grève générale de 1945-1946 (Dakar, Douala). Ces syndicats deviennent les bras armés des partis nationalistes (RDA en AOF, CPP au Ghana, UPC au Cameroun). L'indépendance du Ghana (1957) et le « Non » de la Guinée (1958) ouvrent la voie à l'« Année de l'Afrique » (1960), où le Cameroun, par sa date symbolique du \textbf{1er janvier 1960}, devient le \cle{premier État d'Afrique subsaharienne francophone} à accéder à la souveraineté internationale.

\begin{savaistu}[Le Cameroun : une indépendance née de la guerre]
Le Cameroun est un cas d'école : mandat de la SDN puis tutelle de l'ONU (Chap. XII), il devait théoriquement préparer l'autonomie sous contrôle international. La frustration post-1945 (interdiction de l'UPC en 1955, guerre de maquis dans le Sanaga-Maritime, le Bamiléké et le Nord, répression française) transforme le processus légal en \textbf{guerre de libération nationale} (1955-1971). L'indépendance du 1er janvier 1960, négociée par Ahmadou Ahidjo, ne règle pas le conflit : l'UPC refuse de déposer les armes, considérant le nouveau pouvoir comme néocolonial. La « réunification » avec le Southern British Cameroons (1961) parachève le territoire mais laisse une mémoire blessée, encore vive aujourd'hui (crise anglophone depuis 2016).
\end{savaistu}

\courssub{La Déclaration universelle des Droits de l'Homme (1948) : la réponse juridique à la barbarie}

\begin{popfigure}
\begin{tikzpicture}[scale=0.95, every node/.style={font=\small, align=center}]
% Nœuds principaux
\node[draw=popblue, thick, fill=popblueL, rounded corners, minimum width=3.5cm, minimum height=1.2cm, align=center] (G1864) at (0,3){\textbf{Convention de Genève 1864}\\\textit{Protection blessés armées}};
\node[draw=popblue, thick, fill=popblueL, rounded corners, minimum width=3.5cm, minimum height=1.2cm, align=center] (G1949) at (0,1){\textbf{Conventions de Genève 1949 (4)}\\\textit{DIH : blessés, prisonniers, civils}};
\node[draw=poporange, thick, fill=poporangeL, rounded corners, minimum width=4cm, minimum height=1.2cm, align=center] (Nure) at (6,2){\textbf{Procès de Nuremberg 1945-46}\\\textit{Crimes contre la paix, de guerre, contre l'humanité}\\\textit{\textit{Juridiction pénale individuelle}}};
\node[draw=poppurple, thick, fill=poppurpleL, rounded corners, minimum width=4cm, minimum height=1.2cm, align=center] (DUDH) at (12,2){\textbf{DUDH 1948}\\\textit{Universalité, Indivisibilité, Inconditionnalité}\\\textit{\textit{Base morale \& juridique moderne}}};
\node[draw=popred, thick, fill=popredL, rounded corners, minimum width=3.5cm, minimum height=1.2cm, align=center] (CPI) at (18,2){\textbf{Statut de Rome 1998}\\\textit{CPI (2002)}\\\textit{\textit{Juridiction permanente}}};
% Flèches
\draw[->, thick, popblue] (G1864) -- (G1949) node[midway, left] {Évolution DIH};
\draw[->, thick, popblue] (G1949) -- node[above, sloped, pos=0.3] {Codification post-guerre} (Nure);
\draw[->, thick, poporange] (Nure) -- node[above] {Précédent jurisprudentiel} (DUDH);
\draw[->, thick, poppurple] (DUDH) -- node[above] {Inspiration normative} (CPI);
\draw[->, thick, dashed, popgreen] (G1949) to[bend right=20] node[below, sloped] {Droit humanitaire $\to$ Droit pénal} (CPI);
% Annotations
\node[fill=popcream, draw=popmuted, rounded corners, text width=3.5cm, font=\footnotesize] at (3,4) {\textbf{Rappel :} Le DIH (Genève) régit la \textit{conduite} des hostilités (\textit{jus in bello}). Nuremberg/DUDH/CPI sanctionnent les \textit{crimes} et protègent la personne (\textit{jus ad bellum} / Droit int. des DH).};
\draw[popmuted, ->] (3,3.6) -- (G1949.north);
\draw[popmuted, ->] (3,3.6) -- (Nure.north);
\end{tikzpicture}
\legende{Généalogie du Droit International Humanitaire (DIH) et des Droits de l'Homme : de la protection des combattants (Genève 1864) à la responsabilité pénale individuelle permanente (CPI 2002), en passant par la rupture normative de Nuremberg et la consécration universelle de 1948.}
\end{popfigure}

Si la Charte de l'ONU mentionne les droits de l'homme (Préambule, Art. 1, 55, 56), elle ne les \textit{définit} pas. L'horreur découverte dans les camps (Shoah, camps de concentration, expériences médicales nazies) impose une codification urgente. La \cle{Déclaration universelle des droits de l'homme (DUDH)}, adoptée le \textbf{10 décembre 1948} à Paris (Palais de Chaillot) par 48 voix pour, 0 contre, 8 abstentions (bloc soviétique, Afrique du Sud, Arabie saoudite), constitue l'acte fondateur du droit international moderne des droits de l'homme.

\begin{definition}[DUDH - Trois piliers]
\begin{itemize}
    \item \textbf{Universalité :} S'applique à tout être humain, « sans distinction de race, de sexe, de langue, de religion, d'opinion politique ou de toute autre opinion » (Art. 2).
    \item \textbf{Indivisibilité :} Droits civils/politiques (Art. 3-21) et droits économiques/sociaux/culturels (Art. 22-27) forment un tout inséparable.
    \item \textbf{Inconditionnalité :} Nul ne peut s'en prévaloir pour détruire les droits d'autrui (Art. 29-30) ; ils ne dépendent pas de la citoyenneté nationale.
\end{itemize}
\end{definition}

\paragraph{Genèse et acteurs clés.} La Commission des droits de l'homme (ECOSOC), présidée par \textbf{Eleanor Roosevelt} (USA), avec \textbf{René Cassin} (France, prix Nobel 1968) comme vice-président et rapporteur, \textbf{Peng Chun Chang} (Chine) et \textbf{Charles Malik} (Liban), rédige le texte en deux ans (1947-1948). Le compromis Est/Ouest porte sur la primauté de l'individu (Occident) vs les devoirs envers la communauté (URSS) ; le texte final équilibre les deux.

\paragraph{De la déclaration au droit contraignant.} La DUDH n'est qu'une \textbf{résolution d'Assemblée générale} (non contraignante \textit{stricto sensu}), mais elle a acquis force coutumière. Elle est complétée par deux \textbf{Pactes de 1966} (entrés en vigueur 1976) :
\begin{itemize}
    \item PIDCP (Droits civils et politiques) : Comité des droits de l'homme (rapports, communications individuelles).
    \item PIDESC (Droits économiques, sociaux, culturels) : Comité DESC (rapports, procédure de communications 2008).
\end{itemize}
Ensemble, ils forment la \cle{Charte internationale des droits de l'homme}. La Conférence de Vienne (1993) réaffirme leur universalité face au relativisme culturel.

\begin{propriete}[Exigibilité au Probatoire]
La distinction entre \textbf{DUDH (1948, déclaration, valeur morale/coutumière)} et \textbf{Pactes 1966 (traité, valeur juridique contraignante pour les États parties)} est un classique de l'examen. Savoir citer 3 articles emblématiques (Art. 1 : dignité ; Art. 3 : vie/liberté/sûreté ; Art. 18 : liberté pensée/conscience/religion ; Art. 19 : opinion/expression).
\end{propriete}

\begin{attention}[Piège fréquent]
Ne pas confondre \textbf{Droit International Humanitaire (DIH / Conventions de Genève)} qui s'applique \textit{en temps de guerre} (\textit{jus in bello}) et \textbf{Droit International des Droits de l'Homme (DIDH)} qui s'applique \textit{en tout temps} (paix et guerre), avec des régimes de dérogation différents. Nuremberg fait le pont : il juge des crimes de guerre (DIH) et des crimes contre l'humanité (DIDH).
\end{attention}

\courssub{La construction européenne : la paix par le droit et le marché}

\begin{popfigure}
\begin{tikzpicture}[scale=0.85, every node/.style={font=\small, align=center}]
% Étapes majeures
\node[draw=popblue, thick, fill=popblueL, rounded corners, minimum width=3.2cm, minimum height=1.1cm, align=center] (Schuman) at (0,3){\textbf{Déclaration Schuman}\\\textit{9 mai 1950}\\\textit{Charbon \& Acier}};
\node[draw=popblue, thick, fill=popblueL, rounded corners, minimum width=3.2cm, minimum height=1.1cm, align=center] (CECA) at (4,3){\textbf{CECA 1951}\\\textit{6 États : F, RFA, I, BNL}\\\textit{Haute Autorité (Supranational)}};
\node[draw=popgreen, thick, fill=popgreenL, rounded corners, minimum width=3.2cm, minimum height=1.1cm, align=center] (CEE) at (8,3){\textbf{Traités de Rome 1957}\\\textit{CEE + Euratom}\\\textit{Marché commun, PAC}};
\node[draw=poporange, thick, fill=poporangeL, rounded corners, minimum width=3.2cm, minimum height=1.1cm, align=center] (UE) at (12,3){\textbf{Traité de Maastricht 1992}\\\textit{Union Européenne (UE)}\\\textit{Piliers : Éco, PESC, JAI}};
\node[draw=poppurple, thick, fill=poppurpleL, rounded corners, minimum width=3.2cm, minimum height=1.1cm, align=center] (Lisbonne) at (16,3){\textbf{Traité de Lisbonne 2007}\\\textit{Personnalité juridique}\\\textit{Charte droits fondamentale}};
% Flèches principales
\draw[->, thick, popblue] (Schuman) -- node[above] {Traité Paris} (CECA);
\draw[->, thick, popgreen] (CECA) -- node[above] {Relance} (CEE);
\draw[->, thick, poporange] (CEE) -- node[above] {Élargissement \& Approfondissement} (UE);
\draw[->, thick, poppurple] (UE) -- node[above] {Réforme institutionnelle} (Lisbonne);
% Acteurs / Moteurs (en dessous)
\node[draw=popgold, thick, fill=popgoldL, rounded corners, minimum width=3.5cm, minimum height=1.5cm, align=center] (Monnet) at (2,0.5){\textbf{Jean Monnet}\\\textit{Architecte méthode fonctionnaliste}\\\textit{\textit{« Faire par le concret »}}};
\node[draw=popgold, thick, fill=popgoldL, rounded corners, minimum width=3.5cm, minimum height=1.5cm, align=center] (Schuman2) at (6,0.5){\textbf{Robert Schuman}\\\textit{Père de l'Europe}\\\textit{\textit{Réconciliation F-RFA}}};
\node[draw=popgold, thick, fill=popgoldL, rounded corners, minimum width=3.5cm, minimum height=1.5cm, align=center] (Adenauer) at (10,0.5){\textbf{Konrad Adenauer}\\\textit{Chancelier RFA}\\\textit{\textit{Ancrage occidental}}};
\node[draw=popgold, thick, fill=popgoldL, rounded corners, minimum width=3.5cm, minimum height=1.5cm, align=center] (Delors) at (14,0.5){\textbf{Jacques Delors}\\\textit{Président Commission}\\\textit{\textit{Acte unique, Euro}}};
\draw[popmuted, dashed] (Schuman) |- (Monnet);
\draw[popmuted, dashed] (Schuman) |- (Schuman2);
\draw[popmuted, dashed] (CECA) |- (Adenauer);
\draw[popmuted, dashed] (UE) |- (Delors);
% Annotation Guerre Froide
\node[fill=poppinkL, draw=poppink, rounded corners, text width=4cm, font=\footnotesize, align=center] at (8,5) {\textbf{Contexte Guerre Froide :}\\\textit{Plan Marshall (1947) $\to$ OECE $\to$ Intégration ouest-européenne sous parapluie US (OTAN 1949). RFA réarmée (1954) via UEO puis OTAN.}};
\draw[poppink, ->] (8,4.5) -- (CEE.north);
\end{tikzpicture}
\legende{Les grandes étapes de la construction européenne (1950-2009) : de la mise en commun du charbon et de l'acier (CECA) à l'Union politique (Lisbonne), portée par la réconciliation franco-allemande et la méthode fonctionnaliste (Monnet).}
\end{popfigure}

L'idée d'une Europe unie préexiste à 1939, mais la Seconde Guerre mondiale en fait une \cle{nécessité vitale} : empêcher le retour de la guerre franco-allemande (trois conflits en 70 ans : 1870, 1914, 1939) et reconstruire l'économie continentale face aux deux superpuissances (USA, URSS).

\paragraph{La méthode fonctionnaliste (Jean Monnet).} Pas de fédération d'emblée, mais des \textbf{réalisations concrètes} créant une \cle{solidarité de fait} : on commence par les secteurs stratégiques de la guerre (charbon, acier) pour rendre un conflit matériel impossible entre la France et la RFA.

\begin{enumerate}
    \item \textbf{1950-1951 : CECA (Communauté européenne du charbon et de l'acier).} Déclaration Schuman (9 mai 1950), Traité de Paris (1951). 6 États (France, RFA, Italie, Belgique, Pays-Bas, Luxembourg). \textbf{Haute Autorité} : premier organe \textbf{supranational} (décisions s'imposent aux États). Échec de la CED (Communauté européenne de défense, 1954) rejetée par l'Assemblée nationale française.
    \item \textbf{1957 : Traité de Rome $\to$ CEE + Euratom.} Marché commun (libre circulation marchandises, personnes, services, capitaux), Politique agricole commune (PAC). Institutions : Commission (exécutif), Conseil (législatif intergouvernemental), Assemblée (parlementaire), Cour de justice.
    \item \textbf{1986 : Acte unique européen.} Marché unique (1993), vote à la majorité qualifiée au Conseil (fin du veto systématique).
    \item \textbf{1992 : Traité de Maastricht $\to$ Union Européenne (UE).} Trois piliers : 1. Communautés (supranational), 2. PESC (Politique étrangère/sécurité, intergouvernemental), 3. JAI (Justice/Affaires intérieures). Citoyenneté européenne, critères de convergence pour l'Euro.
    \item \textbf{2002/2009 : Euro et Traité de Lisbonne.} Monnaie unique (zone euro), personnalité juridique unique, Charte des droits fondamentaux (valeur contraignante), Président du Conseil européen stable, Haut Représentant PESC.
\end{enumerate}

\begin{methode}[Analyse d'un traité européen (Bac)]
\begin{enumerate}
    \item \textbf{Contexte :} Crise ou relance ? (ex: relance post-CED pour Rome ; fin Guerre froide pour Maastricht).
    \item \textbf{Acteurs :} États moteurs (couple franco-allemand), Commission (Monnet, Delors), Parlement (élections directes 1979).
    \item \textbf{Innovation institutionnelle :} Supranationalité accrue ? (QMV, Parlement co-législateur, Cour de justice).
    \item \textbf{Portée sectorielle :} Économique seulement (Rome) ou politique (Maastricht, Lisbonne) ?
    \item \textbf{Limites :} « Démocratie déficitaire », élargissements sans approfondissement (2004, 2007), Brexit (2020).
\end{enumerate}
\end{methode}

\begin{aretenir}[Synthèse des trois héritages (1945 - aujourd'hui)]
\begin{center}
\begin{tabularx}{\linewidth}{|l|X|X|X|}\hline
\textbf{Héritage} & \textbf{Décolonisation} & \textbf{Droits de l'Homme} & \textbf{Construction européenne} \\ \hline
\textbf{Moteur principal} & Choc WW2 + Vétérans + ONU & Shoah + Nuremberg + ONU & Réconciliation F-RFA + Guerre Froide \\ \hline
\textbf{Acte fondateur} & Charte ONU (Chap. XI/XII) & DUDH (1948) & Déclaration Schuman (1950) \\ \hline
\textbf{Innovation juridique} & Droit des peuples à disposer d'eux-mêmes (Art. 1.2 Charte) & Universalité / Indivisibilité / Responsabilité individuelle (Nuremberg) & Supranationalité (Haute Autorité CECA $\to$ Commission UE) \\ \hline
\textbf{Cas camerounais} & Indépendance 1960 (Tutelle ONU) + Guerre UPC (1955-71) & Ratification Pactes 1984 ; Commission nationale DH 2004 & Partenariat ACP-UE (Lomé $\to$ Cotonou $\to$ Samoa) ; Accords d'association \\ \hline
\end{tabularx}
\end{center}
\end{aretenir}

\begin{exercice}[Sujet type Probatoire : Commentaire de documents]
\textbf{Document 1 :} Extrait de la Déclaration Schuman (9 mai 1950). « La paix mondiale ne saurait être sauvegardée sans des efforts créateurs à la mesure des dangers qui la menacent. L'Europe ne se fera pas d'un coup, ni dans une construction d'ensemble : elle se fera par des réalisations concrètes créant d'abord une solidarité de fait. La mise en commun des productions de charbon et d'acier [...] changera les destinées de ces régions qui depuis longtemps vouées à la fabrication des armes de guerre [...] ».
\textbf{Document 2 :} Carte schématique de l'élargissement de la CEE/UE (1957-2013) : des 6 fondateurs à 28 membres (avant Brexit), avec flèches d'élargissement Sud (Grèce, Espagne, Portugal), Nord (UK, Danemark, Irlande), Est (ex-Pacte de Varsovie).
\textbf{Document 3 :} Extrait du discours de Robert Schuman à l'Assemblée de Strasbourg (1958) : « Nous entreprenons une œuvre de paix. [...] L'Europe que nous voulons, c'est une Europe qui assure la liberté, la prospérité et la paix. »

\textbf{Questions :}
\begin{enumerate}
    \item Situez le document 1 (auteur, date, contexte historique immédiat).
    \item Expliquez l'expression « solidarité de fait » en la reliant à la méthode fonctionnaliste.
    \item Montrez, à l'aide des trois documents, comment la construction européenne répond à un double objectif de \textbf{paix} et de \textbf{prospérité}.
    \item En vous appuyant sur vos connaissances, précisez les principales étapes de l'approfondissement institutionnel entre 1957 et 1992.
\end{enumerate}
\end{exercice}

\begin{corrige}[Exercice n°1 : Commentaire de documents]
\begin{enumerate}
    \item \textbf{Document 1 :} Robert Schuman, Ministre français des Affaires étrangères, déclaration du 9 mai 1950 (Salon de l'Horloge, Quai d'Orsay). Contexte : Guerre froide (blocus de Berlin 1948-49, guerre de Corée 1950), échec imminent de la CED (rejetée 1954), nécessité d'ancrer la RFA (créée 1949) dans l'Ouest. Inspiré par Jean Monnet.
    \item \textbf{« Solidarité de fait » :} Concept monnétien. L'intégration ne se décrète pas par un traité constitutionnel (méthode fédéraliste), elle se \textit{construit} par des actes concrets (marché commun charbon/acier) qui lient économiquement les États (interdépendance), rendant la guerre matériellement impossible et politiquement impensable. C'est le pragmatisme du « faire par le concret ».
    \item \textbf{Paix :} Doc 1 (« paix mondiale », « régions vouées à la fabrication des armes »), Doc 3 (« œuvre de paix », « assure la paix »). Réconciliation franco-allemande (Schuman/Adenauer). Doc 2 : Élargissement aux pays du Sud (sortie dictatures : Grèce 74, Espagne/Portugal 75) et de l'Est (sortie communisme 89) = consolidation paix continentale.
    \item \textbf{Prospérité :} Doc 1 (« prospérité » implicite via modernisation industrie), Doc 3 (« prospérité »). Doc 2 : Marché unique (Acte unique 86), PAC, Fonds structurels (cohesion). Étapes : Traité Rome 57 (CEE, marché commun) $\to$ Fusion exécutifs 67 $\to$ Acte unique 86 (Marché unique, QMV) $\to$ Maastricht 92 (UE, Euro, Citoyenneté, PESC).
\end{enumerate}
\end{corrige}

\newpage

\coursec{Activité d'intégration}

\courssub{Objectifs de l'activité}

Cette activité d'intégration clôt la leçon sur la Deuxième Guerre mondiale en replaçant l'élève dans la posture de l'historien : confronter des sources hétérogènes, croiser les points de vue, distinguer le discours officiel de la réalité vécue. Elle mobilise l'ensemble des savoirs de l'unité (genèse du conflit, violence de guerre, rôle de l'Afrique et du Cameroun, conséquences de 1945) et les deux savoir-faire du programme : \cle{appliquer les notions de l'unité à des situations concrètes} et \cle{rédiger et justifier une démarche, une réponse ou une conclusion}.

À l'issue de l'activité, l'élève doit être capable de :
\begin{itemize}
\item identifier la nature, l'origine et l'intention d'un document historique ;
\item évaluer la fiabilité d'une source en la croisant avec d'autres ;
\item construire un tableau comparatif synthétique ;
\item calculer et interpréter un taux de variation à partir de données statistiques ;
\item rédiger une réponse structurée et argumentée à une question problématisée ;
\item participer à un débat réglé sur une question de mémoire contemporaine.
\end{itemize}

\courssub{Situation}

\textbf{Enquête historienne : « Le Cameroun, base arrière et fournisseur d'hommes : réalité vécue vs discours officiel ».}

Tu es élève en classe de Terminale technique au lycée de Nkongsamba. Dans le cadre de la semaine de la mémoire organisée par le club d'histoire, ton enseignant te confie, avec ton groupe, un dossier d'archives reconstitué portant sur l'engagement du Cameroun dans la Deuxième Guerre mondiale. La question qui guide l'enquête est la suivante : le Cameroun est présenté dans les discours officiels de l'époque comme un territoire « volontairement » rallié à la France libre et « uni » à l'Empire pour la victoire. Mais que disent réellement les documents ?

\textbf{Le dossier documentaire (6 documents) :}

\textbf{Document 1 — Décret de ralliement (août 1940).} Extrait du texte par lequel le gouverneur et les notables du Cameroun français, à la suite de l'appel du général de Gaulle du 18 juin 1940 et de l'action de Philippe Leclerc de Hauteclocque débarqué à Douala, proclament le ralliement du territoire à la France libre. Le texte célèbre « l'adhésion spontanée des populations » à la cause de la libération.

\textbf{Document 2 — Tableau statistique : production et exportations du Cameroun, 1939-1945.}

\begin{center}
\begin{tabularx}{\linewidth}{|l|X|X|X|}\hline
\rowcolor{popblueL} \textbf{Produit} & \textbf{1939} & \textbf{1943} & \textbf{1945} \\ \hline
Caoutchouc (tonnes) & 2\,100 & 7\,800 & 6\,500 \\ \hline
Cacao (tonnes) & 38\,000 & 24\,000 & 21\,000 \\ \hline
Bois (grumes, m\textsuperscript{3}) & 45\,000 & 92\,000 & 70\,000 \\ \hline
Palmistes/huile de palme (tonnes) & 30\,000 & 41\,000 & 36\,000 \\ \hline
\end{tabularx}
\end{center}

Source : données administratives coloniales reconstituées (valeurs indicatives). Le caoutchouc et le bois, matières stratégiques pour l'effort de guerre, explosent pendant le conflit ; le cacao, culture commerciale paysanne, chute faute de débouchés et de main-d'œuvre.

\textbf{Document 3 — Rapport d'un administrateur de circonscription (1943).} « Les réquisitions de vivres (ignames, maïs, arachides) destinées au ravitaillement des troupes et des chantiers ont été renouvelées. Les prestations de portage et de travail sur les pistes et les plantations de caoutchouc pèsent lourdement sur les villages. Plusieurs chefs signalent des difficultés de subsistance et le départ des jeunes gens vers la brousse pour échapper aux corvées. »

\textbf{Document 4 — Lettre censurée d'un tirailleur camerounais, Provence, septembre 1944.} « Nous avons débarqué avec les troupes françaises. Le froid, la neige, nous ne connaissions pas. Beaucoup de camarades sont tombés. On nous dit que nous combattons pour la liberté, mais notre solde n'est pas celle des soldats français, et nos familles écrivent que les réquisitions continuent au pays. »

\textbf{Document 5 — Affiche de propagande (1943-1944).} « L'Empire français uni pour la victoire » : on y voit des soldats de toutes les colonies, dont un tirailleur africain, marchant côte à côte derrière le drapeau français et la croix de Lorraine, avec les mots « Union », « Victoire », « Volontariat ».

\textbf{Document 6 — Témoignage oral transcrit (recueilli en 2005 auprès d'un vétéran de Bafoussam).} « Je suis parti en 1942. On disait que c'était pour la France. Certains sont partis volontairement, d'autres ont été désignés par les chefs. À mon retour, j'ai attendu des années une pension qui n'arrivait pas, ou qui était plus petite que celle des anciens combattants français. C'est seulement très tard qu'on a parlé d'égalité. »

\courssub{Consignes}

Par groupes de 4, tu réaliseras les tâches suivantes, dans l'ordre :

\begin{enumerate}
\item \textbf{Classer et critiquer.} Classe les six documents par nature (document officiel, document statistique, document intime/privé, document de propagande, témoignage) et remplis la grille d'analyse fournie : auteur, date, destinataire, intention, fiabilité, limites.
\item \textbf{Croiser les sources.} Construis un tableau comparatif à deux colonnes : « Discours officiel (Victoire, Union, Volontariat) » et « Réalité vécue (Contrainte, Pénurie, Inégalité, Séparation) ». Chaque case doit être justifiée par au moins un document précis.
\item \textbf{Exploiter les chiffres.} À partir du document 2, calcule le taux de variation de la production de caoutchouc et de cacao entre 1939 et 1943, et explique ce que ces chiffres révèlent de la place du Cameroun dans l'effort de guerre.
\item \textbf{Rédiger.} Rédige collectivement une réponse structurée (introduction, deux parties argumentées, conclusion) à la question : « Dans quelle mesure l'engagement du Cameroun dans la Seconde Guerre mondiale relève-t-il d'une adhésion libre ou d'une contrainte coloniale réactivée ? »
\item \textbf{Présenter et débattre.} Présente la conclusion de ton groupe à la classe, puis participe au débat final : quelle place pour la mémoire des vétérans africains aujourd'hui (commémorations, pensions, reconnaissance) ?
\end{enumerate}

\courssub{Ressources à mobiliser}

\begin{aretenir}[Ce que tu dois réactiver]
\begin{itemize}
\item \textbf{Méthode d'analyse de document} : identifier la nature, l'auteur, le contexte, le destinataire et l'intention ; distinguer le fait de l'opinion ; croiser les sources pour vérifier une information.
\item \textbf{Savoirs de la leçon} : le ralliement du Cameroun à la France libre (août 1940, action de Leclerc) ; le Cameroun comme base arrière et réservoir d'hommes et de matières premières ; les tirailleurs dans les campagnes d'Afrique du Nord, d'Italie et de Provence ; les conséquences de la guerre (affaiblissement des puissances coloniales, création de l'ONU en 1945, aspirations nationalistes).
\item \textbf{Outil statistique} : le taux de variation se calcule ainsi :
\[ \text{Taux de variation} = \frac{V_{\text{finale}} - V_{\text{initiale}}}{V_{\text{initiale}}} \times 100 \]
\item \textbf{Méthode de la réponse argumentée} : introduction (accroche, problématique, annonce du plan), deux parties équilibrées, conclusion ouverte.
\end{itemize}
\end{aretenir}

\courssub{Démarche et corrigé commenté}

\begin{exoresolu}[Corrigé commenté de l'enquête]
\textbf{Étape 1 — Classement et critique des documents.}

Doc 1 (décret) : document \cle{officiel}, émanant de l'autorité coloniale ; intention de légitimer le ralliement ; le mot « spontané » doit être interrogé. Doc 2 (statistiques) : source administrative, factuelle mais à croiser (qui a produit, dans quelles conditions de travail ?). Doc 3 (rapport) : source officielle interne, donc plus franche que le discours public ; elle révèle la contrainte. Doc 4 (lettre) : source \cle{intime}, ayant échappé partiellement à la censure, donc précieuse car elle dit le vécu malgré le contrôle. Doc 5 (affiche) : document de \cle{propagande} ; il ne décrit pas la réalité, il cherche à la fabriquer. Doc 6 (témoignage) : source orale tardive (2005) ; risque de déformation de la mémoire, mais irremplaçable sur les pensions et la reconnaissance.

\tcblower

\textbf{Étape 2 — Tableau comparatif (exemple de résultat attendu).}

\begin{center}
\begin{tabularx}{\linewidth}{|X|X|}\hline
\rowcolor{popblueL} \textbf{Discours officiel} & \textbf{Réalité vécue} \\ \hline
Adhésion « spontanée » des populations (doc 1) & Ralliement décidé par l'autorité coloniale et les militaires, sans consultation des populations (docs 1 et 3 croisés) \\ \hline
« Union » de l'Empire pour la victoire (doc 5) & Inégalité de solde et de pensions entre tirailleurs africains et soldats français (docs 4 et 6) \\ \hline
« Volontariat » des combattants (doc 5) & Désignation par les chefs, pressions, réquisitions et corvées (docs 3 et 6) \\ \hline
Effort partagé et récompensé (doc 5) & Pénuries vivrières, fuite des jeunes vers la brousse, séparation des familles (docs 3 et 4) \\ \hline
\end{tabularx}
\end{center}

\textbf{Étape 3 — Exploitation statistique.}

Caoutchouc : $\dfrac{7800 - 2100}{2100} \times 100 = \dfrac{5700}{2100} \times 100 \approx +271\,\%$. Cacao : $\dfrac{24000 - 38000}{38000} \times 100 = \dfrac{-14000}{38000} \times 100 \approx -37\,\%$. Interprétation : la guerre réoriente l'économie camerounaise vers les matières \cle{stratégiques} (caoutchouc, bois) au détriment des cultures commerciales paysannes. Le Cameroun fonctionne comme une \cle{base arrière} : il fournit hommes, vivres et matières premières à l'effort de guerre de la France libre puis des Alliés.

\textbf{Étape 4 — Réponse structurée (trame rédigée).}

\emph{Introduction} : en août 1940, le Cameroun se rallie à la France libre et devient une base arrière de l'effort de guerre. Mais cet engagement fut-il le fruit d'une adhésion libre ou d'une contrainte coloniale réactivée ? Nous verrons d'abord qu'une adhésion réelle a existé, puis que la contrainte coloniale a dominé le vécu des populations.

\emph{Partie 1 — Une adhésion réelle mais encadrée.} Le ralliement d'août 1940 donne au Cameroun un rôle stratégique : territoire de la France libre, base de départ des opérations de Leclerc vers le Tchad et l'Afrique du Nord. Des volontaires s'engagent (doc 6 : « certains sont partis volontairement ») ; les tirailleurs participent aux campagnes de libération (doc 4, Provence 1944). Le discours de l'union (doc 5) traduit aussi une adhésion idéologique à la lutte contre le nazisme.

\emph{Partie 2 — Une contrainte coloniale réactivée.} Le ralliement est décidé d'en haut, sans consultation des populations (doc 1, discours de légitimation). L'effort de guerre repose sur les réquisitions, les corvées et le portage (doc 3), provoquant pénuries et fuites. Les inégalités persistent jusque dans l'armée : solde inférieure (doc 4), pensions différées et moindres (doc 6). La statistique montre une économie mise au service exclusif de la métropole (doc 2).

\emph{Conclusion} : l'engagement camerounais combine une adhésion partielle des élites et de volontaires et une contrainte coloniale massive pour la majorité. Cette contradiction nourrira, après 1945, les revendications nationalistes et le combat pour la reconnaissance des vétérans.

\textbf{Étape 5 — Débat de mémoire.} La classe confronte les conclusions : place des tirailleurs dans les commémorations, « cristallisation » puis décristallisation des pensions, devoir de mémoire au Cameroun.
\end{exoresolu}

\begin{attention}[Pièges fréquents]
Ne pas prendre l'affiche de propagande pour une description de la réalité : elle exprime une \cle{intention}. Ne pas opposer de façon manichéenne « tous volontaires » ou « tous forcés » : le corrigé montre une nuance. Dans les calculs, bien identifier la valeur initiale (1939) et la valeur finale (1943) avant d'appliquer la formule. Enfin, citer précisément les documents dans la rédaction : une affirmation sans référence n'est pas un argument historique.
\end{attention}

\courssub{Bilan de l'activité}

Cette enquête a permis de mettre en évidence trois acquis essentiels. D'abord, la \cle{distance critique} entre le discours officiel (ralliement spontané, Empire uni, volontariat) et la réalité vécue (contrainte, réquisitions, inégalités), que seul le croisement de sources hétérogènes permet d'établir. Ensuite, la double échelle d'analyse : au niveau \cle{macro}, le Cameroun est un enjeu stratégique et un fournisseur d'hommes et de matières premières ; au niveau \cle{micro}, la guerre est vécue comme une épreuve de pénuries, de séparations et d'injustices. Enfin, l'activité relie le passé au présent : la question des pensions et de la place des vétérans dans la mémoire nationale montre que la Deuxième Guerre mondiale n'est pas un chapitre clos, et qu'elle a préparé, par la prise de conscience des inégalités, les dynamiques de décolonisation étudiées dans la suite du programme. L'élève a ainsi pratiqué concrètement le métier d'historien : classer, critiquer, croiser, chiffrer, argumenter et débattre.

\newpage

\coursec{Méthodes et exercices résolus}

\courssub{Méthodes et exercices résolus}

\begin{methode}[Analyser un document historique (texte, image, carte) — Démarche en 4 temps]
\textbf{Objectif :} Répondre à une question de type « Présentez et commentez » ou « Que montre ce document ? » (Épreuve écrite Bac/Probatoire).
\begin{enumerate}
    \item \textbf{Identification (Nature, Auteur, Date, Contexte) :} Préciser la nature (discours, affiche, photo, article de presse, carte), l'auteur, la date exacte, le contexte historique immédiat (quoi ? qui ? quand ? où ? pourquoi ?).
    \item \textbf{Description / Analyse interne (Le « Quoi ») :} Décrire objectivement le contenu : idées forces (texte), composition/couleurs/symboles (image), légende/échelle/informations portées (carte). Utiliser le vocabulaire précis (\cle{propagande}, \cle{mobilisation}, \cle{ralliement}, \cle{traité}, \cle{front}).
    \item \textbf{Explication / Mise en perspective (Le « Pourquoi » et « Comment ») :} Mobiliser les connaissances du cours pour expliquer le sens, les intentions, la portée. Relier le document aux grandes lignes du chapitre (extension du conflit, violence paroxystique, rôle de l'Afrique/Cameroun, naissance ONU).
    \item \textbf{Conclusion critique (Portée et Limites) :} Résumer l'apport du document pour comprendre l'histoire. Signaler les limites (subjectivité, partialité, cadre restreint, censure).
\end{enumerate}
\textbf{Reflexe Bac :} Toujours citer le document (« Le document montre que... », « L'auteur affirme que... ») et dater les connaissances mobilisées.
\end{methode}

\begin{exoresolu}[Analyse d'une affiche de propagande : Le ralliement du Cameroun à la France Libre (1940)]
\textbf{Énoncé :} Vous présentez l'affiche ci-dessous (description textuelle fournie) diffusée au Cameroun en août 1940.
\begin{quote}
\textit{Description :} Affiche en couleurs. Au centre, un grand drapeau tricolore français flottant, surmonté de la Croix de Lorraine. Au premier plan, un soldat africain (tirailleur) en uniforme kaki, casque colonial, fusil à la main, regarde l'horizon avec détermination. En arrière-plan, silhouettes d'usines et de palmiers. En haut, en lettres rouges majuscules : « CAMEROUNAIS, LA FRANCE VOUS APPELLE ! ». En bas : « REJOIGNEZ LES FORCES FRANÇAISES LIBRES — GÉNÉRAL DE GAULLE ».
\end{quote}
\textbf{Question :} Présentez ce document et expliquez comment il illustre l'engagement du Cameroun dans la Deuxième Guerre mondiale aux côtés des Alliés (1940-1945).
\tcblower
\textbf{Solution détaillée}

\textbf{1. Identification}
\begin{itemize}
    \item \textbf{Nature :} Affiche de propagande politique et de recrutement militaire.
    \item \textbf{Auteur / Commanditaire :} Services d'information de la \cle{France Libre} (Comité national français / Gaulle), relayés par l'administration coloniale ralliée (Gouverneur Félix Éboué à Brazzaville, puis Pierre Cournarie au Cameroun).
    \item \textbf{Date :} Août 1940 (contexte immédiat : \cle{Appel du 18 Juin}, ralliement du Tchad le 26 août, ralliement du Cameroun le 27 août 1940).
    \item \textbf{Contexte :} Après l'armistice du 22 juin 1940 (Vichy), la France est divisée. L'AEF (Afrique-Équatoriale Française) bascule dans le camp allié grâce à Félix Éboué. Le Cameroun, sous mandat français, hésite puis rallie la France Libre le 27 août 1940 sous l'impulsion du gouverneur Cournarie et de Leclerc. C'est le point de départ de l'engagement camerounais.
\end{itemize}

\textbf{2. Description et Analyse interne}
\begin{itemize}
    \item \textbf{Composition visuelle :} Triangulaire (Drapeau/Croix de Lorraine au sommet $\rightarrow$ Soldat au centre $\rightarrow$ Arrière-plan économique). Couleurs symboliques : \textbf{Bleu-Blanc-Rouge} (République), \textbf{Rouge} pour l'appel urgent (texte), \textbf{Kaki/Vert} (armée, nature africaine).
    \item \textbf{Personnages/Symboles :} Le \textbf{tirailleur africain} (représentant le « \cle{Sénégalais} » générique, figure du colonisé loyal), la \textbf{Croix de Lorraine} (emblème de la France Libre, opposée à la Francisque de Vichy), le \textbf{Drapeau tricolore} (légitimité républicaine).
    \item \textbf{Textes :} Impératif (« REJOIGNEZ »), appel identitaire (« CAMEROUNAIS »), référence à l'autorité légitime (« GÉNÉRAL DE GAULLE »). L'arrière-plan (usines + palmiers) suggère la \cle{mobilisation économique} (matières premières : caoutchouc, bois, minerais) et humaine.
\end{itemize}

\textbf{3. Explication et Mise en perspective (Connaissances mobilisées)}
\begin{itemize}
    \item \textbf{L'engagement militaire :} L'affiche appelle au recrutement. \textbf{Fait :} Environ \num{18000} Camerounais s'engagent dans les \cle{Forces Françaises Libres (FFL)}. Ils forment le Bataillon de Marche n°1 (BM1), qui deviendra le IIIe bataillon de la \cle{13e DBLE}. Participation aux campagnes du \cle{Fezzan} (Libye, 1941), \cle{Bir Hakeim} (Libye, 1942), Italie (1944), Libération de la France (1944-45).
    \item \textbf{L'engagement économique :} L'arrière-plan industriel/agricole renvoie à l'effort de guerre. \textbf{Fait :} Le Cameroun fournit \textbf{caoutchouc} (indispensable pneus/avions), \textbf{bois}, \textbf{étain}, \textbf{aluminium}, \textbf{vivres}. L'économie est réorientée (contrôle des changes, réquisition, travail forcé/\cle{prestations} accrues).
    \item \textbf{Portée politique :} L'affiche gomme le statut de « mandat SDN » pour assimiler le Cameroun à la France (« LA FRANCE VOUS APPELLE »). Elle légitime l'autorité de Gaulle face à Vichy. C'est un outil de \cle{mobilisation idéologique} (défense de la « Liberté » vs « Oppression » nazie/vichyste).
\end{itemize}

\textbf{4. Conclusion critique}
\begin{itemize}
    \item \textbf{Apport :} Document clé pour comprendre la \cle{propagande} comme levier de ralliement et de recrutement. Montre la construction d'une image du « bon colonisé » combattant pour la « Mère Patrie ».
    \item \textbf{Limites :} Vision unilatérale (officielle). Ne montre pas : les réticences initiales, les tensions entre Vichy/Gaullistes au sein de l'administration, la contrainte (prestations obligatoires), le sort des vétérans oubliés après 1945, la censure de l'opposition nationaliste naissante (ex: USPC, UPC plus tard).
\end{itemize}
\textbf{Phrase de conclusion :} Cette affiche résume la double mobilisation (humaine et économique) du Cameroun, transformant un territoire de mandat en base arrière stratégique de la France Libre, tout en masquant les contradictions de l'ordre colonial.
\end{exoresolu}

\begin{methode}[Rédiger une réponse argumentée structurée (A.R.E.S.) — Le paragraphe type Bac]
\textbf{Objectif :} Répondre à une question explicative type « Expliquez pourquoi... », « Quelles sont les causes/conséquences de... », « Montrez que... » en 15-20 lignes.
\begin{enumerate}
    \item \textbf{Analyser le sujet :} Souligner le \cle{sujet} (quoi ?), l'\cle{aspect} (causes ? conséquences ? rôle ?), la \cle{chronologie} (dates bornes), l'\cle{espace} (monde, Afrique, Cameroun). Définir les termes clés.
    \item \textbf{Construire le plan (2 ou 3 arguments majeurs) :} Hiérarchiser : Cause principale / Cause secondaire ; Conséquence immédiate / Conséquence structurelle. Un argument = une idée force.
    \item \textbf{Rédiger le paragraphe (Structure A.R.E.S.) :}
    \begin{itemize}
        \item \textbf{A}ccroche / \textbf{A}nnonce du plan : Phrase d'introduction situant le contexte + thèse directe (réponse à la question) + annonce des arguments (« D'abord... Ensuite... Enfin... »).
        \item \textbf{R}easoning (Développement) : \textbf{Un paragraphe par argument}. Structure : \textbf{Idée force} (phrase thèse) $\rightarrow$ \textbf{Preuves/Faits précis} (dates, chiffres, noms, événements, concepts du cours) $\rightarrow$ \textbf{Explication/Lien} (pourquoi ce fait prouve l'idée force).
        \item \textbf{S}ynthèse / \textbf{S}ortie : Phrase de conclusion qui résume la réponse et ouvre (lien avec la suite du programme : Guerre froide, Décolonisation, ONU).
    \end{itemize}
    \item \textbf{Relire :} Vérifier : Réponse directe ? Dates exactes ? Vocabulaire spécifique ? Connecteurs logiques (\textit{donc, par conséquent, en outre, cependant}) ? Pas de « je » ?
\end{enumerate}
\textbf{Formule magique :} \textbf{1 Idée Force = 1 Paragraphe = 1 Titre (mental) + 2 à 3 Faits Précis + 1 Lien explicatif.}
\end{methode}

\begin{exoresolu}[Réponse argumentée : L'extension mondiale du conflit (1941)]
\textbf{Énoncé :} « Expliquez comment et pourquoi la guerre européenne devient un conflit mondial en 1941. » (Attendu : ~15-20 lignes).
\tcblower
\textbf{Solution détaillée}

\textbf{Analyse du sujet :}
\begin{itemize}
    \item \textbf{Sujet :} Extension géographique de la WWII.
    \item \textbf{Aspect :} Mécanismes (\textit{comment}) et Causes (\textit{pourquoi}).
    \item \textbf{Chronologie :} Année charnière \textbf{1941} (avant = guerre européenne/africaine ; après = guerre mondiale).
    \item \textbf{Espace :} Monde entier.
    \item \textbf{Mots-clés :} Opération Barbarossa, Pearl Harbor, Axe (Allemagne, Italie, Japon), Alliés (UK, URSS, USA), Idéologie, Ressources, Alliances.
\end{itemize}

\textbf{Plan (2 arguments majeurs + 1 structurel) :}
\begin{enumerate}
    \item L'élargissement idéologique et stratégique par l'attaque de l'URSS (Juin 1941).
    \item La mondialisation effective par l'entrée en guerre des USA (Décembre 1941).
    \item (Transversal) L'engrenage des alliances et la guerre des ressources.
\end{enumerate}

\textbf{Rédaction (A.R.E.S.) :}

\textit{Accroche/Annonce :} En 1941, le conflit bascule d'une guerre hégémonique européenne à une \cle{guerre mondiale} totale par la conjonction de deux événements majeurs : l'invasion de l'URSS par le Reich et l'attaque japonaise sur Pearl Harbor, transformant les théâtres d'opérations distincts en un front unique planétaire.

\textit{Argument 1 : Le tournant continental — L'Opération Barbarossa (22 juin 1941).}
\textbf{Idée force :} Hitler ouvre un second front à l'Est par obsession idéologique (\cle{lebensraum}, anticommunisme) et nécessité économique (pétrole Caucase, blé Ukraine), brisant le pacte germano-soviétique.
\textbf{Preuves :} 3,8 millions de soldats (Wehrmacht + alliés roumains, finlandais, hongrois), front de 2000 km. Bataille de \cle{Moscou} (déc. 1941) : premier échec majeur allemand. Entrée en guerre de l'URSS aux côtés du Royaume-Uni (Accords de Moscou, juillet 1941).
\textbf{Lien :} Ce n'est plus une guerre de frontières mais une \cle{guerre d'anéantissement} (commissaires politiques, Einsatzgruppen, famine programmée), qui lie le sort de l'Europe à l'immense espace eurasien.

\textit{Argument 2 : Le tournant océanique — Pearl Harbor (7 décembre 1941).}
\textbf{Idée force :} Le Japon, en quête d'autarcie (\cle{Sphère de coprospérité}), attaque la flotte US pour lever l'embargo pétrolier (juillet 1941), entraînant l'entrée en guerre de la première puissance industrielle.
\textbf{Preuves :} Attaque surprise sur \cle{Pearl Harbor} (Hawaï). Déclaration de guerre US au Japon (8 déc.), puis Allemagne/Italie aux USA (11 déc. - erreur stratégique de Hitler). Création de l'axe \cle{Rome-Berlin-Tokyo} (Pacte tripartite 1940) devenu opérationnel.
\textbf{Lien :} Le conflit relie désormais l'Europe, l'URSS, l'Asie-Pacifique et l'Amérique du Nord. Les \cle{Alliés} (Churchill, Roosevelt, Staline) forment la « \cle{Grande Alliance} » (Conférence de Washington, déc. 1941 : principe « Germany first »).

\textit{Synthèse/Ouverture :} 1941 scelle la \cle{mondialisation} du conflit : géographique (5 continents), humaine (mobilisation coloniale massive, dont l'Afrique), industrielle (économie de guerre US/URSS) et idéologique (fascisme vs démocratie vs communisme). Cette extension rend inévitable la \cle{conférence de l'ONU} (1942) et pose les bases de l'ordre de 1945.
\end{exoresolu}

\begin{methode}[Construire et exploiter une frise chronologique commentée / Identifier les phases]
\textbf{Objectif :} Situer des événements dans le temps, mettre en évidence des ruptures, des simultanéités, des durées (compétence transversale, souvent évaluée en TD ou partie « Repères » du Bac).
\begin{enumerate}
    \item \textbf{Délimiter l'échelle :} Choisir l'amplitude (ex: 1939-1945) et l'unité (mois ou trimestres). Tracer l'axe horizontal (flèche vers la droite).
    \item \textbf{Sélectionner les repères obligatoires (Programme) :} \cle{1er sept. 1939} (invasion Pologne), \cle{10 mai 1940} (offensive Ouest), \cle{22 juin 1940} (armistice), \cle{27 août 1940} (ralliement Cameroun), \cle{22 juin 1941} (Barbarossa), \cle{7 déc. 1941} (Pearl Harbor), \cle{janv. 1942} (Conférence Wannsee / Déclaration ONU), \cle{6 juin 1944} (Débarquement Normandie), \cle{8 mai 1945} (Capitulation Allemagne), \cle{6/9 août 1945} (Hiroshima/Nagasaki), \cle{2 sept. 1945} (Capitulation Japon), \cle{24 oct. 1945} (Entrée en vigueur Charte ONU).
    \item \textbf{Colorier / Grouper en PHASES (niveau supérieur) :} Utiliser des codes couleur ou des accolades au-dessus de l'axe pour nommer les 3-4 grandes phases du cours :
    \begin{itemize}
        \item \textbf{Phase 1 : L'éclatement et l'hégémonie de l'Axe (sept. 1939 - juin 1941).} Guerre éclair, chute de la France, Bataille d'Angleterre.
        \item \textbf{Phase 2 : La mondialisation et le paroxysme (juin 1941 - fin 1942/43).} Barbarossa, Pearl Harbor, Shoah (Wannsee), tournant Stalingrad/El Alamein/Midway.
        \item \textbf{Phase 3 : La reconquête alliée et l'effondrement de l'Axe (1943 - sept. 1945).} Débarquements, libération camps, conférences (Téhéran, Yalta, Potsdam), bombe atomique.
        \item \textbf{Phase 4 : L'après-guerre immédiat (1945-...).} Naissance ONU, procès Nuremberg, début décolonisation/Guerre froide.
    \end{itemize}
    \item \textbf{Annoter les simultanéités (verticalité) :} Placer au-dessus de l'axe (Europe/Monde) et en-dessous (Afrique/Cameroun) pour voir les interactions (ex: 1940 : Armistice ET Ralliement Cameroun ; 1944 : Débarquement ET Conférence de Brazzaville).
\end{enumerate}
\textbf{Reflexe Bac :} On ne note JAMAIS une date sans l'événement, ni un événement sans sa date. Une frise sans phases nommées est une liste, pas un outil d'analyse.
\end{methode}

\begin{exoresolu}[Frise commentée : Le Cameroun dans la WWII (1939-1945)]
\textbf{Énoncé :} Construisez une frise chronologique commentée (axe + 10 repères datés minimum + 3 phases nommées + annotations « Afrique/Cameroun ») pour illustrer la participation du Cameroun à la Deuxième Guerre mondiale.
\tcblower
\textbf{Solution détaillée}

\textbf{Code TikZ de la frise (à insérer dans le cahier / copie) :}
\begin{popfigure}
\begin{tikzpicture}[scale=0.85, every node/.style={font=\footnotesize}]
    % Axe principal
    \draw[thick, ->] (0,0) -- (16,0) node[right] {Temps};
    % Graduations annuelles
    \foreach \x/\year in {0/1939, 1.6/1940, 3.2/1941, 4.8/1942, 6.4/1943, 8/1944, 9.6/1945, 11.2/1946} {
        \draw (\x, -0.1) -- (\x, 0.1) node[below=2pt] {\year};
    }
    % PHASES (Accolades manuelles au-dessus)
    \draw[popblue, thick] (0, 1.2) -- (0, 1.4) -- (3.2, 1.4) -- (3.2, 1.2);
    \node[popblue, above=8pt] at (1.6, 1.4) {\textbf{Phase 1 : Ralliement \& Base Arrière (1939-1941)}};
    \draw[poporange, thick] (3.2, 1.2) -- (3.2, 1.4) -- (6.4, 1.4) -- (6.4, 1.2);
    \node[poporange, above=8pt] at (4.8, 1.4) {\textbf{Phase 2 : Mobilisation Totale (1941-1943)}};
    \draw[popgreen, thick] (6.4, 1.2) -- (6.4, 1.4) -- (9.6, 1.4) -- (9.6, 1.2);
    \node[popgreen, above=8pt] at (8.0, 1.4) {\textbf{Phase 3 : Libération \& Mutations (1944-1945)}};
    % EVENEMENTS MONDE / EUROPE (Au-dessus phases)
    \node[align=center, anchor=south] at (0.2, 2.5) {\textbf{1 sept. 39}\\\textit{Invasion Pologne}};
    \node[align=center, anchor=south] at (1.6, 3.0) {\textbf{22 juin 40}\\\textit{Armistice France}};
    \node[align=center, anchor=south] at (3.2, 2.5) {\textbf{22 juin 41}\\\textit{Barbarossa}};
    \node[align=center, anchor=south] at (4.0, 3.0) {\textbf{7 déc. 41}\\\textit{Pearl Harbor}};
    \node[align=center, anchor=south] at (5.0, 2.5) {\textbf{Janv. 42}\\\textit{Wannsee / Décl. ONU}};
    \node[align=center, anchor=south] at (8.0, 3.0) {\textbf{6 juin 44}\\\textit{Débarquement}};
    \node[align=center, anchor=south] at (9.2, 2.5) {\textbf{8 mai / 2 sept. 45}\\\textit{Capitulations}};
    \node[align=center, anchor=south] at (10.5, 3.0) {\textbf{24 oct. 45}\\\textit{ONU effective}};
    % EVENEMENTS CAMEROUN / AFRIQUE (En-dessous de l'axe)
    \node[align=center, anchor=north, poppurple] at (1.7, -1.5) {\textbf{27 août 40}\\\textit{Ralliement Cameroun (Cournarie/Leclerc)}};
    \node[align=center, anchor=north, poppurple] at (2.5, -2.2) {\textbf{Nov. 40}\\\textit{BM1 créé (Yaoundé)}};
    \node[align=center, anchor=north, poppurple] at (3.5, -1.5) {\textbf{1941-42}\\\textit{Campagnes Fezzan / Bir Hakeim (BM1)}};
    \node[align=center, anchor=north, poppurple] at (4.5, -2.2) {\textbf{1942-43}\\\textit{Effort max : Caoutchouc, Étain, Prestations}};
    \node[align=center, anchor=north, poppurple] at (6.5, -1.5) {\textbf{Janv. 44}\\\textit{Conf. Brazzaville (Réforme coloniale)}};
    \node[align=center, anchor=north, poppurple] at (8.2, -2.2) {\textbf{1944-45}\\\textit{BM1 en Italie / France (Libération)}};
    \node[align=center, anchor=north, poppurple] at (9.6, -1.5) {\textbf{1945}\\\textit{Retour vétérans / Naissance syndicalisme (USC)}};
    % Flèches de liaison
    \draw[->, popmuted, dashed] (1.6, 0.1) -- (1.7, -1.4);
    \draw[->, popmuted, dashed] (8.0, 0.1) -- (8.2, -2.1);
\end{tikzpicture}
\legende{Frise chronologique commentée : Le Cameroun dans la WWII. Au-dessus : repères mondiaux. En-dessous (violet) : repères camerounais. Phases colorées. Flèches pointillées = interactions.}
\end{popfigure}

\textbf{Commentaire écrit attendu (exemple pour l'oral ou la copie) :}
« Cette frise met en évidence la \textbf{précocité} de l'engagement camerounais (août 1940, phase 1) par rapport à l'entrée en guerre de l'URSS et des USA (1941). Elle montre la \textbf{double chronologie} : le pic de l'effort économique camerounais (1942-43) coïncide avec le tournant militaire allié (Stalingrad, El Alamein). Enfin, elle révèle la \textbf{porosité} entre guerre et décolonisation : la Conférence de Brazzaville (janv. 1944) prépare l'après-guerre avant même la fin des combats en Europe. »
\end{exoresolu}

\begin{methode}[Expliquer par la causalité multi-scalaire (Causes / Facteurs / Conséquences) — Tableau d'analyse]
\textbf{Objectif :} Répondre à « Quelles sont les causes de... ? » ou « Analysez les conséquences de... » en évitant la liste de courses. Distinguer niveaux (structurels, conjoncturels, déclencheurs) et échelles (international, national, local).
\begin{enumerate}
    \item \textbf{Identifier le phénomène central :} Ex: « L'entrée en guerre du Cameroun », « La naissance de l'ONU », « Le génocide des Juifs ».
    \item \textbf{Construire un tableau 3 colonnes (ou 3 paragraphes distincts) :}
    \begin{center}
    \begin{tabularx}{\linewidth}{|X|X|X|}\hline
    \rowcolor{popblueL}
    \textbf{Causes Structurelles (Long terme)} & \textbf{Causes Conjoncturelles (Moyen terme)} & \textbf{Déclencheurs / Facteurs immédiats (Court terme)} \\ \hline
    Ex: Traité de Versailles, Crise 1929, Idéologies totalitaires, Mandat SDN Cameroun. & Ex: Réarmement, Anschluss, Munich, Politique d'apaisement, Crise du ralliement 1940. & Ex: Invasion Pologne (1er sept 39), Appel 18 Juin, Ralliement Éboué/Cournarie (août 40). \\ \hline
    \end{tabularx}
    \end{center}
    \item \textbf{Hiérarchiser et Lier :} Expliquer \textit{comment} le structurel pèse sur le conjoncturel, et \textit{comment} le déclencheur fait basculer la situation. Utiliser : \textit{« Crée les conditions de... », « Précipite... », « Détermine le choix de... »}.
    \item \textbf{Traiter les Conséquences (Même logique) :} Immédiates (mortalité, destruction), Structurelles (ONU, Guerre froide, Décolonisation, Droits de l'Homme), Locales (Cameroun : vétérans, syndicalisme, nationalisme, réforme Brazzaville).
\end{enumerate}
\textbf{Attention :} Ne pas confondre \textbf{Cause} (origine) et \textbf{Prétexte} (justification officielle). Ne pas confondre \textbf{Conséquence} (résultat direct/indirect) et \textbf{Événement postérieur} (simple chronologie).
\end{methode}

\begin{exoresolu}[Analyse causale : Les conséquences de la WWII sur le Cameroun]
\textbf{Énoncé :} « Montrez que la Deuxième Guerre mondiale est un accélérateur de l'histoire politique et sociale au Cameroun. » (Utilisez un tableau causes/conséquences ou paragraphes structurés).
\tcblower
\textbf{Solution détaillée}

\textbf{Tableau d'analyse multi-scalaire : Conséquences de la WWII sur le Cameroun}

\begin{center}
\begin{tabularx}{\linewidth}{|>{\bfseries}X|X|X|}\hline
\rowcolor{popblueL}
\textbf{Échelle / Domaine} & \textbf{Conséquences Immédiates (1945-1947)} & \textbf{Conséquences Structurelles (Long terme : Décolonisation)} \\ \hline
\textbf{Politique / Institutionnel} & 
- Fin du régime de l'indigénat (décret 1946, suite Brazzaville).
- Création d'assemblées locales (ARTCAM 1946) et représentation à l'Assemblée française (2 députés : Louis-Paul Aujoulat, Alexandre Douala Manga Bell).
- Conférence de Brazzaville (janv. 1944) : refus de l'indépendance mais promesse d'évolution. & 
- Échec de l'assimilation $\rightarrow$ Montée du nationalisme (\cle{UPC} créée 1948 par Ruben Um Nyobè).
- Revendication de l'indépendance immédiate (pas d'autonomie interne).
- Guerre de décolonisation (1955-1971) : racine dans la frustration post-1945 (promesses non tenues). \\ \hline
\textbf{Social / Syndical} & 
- Retour de \num{~18000} vétérans (BM1, travailleurs forcés) : nouvelles compétences, conscience politique, refus de l'ordre colonial.
- Création de l'\cle{Union des Syndicats Confédérés du Cameroun (USCC)} (1945), puis \cle{USC} (1947) : premières grèves (Douala 1945, 1947). & 
- Émergence d'une \cle{classe ouvrière} et d'une \cle{élite intellectuelle} politisée (syndicalistes $\rightarrow$ nationalistes).
- Lutte pour les droits sociaux (code du travail 1952) liée aux droits de l'homme (Décl. 1948). \\ \hline
\textbf{Économique} & 
- Fin de l'économie de guerre (contrôles, réquisitions).
- Pénuries, inflation, développement du marché noir.
- Infrastructure : route Douala-Yaoundé, port de Douala, chemin de fer (Transcam) développés pour l'effort allié. & 
- Héritage d'une économie d'extraction (plantation, mine) peu diversifiée.
- Dépendance structurelle vis-à-vis de la France (Zone Franc, accords de coopération 1960). \\ \hline
\textbf{International / Symbolique} & 
- Cameroun membre fondateur de l'\cle{ONU} (1945) via la France.
- Passage du \cle{Mandat SDN} au \cle{Tutelle ONU} (1946) : obligation de rapport annuel, pétitions recevables (Um Nyobè à l'ONU 1952-58). & 
- Internationalisation de la question camerounaise (Tribune ONU).
- Modèle \cle{indépendance par référendum} (1960/1961) sous supervision ONU. \\ \hline
\end{tabularx}
\end{center}

\textbf{Synthèse argumentée (Réponse type Bac) :}
La WWII agit comme un \textbf{révélateur et un accélérateur}.
\begin{enumerate}
    \item \textbf{Révélateur :} La contribution massive (sang, caoutchouc, étain) prouve la capacité d'auto-administration et contredit l'idéologie coloniale de l'« inaptitude » des Africains. La Conférence de Brazzaville (1944) acte la fin de l'indigénat mais refuse l'indépendance, créant un fossé infranchissable avec les nouvelles élites (vétérans, syndicalistes, UPC).
    \item \textbf{Accélérateur :} Le statut de \cle{Territoire sous Tutelle ONU} (1946) offre un cadre juridique international (pétitions, missions de visite) absent sous le Mandat. Les vétérans (ex: \cle{André-Marie Mbida}, \cle{Ernest Ouandié}) deviennent les cadres de la décolonisation. La guerre froide (contexte ONU) transforme la revendication nationaliste en enjeu géopolitique.
\end{enumerate}
\textbf{Conclusion :} La WWII ne \textit{crée} pas le nationalisme camerounais, mais elle en fournit les \textbf{acteurs} (vétérans), les \textbf{outils} (syndicats, ONU, Tribune), et la \textbf{légitimité} (droits de l'homme, contribution à la victoire), rendant la décolonisation inéluctable.
\end{exoresolu}

\begin{methode}[Réaliser / Analyser un croquis de carte (Schéma géographique) — Méthode en 5 étapes]
\textbf{Objectif :} « Réalisez un croquis pour montrer... » ou « Analysez le croquis ci-dessous ». Compétence géohistorique majeure au Bac Tech.
\begin{enumerate}
    \item \textbf{Analyser le sujet :} Identifier le \textbf{Thème} (ex: Contribution du Cameroun), l'\textbf{Échelle} (Cameroun, Afrique, Monde), la \textbf{Période} (1940-1945). Définir le \textbf{Titre} précis (Thème + Espace + Dates).
    \item \textbf{Sélectionner l'information (Légende) :} Choisir 4 à 6 items max, regroupés en \textbf{2 ou 3 catégories logiques} (ex: 1. Engagement militaire, 2. Effort économique, 3. Enjeux stratégiques). Hiérarchiser : Figures de surface (zones), Figures linéaires (axes/flux), Figures ponctuelles (villes/sites).
    \item \textbf{Choisir les signes graphiques (Sémiologie) :}
    \begin{itemize}
        \item \textbf{Surfaces (Zones) :} Hachures, pointillés, couleurs claires (ex: Zone de recrutement, Zone de production caoutchouc).
        \item \textbf{Lignes (Flux/Axes) :} Flèches \textbf{épaisseur proportionnelle} (quantité), couleur distincte (ex: Flèche hommes $\rightarrow$ Nord/Est ; Flèche matières premières $\rightarrow$ Douala $\rightarrow$ Europe).
        \item \textbf{Points (Villes/Sites) :} Carrés (capitales), Triangles (mines), Cercles (ports), Étoiles (batailles). Taille = importance.
    \end{itemize}
    \item \textbf{Construire le fond de carte (TikZ/Esquisse) :} Contour simplifié du Cameroun (polygone gros traits). Repères indispensables : \cle{Douala} (port), \cle{Yaoundé} (capitale), \cle{Ngaoundéré} (Nord), \cle{Bertoua} (Est), \cle{Kribi} (port minier), \cle{Édéa} (aluminium/énergie). Fleuve \cle{Sanaga}, \cle{Bénoué}, \cle{Logone}. Frontières (Nigéria, Tchad, RCA, Congo, Guinée Eq.).
    \item \textbf{Finaliser :} \textbf{Légende organisée} (Titre, Items classés par catégorie, Signes exacts), \textbf{Flèche Nord}, \textbf{Échelle approximative}, \textbf{Titre global} en haut.
\end{enumerate}
\textbf{Reflexe Bac :} Un croquis sans légende organisée = 0 point. Un croquis sans titre = 0 point. Les flèches de flux DOIVENT avoir une épaisseur variable si les quantités diffèrent.
\end{methode}

\begin{exoresolu}[Croquis : La contribution du Cameroun à l'effort de guerre allié (1940-1945)]
\textbf{Énoncé :} Réalisez un croquis légendé intitulé : « Le Cameroun, base arrière et pourvoyeur de l'effort de guerre allié (1940-1945) ».
\tcblower
\textbf{Solution détaillée : Description du croquis attendu (Code TikZ fourni pour l'enseignant / élève avancé)}

\textbf{1. Titre :} \textbf{LE CAMEROUN, BASE ARRIÈRE ET POURVOYEUR DE L'EFFORT DE GUERRE ALLIE (1940-1945)}

\textbf{2. Légende organisée (À tracer à droite de la carte) :}
\begin{center}
\begin{tabular}{|l|l|l|}\hline
\rowcolor{popblueL} \textbf{Catégorie} & \textbf{Information} & \textbf{Signe graphique} \\ \hline
\multicolumn{3}{|c|}{\textbf{I. ENGAGEMENT MILITAIRE}} \\ \hline
Zone & Zone de recrutement principal (Sud/Centre) & \textcolor{popblue}{\rule{1cm}{4pt}} Hachures bleu clair \\ \hline
Point & Yaoundé : Centre d'instruction / BM1 & $\blacksquare$ Carré bleu foncé \\ \hline
Flux & Départ volontaires $\rightarrow$ Tchad (Leclerc) / Égypte & $\xrightarrow{\text{\textbf{---}}} $ Flèche bleue fine pointillée \\ \hline
Point & Batailles extérieures : Fezzan, Bir Hakeim, Italie & $\star$ Étoile rouge (hors carte, encart) \\ \hline
\multicolumn{3}{|c|}{\textbf{II. EFFORT ÉCONOMIQUE (MATIÈRES PREMIÈRES)}} \\ \hline
Zone & Zone caoutchouc / bois (Forêt équatoriale Sud/Est) & \textcolor{popgreen}{\rule{1cm}{4pt}} Pointillés verts \\ \hline
Zone & Zone minière Étain (Bertoua) / Aluminium (Édéa) & \textcolor{poporange}{\rule{1cm}{4pt}} Hachures orange \\ \hline
Point & Mines de \cle{Bertoua} (Étain) & $\blacktriangle$ Triangle orange \\ \hline
Point & Usine \cle{Édéa} (Aluminium - Alucam) & $\blacklozenge$ Losange orange \\ \hline
Flux & Évacuation productions $\rightarrow$ \cle{Douala} (Port) & $\xrightarrow{\text{\textbf{====}}} $ Flèche verte \textbf{ÉPAISSE} \\ \hline
Flux & Export Douala $\rightarrow$ Alliés (Libreville/UK/US) & $\xrightarrow{\text{\textbf{====}}} $ Flèche verte \textbf{ÉPAISSE} (mer) \\ \hline
\multicolumn{3}{|c|}{\textbf{III. ENJEUX STRATÉGIQUES}} \\ \hline
Point & \cle{Douala} : Port principal, Base navale/alliée & $\bigcirc$ Grand cercle bleu + ancre \\ \hline
Ligne & Chemin de fer (Transcam) Douala - Yaoundé - Ngaoundéré & \texttt{---$\bullet$---} Trait pointillé noir \\ \hline
Ligne & Route stratégique Douala - Yaoundé - Bertoua & \texttt{---} Trait continu gris \\ \hline
\end{tabular}
\end{center}

\textbf{3. Code TikZ du fond de carte et du croquis (Modèle simplifié) :}
\begin{popfigure}
\begin{tikzpicture}[scale=1.1, every node/.style={font=\scriptsize}]
    % CONTOUR CAMEROUN (Polygone très simplifié)
    \draw[thick, popdark] (0,0) -- (1, -0.5) -- (2.5, -1) -- (4, -0.5) -- (5, 0.5) -- (5.5, 2) -- (5, 3.5) -- (4, 4.5) -- (2.5, 5) -- (1, 4.5) -- (0, 3.5) -- (-0.5, 2) -- cycle;
    % VILLES / POINTS CLÉS
    \node[popblue, fill=popblue, circle, inner sep=1.5pt, label=left:\textbf{Douala}] (Douala) at (0.5, 1.2) {};
    \node[popblue, fill=popblue, rectangle, inner sep=2pt, label=right:\textbf{Yaoundé}] (Yaounde) at (2.5, 2.5) {};
    \node[poporange, fill=poporange, regular polygon, regular polygon sides=3, inner sep=2pt, label=right:\textbf{Bertoua (Étain)}] (Bertoua) at (4.2, 3.0) {};
    \node[poporange, fill=poporange, diamond, inner sep=2pt, label=right:\textbf{Édéa (Alu)}] (Edea) at (1.5, 1.8) {};
    \node[popdark, fill=popdark, circle, inner sep=1.5pt, label=above:\textbf{Ngaoundéré}] (Ngaoundere) at (2.5, 4.2) {};
    \node[popdark, fill=popdark, circle, inner sep=1.5pt, label=right:\textbf{Kribi}] (Kribi) at (0.8, 0.2) {};
    % FLEUVES
    \draw[popblue!50, thick] (2.5, 2.5) .. controls (1.5, 1.5) .. (0.5, 1.2) node[pos=0.5, left] {Sanaga};
    \draw[popblue!50, thick] (2.5, 4.2) -- (3.5, 2.5) node[pos=0.5, right] {Bénoué};
    % CHEMIN DE FER (Transcam)
    \draw[densely dashed, thick, popdark] (Douala) -- (Yaounde) -- (Ngaoundere);
    % ROUTE
    \draw[densely dotted, thick, popmuted] (Douala) -- (Yaounde) -- (Bertoua);
    % ZONES (Surfaces - Pattern) - Utilisation de \fill au lieu de \pattern
    % Zone Caoutchouc/Bois (Sud-Est)
    \fill[pattern=north east lines, pattern color=popgreen!50] (2.5, 2.5) .. controls (3.5, 3.5) and (4.5, 3) .. (5, 2) -- (5.5, 2) -- (5, 3.5) -- (4, 4.5) -- (2.5, 5) -- cycle;
    \node[popgreen, rotate=30] at (4, 3.5) {\textbf{Caoutchouc / Bois}};
    % Zone Minière (Est)
    \fill[pattern=north west lines, pattern color=poporange!50] (3.5, 2.5) -- (4.5, 3) -- (4.2, 3.5) -- (3.2, 3) -- cycle;
    \node[poporange] at (3.8, 2.8) {\textbf{Mines}};
    % FLUX EXPORT (Douala -> Mer)
    \draw[popgreen, very thick, ->, >=Stealth] (Douala) -- ++(-1.5, 0) node[left, align=center] {Export\\Alliés};
    % FLUX RECRUTEMENT (Yaoundé -> Nord/Tchad)
    \draw[popblue, dashed, ->, >=Stealth] (Yaounde) -- ++(0, 1.5) node[above] {Vers Tchad/Leclerc};
    % NORD
    \node[anchor=south] at (2.5, 5.5) {\textbf{N}};
    \draw[->, thick] (2.5, 5.3) -- (2.5, 5.7);
    % ÉCHELLE
    \node[anchor=north] at (5.5, -1.2) {Échelle approx. 1 cm $\approx$ 150 km};
\end{tikzpicture}
\legende{Croquis : Le Cameroun dans l'effort de guerre (1940-1945). Légende organisée en 3 catégories. Signes respectant la sémiologie (surfaces=zones, flèches épaisseur=quantité flux, points=localisation).}
\end{popfigure}

\textbf{Commentaire type (Analyse du croquis) :} « Ce croquis met en évidence la \textbf{double contribution} : humaine (flèche pointillée vers le Nord/Tchad, base du BM1) et économique (flèches \textbf{très épaisses} vertes convergent vers Douala, port unique d'exportation). Il montre la \textbf{dépendance infrastructurelle} (rail/route unique Nord-Sud) et la \textbf{spatialisation} des ressources : forêt au Sud (caoutchouc), mines à l'Est (étain), énergie/industrie à l'Ouest (Édéa). Douala apparaît comme le \textbf{nœud stratégique} incontournable (port, rail, route, aluminium). »
\end{exoresolu}

\begin{methode}[Traiter une question de synthèse / Dissertation courte (Plan détaillé) — Structure en 3 parties]
\textbf{Objectif :} Répondre à un sujet de type « Dans quelle mesure... », « Montrez que... », « Étudier... » en construisant un plan logique (Thèse/Antithèse/Synthèse ou Causes/Manifestations/Consequences) avec introduction, développement en 2 ou 3 parties équilibrées, conclusion.
\begin{enumerate}
    \item \textbf{Analyse du sujet (Brouillon) :} Définir termes, bornes spatio-temporelles, problème central (problématique). Ex: « L'Afrique dans la WWII » $\rightarrow$ Problématique : « L'Afrique a-t-elle été seulement un enjeu ou un acteur décisif de la victoire, et quelles en ont été les conséquences politiques ? »
    \item \textbf{Plan détaillé (Niveau Bac Tech : 2 parties principales + sous-parties) :}
    \begin{itemize}
        \item \textbf{Introduction :} Accroche (chiffre/citation) $\rightarrow$ Définitions $\rightarrow$ Contexte $\rightarrow$ \textbf{Problématique} $\rightarrow$ \textbf{Annonce du plan}.
        \item \textbf{Partie I : L'Afrique, enjeu et théâtre d'opérations (1939-1942).} (A. Convoitée par l'Axe / B. Base de la France Libre / C. Théâtre secondaire mais stratégique).
        \item \textbf{Partie II : L'Afrique, actrice de la victoire et creuset des mutations (1942-1945+).} (A. Mobilisation humaine/économique massive / B. Brazzaville 1944 : rupture réformiste / C. Vétérans et naissance du nationalisme).
        \item \textbf{Conclusion :} Bilan synthétique (Réponse à la problématique) $\rightarrow$ Ouverture (Décolonisation, Guerre froide, Mémoire).
    \end{itemize}
    \item \textbf{Rédiger l'introduction et la conclusion au propre (obligatoire au Bac).} Le développement peut rester en plan détaillé (arguments + exemples précis) si temps court, ou être rédigé.
    \item \textbf{Connecteurs :} \textit{Premièrement, En outre, Cependant, Par conséquent, En définitive.}
\end{enumerate}
\textbf{Critères de réussite :} Problématique explicite, Plan équilibré (pas 3 lignes / 1 page), Exemples \textbf{précis et datés} (noms, chiffres, lieux), Vocabulaire du cours (\cle{indigénat}, \cle{tutelle}, \cle{prestations}, \cle{BM1}, \cle{Brazzaville}).
\end{methode}

\begin{exoresolu}[Plan détaillé de dissertation : L'Afrique dans la Deuxième Guerre mondiale]
\textbf{Énoncé :} Sujet : « L'Afrique dans la Deuxième Guerre mondiale : enjeux, contributions, mutations. » Faites une introduction rédigée, un plan détaillé équilibré (2 parties, 2 sous-parties chacune) et une conclusion rédigée.
\tcblower
\textbf{Solution détaillée}

\textbf{INTRODUCTION RÉDIGÉE}
\textit{[Accroche]} « L'Afrique a été le premier continent où la France libre a pu se reconstruire », déclare le Général de Gaulle en 1944. \textit{[Contexte]} Pourtant, en 1939, l'Afrique coloniale est un empire endormi, régi par le Code de l'indigénat, intégrée à l'économie métropolitaine. \textit{[Déclencheur]} La défaite de juin 1940 et le ralliement de l'AEF (Août 1940) puis de l'AOF (Nov 1942) transforment l'Afrique en \textbf{base arrière vitale} pour les Alliés. \textit{[Problématique]} \textbf{Dans quelle mesure la participation africaine à la victoire alliée (1940-1945) a-t-elle constitué le catalyseur de la décolonisation, transformant des « sujets » en citoyens revendicateurs ?} \textit{[Annonce]} Nous montrerons d'abord que l'Afrique est un \textbf{enjeu stratégique et un théâtre d'opérations} (I), avant d'analyser comment sa \textbf{contribution massive engendre des mutations politiques et sociales irréversibles} (II).

\textbf{PLAN DÉTAILLÉ}

\textbf{I. L'AFRIQUE : ENJEU STRATÉGIQUE ET THÉÂTRE D'OPÉRATIONS (1939-1942)}
\begin{itemize}
    \item \textbf{A. Un enjeu convoité par l'Axe et les Alliés.}
    \begin{itemize}
        \item \textbf{Fait :} Projet allemand \cle{Mittelafrika} (récupération colonies perdues 1919). Italie (Mussolini) vise l'Afrique de l'Est (Éthiopie, Somalie).
        \item \textbf{Fait :} Contrôle des ressources (Or, Uranium \cle{Shinkolobwe} Congo belge $\rightarrow$ Bombe A, Caoutchouc, Étain, Manganèse Gabon/Cameroun).
        \item \textbf{Fait :} Position géostratégique : Route du Cap (contournement Méditerranée), bases aériennes (Dakar, Brazzaville, Douala) pour transit avions US (Ferrying Command).
    \end{itemize}
    \item \textbf{B. Le ralliement : La France renaît d'Afrique (1940-1942).}
    \begin{itemize}
        \item \textbf{Fait :} \cle{18 Juin 1940} : Appel lancé « depuis le sol africain » (BBC). \cle{26 Août 1940} : Ralliement Tchad (Éboué). \cle{27 Août 1940} : Cameroun (Cournarie/Leclerc). \cle{Nov 1942} : AOF ralliée après débarquement allié Afrique du Nord (Opération Torch).
        \item \textbf{Fait :} \cle{Colonne Leclerc} : Du Tchad au Fezzan (Libye) $\rightarrow$ Koufra (mars 1941) $\rightarrow$ Tripoli (janv 1943). Première victoire française.
        \item \textbf{Analyse :} L'Afrique fournit \textbf{légitimité, territoire, hommes, ressources} à la France Libre. Sans l'Afrique, pas de France combattante, pas de siège à l'ONU 1945.
    \end{itemize}
\end{itemize}

\textbf{II. L'AFRIQUE : ACTRICE DE LA VICTOIRE ET CREUSET DES MUTATIONS (1942-1945+)}
\begin{itemize}
    \item \textbf{A. Une mobilisation totale : Hommes et Économie.}
    \begin{itemize}
        \item \textbf{Hommes :} \num{~200000} tirailleurs (AEF/AOF/Maghreb) dont \num{18000} Camerounais. Campagnes : Libye (Bir Hakeim 1942), Italie (Monte Cassino 1944), Provence (Débarquement août 1944), Alsace. \textbf{Prix du sang :} \num{~25000} morts africains.
        \item \textbf{Économie :} \cle{Mobilisation forcée} (prestations, travail obligatoire, réquisitions). Production record : Caoutchouc (Cameroun, Congo), Uranium (Congo), Étain (Nigéria/Cameroun), Arachide (Sénégal). \textbf{Conséquence :} Pénuries alimentaires, inflation, marché noir, révoltes (ex: Thiès 1945, Côte d'Ivoire).
    \end{itemize}
    \item \textbf{B. La rupture politique : De Brazzaville à l'Indépendance.}
    \begin{itemize}
        \item \textbf{Fait :} \cle{Conférence de Brazzaville (Janv-Fév 1944)} : De Gaulle refuse l'indépendance (« hors de question ») mais abolit l'indigénat, instaure citoyenneté, assemblées locales, représentation parlementaire (Lamine Guèye, Houphouët-Boigny).
        \item \textbf{Fait :} \cle{Conférence de la SFIO (1945)} / \cle{Congrès de Bamako (1946)} : Naissance du \cle{RDA} (Rassemblement Démocratique Africain) $\rightarrow$ Parti panafricain indépendantiste.
        \item \textbf{Fait :} \cle{Vétérans} : Retour au pays avec nouvelles compétences, conscience des droits (Déclaration ONU 1942, Charte 1945), frustration (pensions inégales, non-reconnaissance). Base du syndicalisme (USC Cameroun, CGT Afrique) et nationalisme (UPC, RDA, PAIGC).
    \end{itemize}
\end{itemize}

\textbf{CONCLUSION RÉDIGÉE}
\textit{[Bilan]} L'Afrique n'a pas été un simple « réservoir à hommes » passif. Par son ralliement précoce (1940), elle a sauvé l'honneur et la souveraineté française. Par son effort de guerre (sang, matières premières), elle a payé le prix du sang pour une liberté qu'on lui refusait encore. \textit{[Réponse]} La WWII est bien le \textbf{point de bascule} : la Conférence de Brazzaville (1944) ouvre la voie réformiste (Loi Lamine Guèye, Loi Houphouët-Boigny 1946), mais la frustration des vétérans et l'internationalisation (ONU, Guerre froide) poussent au-delà, vers l'indépendance (1960). \textit{[Ouverture]} Cette histoire partagée (Français/Africains) reste au cœur des \textbf{questions mémorielles} actuelles (restitution objets, reconnaissance tirailleurs, pensions, place de l'Afrique dans l'UE/Francophonie).

\end{exoresolu}

\begin{methode}[Comparer deux situations historiques (Tableau comparatif / Parallèle structuré)]
\textbf{Objectif :} Répondre à « Comparez... », « Mettez en parallèle... », « Quelles différences/similitudes entre... ? » (Ex: Vichy vs France Libre au Cameroun ; Europe vs Afrique ; 1914-18 vs 1939-45).
\begin{enumerate}
    \item \textbf{Définir les critères de comparaison (Grille d'analyse) :} Choisir 4 à 5 critères pertinents (Politique, Militaire, Économique, Social, Idéologique, International). \textbf{Ne pas comparer n'importe quoi n'importe comment.}
    \item \textbf{Construire le tableau (2 colonnes + 1 colonne critères) :} Remplir avec des \textbf{faits précis, datés, chiffrés}. Éviter les adjectifs vagues (« important », « fort ») $\rightarrow$ Préférer « \num{18000} engagés », « Décret du 15 mars 1944 ».
    \item \textbf{Rédiger la synthèse (3 phrases types) :}
    \begin{enumerate}
        \item \textbf{Similitude majeure :} « Les deux situations partagent [Critère X] car [Explication structurelle]. »
        \item \textbf{Différence majeure :} « En revanche, elles divergent radicalement sur [Critère Y] : [Situation A] = [Fait], alors que [Situation B] = [Fait], du fait de [Cause profonde]. »
        \item \textbf{Interprétation :} « Cette comparaison montre que [Conclusion historique : rupture/continuité, domination/résistance, etc.]. »
    \end{enumerate}
\end{enumerate}
\textbf{Piège à éviter :} Faire deux paragraphes séparés sans lien (« D'un côté... De l'autre côté... ») sans mot de liaison comparatif (\textit{alors que, contrairement à, de même que, inversement, de façon symétrique}).
\end{methode}

\begin{exoresolu}[Comparaison : L'administration Vichy vs France Libre au Cameroun (1940)]
\textbf{Énoncé :} « Comparez l'attitude et le sort des administrations coloniales vichyste et gaulliste au Cameroun en 1940. » (Tableau + Synthèse).
\tcblower
\textbf{Solution détaillée}

\textbf{Tableau comparatif : Cameroun 1940 — Vichy vs France Libre}

\begin{center}
\begin{tabularx}{\linewidth}{|>{\bfseries}X|X|X|}\hline
\rowcolor{popblueL}
\textbf{Critère de comparaison} & \textbf{Administration VICHY (Juin - Août 1940)} & \textbf{Administration FRANCE LIBRE (Août 1940 +)} \\ \hline
\textbf{Légitimité / Autorité} & Gouverneur \cle{Richard Brunot}. Légalité formelle (décrets Pétain/Laval). Soutien de la majorité des fonctionnaires et planteurs (peur du communisme/gaullisme). & Gouverneur \cle{Pierre Cournarie} (nommé par de Gaulle). Légitimité \textbf{morale et stratégique} (Appel 18 Juin, ralliement Éboué). Soutien de Leclerc (militaire) et d'une minorité active (fonctionnaires, chefs coutumiers ralliés). \\ \hline
\textbf{Position géopolitique} & Alignement sur \cle{Collaboration} (Montoire oct. 1940). Acceptation commission d'armistice allemande/italienne. Neutralité forcée vis-à-vis UK. & Alignement sur \cle{Alliés} (UK/USA). Base arrière pour guerre contre Italie (Libye) et Allemagne. Accords de coopération militaire UK (avions, ravitaillement). \\ \hline
\textbf{Politique indigène / Sociale} & Maintien \cle{Code de l'indigénat}. Travail forcé (\cle{prestations}) intensifié pour l'effort de guerre vichyste (quotas caoutchouc). Répression nationalistes (ex: arrestations USPC). & Maintien initial de l'ordre colonial (pas de rupture immédiate). \textbf{Promesse} réformes (Brazzaville 1944). Recrutement \textbf{volontaire} (BM1) vs forcé. Début égalité soldes (théorique). \\ \hline
\textbf{Économie / Ressources} & Contrôle changes strict (Zone Franc). Exportations vers France métropolitaine (bloquée) ou Allemagne (via commissions). Pénuries, marché noir. & Réorientation \textbf{totale vers Alliés}. Caoutchouc $\rightarrow$ US/UK (accords 1941/42). Étain $\rightarrow$ Alliés. Infrastructure (Route, Rail, Port Douala) développée par/ pour Alliés (Génie US). \\ \hline
\textbf{Issu / Sort} & \textbf{Effondrement sans combat majeur} (27 août 1940). Brunot arrêté, déporté en métropole puis libéré. Fonctionnaires vichystes épurés ou maintenus (continuité administrative). & \textbf{Installation durable (1940-1945)}. Cournarie gouverneur jusqu'en 1943. Base de la \cle{2e DB} (Leclerc). Siège de la Conférence de Brazzaville (1944). \\ \hline
\end{tabularx}
\end{center}

\textbf{Synthèse rédigée (Modèle Bac) :}
\textbf{Similitude structurelle :} Les deux administrations maintiennent l'\textbf{ordre colonial} (hiérarchie raciale, économie d'extraction, chefferies traditionnelles instrumentalisées) et répriment l'opposition nationaliste naissante (UPC/USPC). La rupture n'est pas sociale mais \textbf{géopolitique}.
\textbf{Différence radicale :} L'administration \textbf{Vichy} est \textbf{subie} (légalité formelle, isolement diplomatique, ressources détournées vers l'Axe), tandis que l'administration \textbf{France Libre} est \textbf{choisie et constructive} (ralliement actif, intégration dans l'effort allié, modernisation infrastructurelle, promesse réformatrice de Brazzaville).
\textbf{Interprétation :} Le basculement d'août 1940 illustre la \textbf{fragilité de l'empire colonial} : la légitimité métropolitaine s'effondre face à la force militaire (Leclerc) et à l'intérêt stratégique allié. Le Cameroun devient le \textbf{« berceau de la France Libre »} (selon de Gaulle), transformant un territoire de mandat en acteur fondateur de la IVe République.

\end{exoresolu}

\begin{attention}[Les 5 erreurs qui coûtent des points au Bac (Histoire Terminale Tech)]
\begin{enumerate}
    \item \textbf{Confondre chronologie et causalité :} Écrire « En 1941, Pearl Harbor, donc les USA entrent en guerre » sans expliquer \textit{pourquoi} (embargo pétrole, stratégie japonaise) ni la \textit{conséquence structurelle} (mondialisation, « Germany first »). \textit{Remède :} Utiliser connecteurs logiques (\textit{entraînant, ce qui permit, en réponse à}).
    \item \textbf{Le « Name-dropping » sans analyse :} Lister des noms (De Gaulle, Leclerc, Um Nyobè, Éboué) ou des dates (1940, 1944, 1945) sans les \textbf{lier à l'argument}. \textit{Remède :} « \textbf{Leclerc} incarne la \textbf{continuité militaire} entre l'Afrique (Koufra) et la Métropole (Strasbourg), symbolisant la \textbf{réhabilitation de la France} ». 
    \item \textbf{Oublier l'échelle africaine/camerounaise (Spécificité Tech) :} Traiter la WWII comme un cours d'Histoire générale (Europe/USA/URSS) sans intégrer le \textbf{Dossier Cameroun} (Ralliement 27 août 1940, BM1, Caoutchouc/Étain, Brazzaville 1944, Vétérans/UPC). \textit{Remède :} Au moins 1 argument sur 3 doit concerner l'Afrique/Cameroun dans cette leçon.
    \item \textbf{Anachronisme / Vocabulaire imprécis :} Dire « Hitler a envahi la Russie » (au lieu d'\cle{URSS}), « Les Noirs se sont battus » (au lieu de \cle{Tirailleurs sénégalais} / \cle{BM1} / \cle{Forces Françaises Libres}), « L'ONU a remplacé la SDN » (sans citer \cle{Conférence de San Francisco}, \cle{Charte}, \cle{Conseil de Sécurité}, \cle{Droit de veto}). \textit{Remède :} Fiches de vocabulaire précis par chapitre.
    \item \textbf{Présentation soignée = Points faciles :} Pas de titre, pas de légende sur croquis, pas de paragraphes sautés (intro/dev/conc), ratures, écriture illisible, pas de flèche Nord/échelle sur carte. \textit{Remède :} S'entraîner à rendre une copie « propre » en 30 min (brouillon 10 min, propre 20 min). Le correcteur lit en diagonale : \textbf{structure visible = points acquis}.
\end{enumerate}
\end{attention}

\newpage

\coursec{Exercices}

\coursec{La Deuxième Guerre mondiale}

\courssub{Je vérifie mes connaissances}

\begin{exercice}[QCM : Chronologie et acteurs majeurs]
Parmi les propositions suivantes, une seule est exacte. Recopiez la lettre correspondante.
\begin{enumerate}
\item La date du ralliement du Cameroun à la France libre est :
\begin{itemize}
\item[A.] 27 août 1940
\item[B.] 18 juin 1940
\item[C.] 8 novembre 1942
\item[D.] 6 juin 1944
\end{itemize}
\item La conférence qui a fixé le sort de la Pologne et les nouvelles frontières de l'Europe de l'Est en février 1945 est :
\begin{itemize}
\item[A.] Conférence de Téhéran
\item[B.] Conférence de Yalta
\item[C.] Conférence de Potsdam
\item[D.] Conférence de San Francisco
\end{itemize}
\item Le plan américain de reconstruction de l'Europe lancé en 1947 porte le nom de :
\begin{itemize}
\item[A.] Plan Marshall
\item[B.] Plan Schuman
\item[C.] Plan Morgenthau
\item[D.] Plan Dawes
\end{itemize}
\item L'article de la Charte des Nations Unies qui institue le droit de veto au Conseil de sécurité est l'article :
\begin{itemize}
\item[A.] Article 27
\item[B.] Article 51
\item[C.] Article 1
\item[D.] Article 103
\end{itemize}
\end{enumerate}
\end{exercice}

\begin{exercice}[Vrai / Faux justifié]
Indiquez si les affirmations suivantes sont vraies ou fausses. Justifiez brièvement chaque réponse par un fait précis.
\begin{enumerate}
\item Le Bataillon de Marche n°13 (BM13) camerounais a combattu uniquement sur le front africain.
\item La Conférence de Brazzaville (janvier-février 1944) a posé le principe de l'autonomie progressive des colonies françaises.
\item Le crime contre l'humanité a été défini pour la première fois dans le Statut du Tribunal de Nuremberg (1945).
\item La Déclaration universelle des droits de l'homme (DUDH) a été adoptée par l'Assemblée générale de l'ONU le 10 décembre 1948.
\end{enumerate}
\end{exercice}

\begin{exercice}[Définitions et repères]
Donnez une définition précise (deux à trois lignes) pour chacun des termes suivants en les situant dans le contexte de la Seconde Guerre mondiale :
\begin{enumerate}
\item Guerre d'anéantissement (\emph{Vernichtungskrieg})
\item Solution finale (\emph{Endlösung})
\item Société mobilisée (ou économie de guerre totale)
\item Conférence de San Francisco (avril-juin 1945)
\end{enumerate}
\end{exercice}

\begin{exercice}[Calculs directs et données chiffrées]
À partir des données ci-dessous, effectuez les calculs demandés.
\begin{itemize}
\item Effectifs estimés des tirailleurs camerounais engagés dans la 2\textsuperscript{e} DB et le BM13 : environ 2 500 hommes.
\item Pertes (tués, blessés, disparus) du BM13 lors de la campagne d'Italie et de Provence : 187 hommes.
\item Production de caoutchouc du Cameroun en 1939 : 1 200 tonnes ; en 1944 : 4 800 tonnes.
\end{itemize}
\begin{enumerate}
\item Calculez le taux de pertes du BM13 (en \%). Arrondissez au dixième.
\item Calculez le taux de variation de la production de caoutchouc entre 1939 et 1944 (en \%). Arrondissez au nombre entier.
\item Si le Cameroun a fourni 15 \% de la production totale de caoutchouc de l'Empire français en 1944, quelle était la production totale de l'Empire cette année-là (en tonnes) ?
\end{enumerate}
\end{exercice}

\courssub{Je m'entraîne}

\begin{exercice}[Analyse de document : Extrait de la Conférence de Brazzaville (1944)]
\textbf{Document :} « Les fins de l'œuvre civilisatrice accomplie par la France dans les colonies excluent toute idée d'autonomie, toute possibilité d'évolution hors le bloc français de l'Empire. [...] Il faut écarter tout système qui aurait pour conséquence de substituer à l'administration française une administration internationale. »
\begin{enumerate}
\item Identifiez la nature de ce document, son auteur (ou l'autorité émettrice), sa date et son lieu de rédaction.
\item Replacez ce document dans son contexte historique immédiat (situation militaire, politique en France et dans l'Empire).
\item Expliquez la phrase : « Les fins de l'œuvre civilisatrice accomplie par la France dans les colonies excluent toute idée d'autonomie, toute possibilité d'évolution hors le bloc français de l'Empire. »
\item Quelles ont été les conséquences politiques de cette conférence pour le Cameroun (statut, représentation, mouvements nationalistes) ?
\end{enumerate}
\end{exercice}

\begin{exercice}[Croquis commenté : L'apport du Cameroun à l'effort de guerre allié (1940-1945)]
Réalisez un croquis légendé et organisé sur une carte de base du Cameroun (contour fourni ou à tracer schématiquement). La légende doit comporter \textbf{trois rubriques} obligatoires :
\begin{enumerate}
\item \textbf{Humain} : zones de recrutement (Nord, Ouest, Sud), parcours du BM13 (Yaoundé $\to$ Douala $\to$ Brazzaville $\to$ Levant $\to$ Italie $\to$ Provence), camps d'instruction (Yaoundé, Dschang, Brazzaville).
\item \textbf{Économique} : principales zones de production (caoutchouc : Sud/Est ; bois : Littoral/Sud ; minerais : Nord-Est ; cultures vivrières : Ouest/Adamaoua), axes d'évacuation (chemin de fer Douala-Ngaoundéré, port de Douala, piste vers le Congo).
\item \textbf{Stratégique} : bases aériennes (Douala, Garoua), stations radio (Douala, Yaoundé), rôle dans le ralliement de l'AEF (août 1940, point d'appui pour la colonne Leclerc vers le Nord).
\end{enumerate}
Le croquis doit comporter un titre, une légende hiérarchisée, une échelle indicative et une flèche nord. Utilisez des symboles conventionnels (points, flèches, hachures, couleurs).
\end{exercice}

\begin{exercice}[Synthèse structurée : De la guerre à la naissance de l'ONU]
Rédigez une réponse structurée (introduction, deux parties équilibrées, conclusion) à la question suivante :
« En quoi la Seconde Guerre mondiale a-t-elle constitué un moment de basculement pour l'ordre international, aboutissant à la création de l'ONU et à la proclamation des droits de l'homme ? »
\begin{itemize}
\item Votre réponse doit mobiliser : la faillite de la SDN, les conférences alliées (Téhéran, Yalta, Potsdam, San Francisco), la Charte de l'ONU (organes, principes, droit de veto), le procès de Nuremberg (crimes contre la paix, crimes de guerre, crimes contre l'humanité), la DUDH (1948).
\item Montrez les continuités et les ruptures par rapport à 1919.
\end{itemize}
\end{exercice}

\begin{exercice}[Étude de cas : Le Cameroun, de la colonie au territoire sous tutelle]
À partir de vos connaissances, complétez le tableau comparatif suivant pour mettre en évidence les mutations du statut politique du Cameroun entre 1939 et 1946.

\begin{center}
\begin{tabularx}{\linewidth}{|l|X|X|}\hline
\textbf{Critères} & \textbf{1939 : Colonie sous mandat français (SDN)} & \textbf{1946 : Territoire sous tutelle de l'ONU} \\ \hline
Autorité administrante & & \\ \hline
Base juridique internationale & & \\ \hline
Représentation politique locale & & \\ \hline
Objectif affiché par la puissance administrante & & \\ \hline
Contrôle international & & \\ \hline
\end{tabularx}
\end{center}
Rédigez ensuite un paragraphe de synthèse (10 lignes) expliquant comment la participation du Cameroun à la guerre a accéléré cette évolution statutaire.
\end{exercice}

\courssub{Je me prépare au Bac}

\begin{exercice}[Sujet type Baccalauréat Technique : Composition guidée / Dissertation]
\textbf{Sujet :} « À partir des documents 1 à 4 et de vos connaissances, montrez que la Seconde Guerre mondiale a été un moment de basculement pour l'Empire colonial français et pour le Cameroun en particulier. »

\textbf{Document 1 :} Carte de l'Europe et de la Méditerranée en 1942 (extension maximale de l'Axe, territoires de la France libre, axes de débarquement alliés).
\textbf{Document 2 :} Extrait du discours du Général de Gaulle à Brazzaville, 30 janvier 1944 (Conférence de Brazzaville).
\textbf{Document 3 :} Graphique : Évolution de la production de caoutchouc au Cameroun (1938-1945) — valeurs en tonnes : 1938 : 1 100 ; 1939 : 1 200 ; 1940 : 1 800 ; 1941 : 2 500 ; 1942 : 3 200 ; 1943 : 4 100 ; 1944 : 4 800 ; 1945 : 3 900.
\textbf{Document 4 :} Photographie : Tirailleurs du Bataillon de Marche n°13 (BM13) en Italie, 1944 (uniforme US, insigne de la 2\textsuperscript{e} DB, décorations).

\textbf{Travail demandé :}
\begin{enumerate}
\item \textbf{Analyse des documents (6 points) :}
\begin{itemize}
\item[a.] Présentez chaque document (nature, auteur, date, contexte).
\item[b.] Pour chaque document, identifiez et expliquez deux informations majeures en lien avec le sujet.
\item[c.] Montrez en quoi les documents 2 et 4 illustrent des contradictions de l'Empire français en guerre.
\end{itemize}
\item \textbf{Mise en relation et exploitation des connaissances (10 points) :}
\begin{itemize}
\item[a.] Expliquez pourquoi le ralliement du Cameroun (août 1940) a été un tournant stratégique pour la France libre.
\item[b.] Analysez les contributions humaines (BM13, 2\textsuperscript{e} DB, porteurs) et économiques (caoutchouc, bois, vivriers, bases) du Cameroun à l'effort de guerre allié.
\item[c.] Montrez comment la guerre a transformé les rapports coloniaux : Conférence de Brazzaville, statut de territoire sous tutelle (1946), émergence du nationalisme (USC, UPC).
\end{itemize}
\item \textbf{Synthèse rédigée (4 points) :} Rédigez une conclusion structurée répondant à la problématique : « Basculement militaire, économique et politique : le Cameroun, révélateur de la mutation de l'Empire français. »
\end{enumerate}
\end{exercice}

\begin{exercice}[Sujet type Probatoire : Analyse de document unique]
\textbf{Document :} Extrait du procès-verbal de la Conférence de Brazzaville (30 janvier - 8 février 1944), discours d'ouverture du Général de Gaulle.
« La France a besoin de ses colonies et les colonies ont besoin de la France. [...] Les fins de l'œuvre civilisatrice accomplie par la France dans les colonies excluent toute idée d'autonomie, toute possibilité d'évolution hors le bloc français de l'Empire. [...] Il faut écarter tout système qui aurait pour conséquence de substituer à l'administration française une administration internationale. »

\textbf{Travail demandé :}
\begin{enumerate}
\item \textbf{Identification et contextualisation (4 points) :}
\begin{itemize}
\item[a.] Nature du document, auteur, date, lieu.
\item[b.] Contexte historique : situation de la France, de l'Empire, de la guerre en janvier 1944.
\end{itemize}
\item \textbf{Analyse du contenu (8 points) :}
\begin{itemize}
\item[a.] Expliquez la phrase : « Les fins de l'œuvre civilisatrice accomplie par la France dans les colonies excluent toute idée d'autonomie, toute possibilité d'évolution hors le bloc français de l'Empire. » (Idée de mission civilisatrice, refus de l'indépendance, cadre fédéral).
\item[b.] Quelles réformes concrètes la Conférence de Brazzaville a-t-elle néanmoins décidées pour les colonies (statut, citoyenneté, représentation, travail forcé) ?
\item[c.] En quoi ce discours révèle-t-il l'ambiguïté de la France libre : rupture avec Vichy sur le plan symbolique, continuité sur le plan colonial ?
\end{itemize}
\item \textbf{Ouverture / Conséquences pour le Cameroun (8 points) :}
\begin{itemize}
\item[a.] Quelles ont été les conséquences immédiates de la Conférence pour le Cameroun (décret du 24 octobre 1946, statut de territoire sous tutelle, Assemblée territoriale) ?
\item[b.] Comment les élites camerounaises (Lamine Kaba, Ruben Um Nyobé, Louis-Paul Aujoulat) ont-elles réagi à ces réformes jugées insuffisantes ?
\item[c.] Montrez que la participation du Cameroun à la guerre (ralliement, BM13, effort économique) a constitué un levier pour revendiquer l'égalité des droits et l'autodétermination.
\end{itemize}
\end{enumerate}
\end{exercice}

\begin{exercice}[Évaluation formative : QCM et questions courtes (type Bac blanc)]
Répondez brièvement aux dix questions suivantes. Chaque réponse attendue tient en une ligne maximum.
\begin{enumerate}
\item Date exacte du ralliement du Cameroun à la France libre (jour, mois, année).
\item Nombre approximatif de tirailleurs camerounais engagés dans les Forces françaises libres (FFL) et l'armée de la Libération (1940-1945).
\item Nom du bataillon camerounais qui a combattu en Italie (campagne du Garigliano) et en Provence (débarquement du 15 août 1944).
\item Nom de la conférence fondatrice de l'Organisation des Nations Unies (lieu et dates).
\item Date de la signature de la Charte des Nations Unies à San Francisco.
\item Numéro de l'article de la Charte de l'ONU qui institue le droit de veto des membres permanents du Conseil de sécurité.
\item Définition du « crime contre l'humanité » telle qu'elle figure dans l'article 6(c) du Statut du Tribunal de Nuremberg (1945).
\item Nom de la conférence (février 1945) qui a fixé le sort de la Pologne (frontières, gouvernement) et décidé de la tenue de la Conférence de San Francisco.
\item Nom du plan américain de reconstruction économique de l'Europe annoncé en juin 1947.
\item Date d'adoption de la Déclaration universelle des droits de l'homme (DUDH) par l'Assemblée générale de l'ONU.
\end{enumerate}
\end{exercice}



\newpage

\coursec{Corrigés des exercices}

\begin{corrige}[Exercice 1 --- QCM : Chronologie et acteurs majeurs]
\begin{enumerate}
\item \textbf{A. 27 août 1940.} Après le ralliement du Tchad par le gouverneur Félix Éboué (26 août 1940), le gouverneur du Cameroun Richard Brunot et le colonel Leclerc organisent le ralliement du Cameroun à la France libre le 27 août 1940. Les autres dates correspondent à d'autres événements : 18 juin 1940 (appel du général de Gaulle), 8 novembre 1942 (débarquement allié en Afrique du Nord, opération Torch), 6 juin 1944 (débarquement en Normandie).
\item \textbf{B. Conférence de Yalta.} Réunie du 4 au 11 février 1945 en Crimée, elle rassemble Roosevelt, Churchill et Staline. Elle règle le sort de la Pologne (frontières déplacées vers l'ouest, gouvernement provisoire), le partage de l'Allemagne en zones d'occupation et convoque la conférence de San Francisco.
\item \textbf{A. Plan Marshall.} Annoncé le 5 juin 1947 par le secrétaire d'État américain George Marshall, l'\emph{European Recovery Program} accorde une aide massive (environ 13 milliards de dollars) à la reconstruction de l'Europe occidentale.
\item \textbf{A. Article 27.} L'article 27, alinéa 3, de la Charte prévoit que les décisions du Conseil de sécurité sur les questions de fond requièrent neuf voix sur quinze, \emph{incluant} les voix des cinq membres permanents : c'est le droit de veto.
\end{enumerate}
\textbf{Points de vigilance :} ne pas confondre Yalta (février 1945, guerre en cours) et Potsdam (juillet-août 1945, après la capitulation allemande) ; ne pas attribuer le plan Marshall à 1945 (il date de 1947).
\end{corrige}

\begin{corrige}[Exercice 2 --- Vrai / Faux justifié]
\begin{enumerate}
\item \textbf{Faux.} Le BM13 a combattu hors d'Afrique : campagne d'Italie (Garigliano, printemps 1944) au sein des forces françaises, puis débarquement de Provence (15 août 1944) et campagne de France (vallée du Rhône, Vosges, Alsace) jusqu'en 1945.
\item \textbf{Faux.} La Conférence de Brazzaville (30 janvier -- 8 février 1944) a explicitement \emph{exclu} « toute idée d'autonomie, toute possibilité d'évolution hors le bloc français de l'Empire ». Elle a décidé des réformes (citoyenneté, représentation, suppression annoncée du travail forcé) mais dans le cadre impérial français.
\item \textbf{Vrai.} L'article 6(c) du Statut du Tribunal militaire international (Accord de Londres, 8 août 1945) définit pour la première fois en droit international le crime contre l'humanité : assassinat, extermination, déportation, persécutions pour motifs politiques, raciaux ou religieux contre toute population civile.
\item \textbf{Vrai.} La DUDH est adoptée le 10 décembre 1948 à Paris (Palais de Chaillot) par l'Assemblée générale de l'ONU (résolution 217 A), par 48 voix pour, 0 contre et 8 abstentions.
\end{enumerate}
\textbf{Points de vigilance :} la justification par un fait précis (date, lieu, texte) est exigée ; une réponse « vrai/faux » seule ne vaut aucun point.
\end{corrige}

\begin{corrige}[Exercice 3 --- Définitions et repères]
\begin{enumerate}
\item \textbf{Guerre d'anéantissement (\emph{Vernichtungskrieg}) :} forme de guerre menée par l'Allemagne nazie, surtout à l'Est contre l'URSS à partir de juin 1941, qui vise non seulement la victoire militaire mais la destruction physique des populations (Juifs, Slaves, Tziganes) et l'effacement des États et cultures adverses. Elle abolit la distinction entre combattants et civils.
\item \textbf{Solution finale (\emph{Endlösung}) :} nom de code nazi désignant l'extermination systématique des Juifs d'Europe, coordonnée à la conférence de Wannsee (20 janvier 1942). Mise en œuvre par les fusillades des \emph{Einsatzgruppen}, les ghettos et les centres de mise à mort (Auschwitz-Birkenau, Treblinka), elle fait environ 6 millions de victimes : la Shoah.
\item \textbf{Société mobilisée (économie de guerre totale) :} mobilisation de toutes les ressources d'un État --- humaines (travail des femmes, des colonisés, des prisonniers), industrielles (reconversion des usines vers l'armement), financières (emprunts, réquisitions) et idéologiques (propagande, censure) --- au service de l'effort de guerre. Exemples : le \emph{Victory Program} américain, l'évacuation des usines soviétiques derrière l'Oural.
\item \textbf{Conférence de San Francisco (25 avril -- 26 juin 1945) :} conférence réunissant 50 États qui rédigent et signent la Charte des Nations Unies (26 juin 1945, en vigueur le 24 octobre 1945). Elle fonde l'ONU et ses organes : Assemblée générale, Conseil de sécurité (avec droit de veto des cinq membres permanents), Secrétariat, Cour internationale de Justice, Conseil économique et social, Conseil de tutelle.
\end{enumerate}
\textbf{Points de vigilance :} chaque définition doit comporter une date ou un repère, un contenu précis et un exemple ; éviter les définitions circulaires (« la solution finale est la solution finale du problème juif »).
\end{corrige}

\begin{corrige}[Exercice 4 --- Calculs directs et données chiffrées]
\begin{enumerate}
\item Taux de pertes du BM13 :
\[ \text{Taux} = \frac{187}{2500} \times 100 = 7{,}48 \% \approx \textbf{7{,}5 \%}. \]
\item Taux de variation de la production de caoutchouc entre 1939 et 1944 :
\[ \text{Taux} = \frac{4800 - 1200}{1200} \times 100 = \frac{3600}{1200} \times 100 = 3 \times 100 = \textbf{300 \%}. \]
La production a donc été \emph{multipliée par 4} (augmentation de 300 \%).
\item Si 4 800 tonnes représentent 15 \% de la production totale de l'Empire :
\[ \text{Production totale} = \frac{4800}{0{,}15} = \textbf{32 000 tonnes}. \]
Vérification : $32\,000 \times 0{,}15 = 4\,800$ tonnes. Le calcul est cohérent.
\end{enumerate}
\textbf{Points de vigilance :} bien distinguer taux de variation (300 \%) et coefficient multiplicateur ($\times 4$) ; pour la question 3, on \emph{divise} par 0,15 et on ne multiplie pas ; arrondir uniquement en fin de calcul.
\end{corrige}

\begin{corrige}[Exercice 5 --- Analyse de document : Conférence de Brazzaville (1944)]
\begin{enumerate}
\item \textbf{Présentation :} il s'agit d'un extrait du discours d'ouverture (texte officiel, discours politique) prononcé par le général de Gaulle, président du Comité français de Libération nationale (CFLN), le 30 janvier 1944 à Brazzaville, capitale de l'Afrique équatoriale française.
\item \textbf{Contexte :} en janvier 1944, la guerre tourne en faveur des Alliés (victoire soviétique à Stalingrad en février 1943, débarquement en Italie, préparation du débarquement de Normandie). La France, partiellement libérée (Corse, Afrique du Nord), est dirigée par le CFLN installé à Alger. L'Empire, socle de la France libre depuis 1940, a payé un lourd tribut à la guerre. De Gaulle réunit les gouverneurs pour préparer l'après-guerre et contrer les vues américaines (Charte de l'Atlantique, 1941) favorables à l'autodétermination des peuples colonisés.
\item \textbf{Explication :} la « mission civilisatrice » est l'idéologie qui justifie la colonisation : la France prétend apporter le progrès aux peuples colonisés. Affirmer que cette œuvre « exclut toute idée d'autonomie » signifie que l'indépendance est refusée par principe : les colonies doivent rester dans le « bloc français ». La phrase vise aussi à écarter toute mise sous tutelle internationale (administration étrangère ou onusienne) des colonies françaises.
\item \textbf{Conséquences pour le Cameroun :} les réformes décidées à Brazzaville y sont appliquées : accès à la citoyenneté française, suppression de l'indigénat (1946) et du travail forcé (1946), création de l'Assemblée représentative du Cameroun (ARC, 1946), élection de représentants camerounais aux assemblées constituantes à Paris (Louis-Paul Aujoulat, Alexandre Douala Manga Bell). Le Cameroun devient territoire sous tutelle de l'ONU (13 décembre 1946). Mais le refus de l'autonomie radicalise les élites : naissance de mouvements revendicatifs (syndicats confédérés, puis UPC en 1948) qui réclament l'indépendance et la réunification des deux Cameroun.
\end{enumerate}
\textbf{Points de vigilance :} ne pas présenter Brazzaville comme une conférence « indépendantiste » : c'est l'inverse ; bien dater la tutelle ONU (décembre 1946) et ne pas la confondre avec le mandat SDN (1922).
\end{corrige}

\begin{corrige}[Exercice 6 --- Croquis commenté : L'apport du Cameroun à l'effort de guerre allié]
\textbf{Attendus de la correction (barème indicatif sur 20) :}
\begin{itemize}
\item \textbf{Organisation générale (5 points) :} titre explicite (1 pt) ; contour schématique cohérent du Cameroun (1 pt) ; flèche du nord (1 pt) ; échelle indicative (1 pt) ; légende hiérarchisée en trois rubriques numérotées (1 pt).
\item \textbf{Rubrique Humain (5 points) :} zones de recrutement localisées (Nord, Ouest, Sud) par points ou cercles (2 pts) ; parcours du BM13 tracé en flèches : Yaoundé $\to$ Douala $\to$ Brazzaville $\to$ Levant $\to$ Italie $\to$ Provence (2 pts) ; camps d'instruction signalés (Yaoundé, Dschang) (1 pt).
\item \textbf{Rubrique Économique (5 points) :} zones de production figurées par hachures ou symboles distincts : caoutchouc (Sud/Est), bois (Littoral/Sud), minerais (Nord-Est), cultures vivrières (Ouest/Adamaoua) (3 pts) ; axes d'évacuation : chemin de fer Douala--Ngaoundéré en trait plein, port de Douala (ancre), piste vers le Congo en pointillés (2 pts).
\item \textbf{Rubrique Stratégique (5 points) :} bases aériennes de Douala et Garoua (symbole avion) (2 pts) ; stations radio de Douala et Yaoundé (symbole ondes) (1 pt) ; flèche commentée « Ralliement août 1940 » et « colonne Leclerc vers le Nord (Tchad, Fezzan) » (2 pts).
\end{itemize}
\textbf{Points de vigilance :} un croquis sans légende hiérarchisée ni titre est sanctionné ; les symboles doivent être conventionnels et identiques entre la carte et la légende ; rester schématique mais cohérent géographiquement (Douala sur la côte, Garoua au Nord, Ngaoundéré sur l'Adamaoua).
\end{corrige}

\begin{corrige}[Exercice 7 --- Synthèse structurée : De la guerre à la naissance de l'ONU]
\textbf{Introduction (attendue) :} accroche (conflit le plus meurtrier de l'histoire, environ 50 à 60 millions de morts) ; définition du basculement (fin de l'hégémonie européenne, bipolarisation, refonte du droit international) ; problématique (comment l'échec de la SDN et l'horreur de la guerre conduisent-ils à un ordre international nouveau fondé sur l'ONU et les droits de l'homme ?) ; annonce du plan.

\textbf{I. De la faillite de la SDN à la fondation de l'ONU.}
\begin{itemize}
\item La SDN, créée en 1919, échoue : absence des États-Unis, pas de force armée, règle de l'unanimité ; impuissante face aux agressions (Mandchourie 1931, Éthiopie 1935, remilitarisation de la Rhénanie 1936, Anschluss 1938).
\item Les « Trois Grands » préparent l'après-guerre : Téhéran (novembre-décembre 1943, ouverture du second front), Yalta (février 1945 : sort de l'Allemagne et de la Pologne, principe de l'ONU et du veto), Potsdam (juillet-août 1945 : application des décisions, ultimatum au Japon).
\item San Francisco (avril-juin 1945) : 50 États signent la Charte. Principes : égalité souveraine, interdiction du recours à la force, règlement pacifique des conflits, coopération. Organes : Assemblée générale, Conseil de sécurité avec droit de veto des cinq membres permanents, Conseil de tutelle qui remplace les mandats de la SDN.
\end{itemize}

\textbf{II. Un droit nouveau : Nuremberg et la DUDH.}
\begin{itemize}
\item Le procès de Nuremberg (novembre 1945 -- octobre 1946) institue la responsabilité pénale individuelle des chefs d'État et définit trois chefs d'accusation : crimes contre la paix, crimes de guerre, crimes contre l'humanité (article 6(c)).
\item La DUDH (10 décembre 1948) proclame des droits universels : droits civils et politiques (vie, liberté, opinion) et droits économiques, sociaux et culturels (travail, éducation, santé). Non contraignante juridiquement, elle fonde les pactes de 1966.
\item Continuités et ruptures avec 1919 : ruptures (universalité plus large, sanction pénale des crimes, dimension sociale, tutelle à visée d'autonomie) ; continuités (concert des grandes puissances, souveraineté des États, veto qui limite l'efficacité comme l'unanimité limitait la SDN).
\end{itemize}

\textbf{Conclusion (attendue) :} bilan nuancé (l'ONU naît imparfaite, paralysée par la Guerre froide, mais constitue un cadre durable) ; ouverture vers la décolonisation, prolongement de la logique d'autodétermination.

\textbf{Points de vigilance :} deux parties équilibrées exigées ; chaque conférence doit être datée et dotée d'au moins une décision précise ; le comparatif 1919/1945 doit apparaître explicitement.
\end{corrige}

\begin{corrige}[Exercice 8 --- Étude de cas : Le Cameroun, de la colonie au territoire sous tutelle]
\textbf{Tableau complété :}
\begin{center}
\begin{tabularx}{\linewidth}{|l|X|X|}\hline
\textbf{Critères} & \textbf{1939 : territoire sous mandat français (SDN)} & \textbf{1946 : territoire sous tutelle de l'ONU} \\ \hline
Autorité administrante & France (gouverneur responsable devant Paris) & France (haut-commissaire, responsable devant Paris et devant l'ONU) \\ \hline
Base juridique internationale & Mandat B de la SDN (article 22 du Pacte, 1919 ; mandat confirmé en 1922) & Accord de tutelle approuvé par l'Assemblée générale de l'ONU le 13 décembre 1946 (articles 75-91 de la Charte) \\ \hline
Représentation politique locale & Quasi inexistante : conseils de notables consultatifs, suffrage très restreint & Assemblée représentative (puis territoriale) du Cameroun élue ; représentants aux assemblées françaises \\ \hline
Objectif affiché & « Mise en valeur » du territoire dans le cadre de la mission civilisatrice, sans échéance d'émancipation & Développement politique, économique et social orienté vers l'autonomie ou l'indépendance (article 76 de la Charte) \\ \hline
Contrôle international & Commission permanente des mandats : rapports annuels écrits, sans visite sur place, pouvoir faible & Conseil de tutelle : rapports annuels, visites périodiques sur place, droit de pétition des populations \\ \hline
\end{tabularx}
\end{center}

\textbf{Paragraphe de synthèse (corrigé type) :} la participation du Cameroun à la guerre a accéléré sa mutation statutaire. Le ralliement à la France libre du 27 août 1940 a fait du territoire un pilier de la résistance gaulliste en Afrique : base de départ de la colonne Leclerc vers le Tchad et le Fezzan, contribution humaine (BM13, renforts de la 2\textsuperscript{e} DB, porteurs) et économique (caoutchouc multiplié par quatre, bois, vivriers). Cette contribution a créé une dette politique que la France ne pouvait ignorer à la Libération. La Conférence de Brazzaville (1944), tout en refusant l'autonomie, a acté des réformes : citoyenneté, suppression de l'indigénat et du travail forcé, représentation politique. Surtout, la création de l'ONU a transformé le cadre international : le mandat SDN devient tutelle (13 décembre 1946), avec un objectif explicite d'autonomie, des visites sur place et un droit de pétition. Les élites camerounaises s'emparent de ces leviers : l'UPC d'Um Nyobé adresse des pétitions à l'ONU pour réclamer l'indépendance et la réunification. L'effort de guerre a ainsi transformé les « sujets » coloniaux en acteurs revendiquant l'égalité, ouvrant la voie à l'indépendance de 1960.

\textbf{Points de vigilance :} employer le vocabulaire exact (mandat SDN $\neq$ tutelle ONU ; sujet $\neq$ citoyen) ; citer au moins deux faits de la participation camerounaise dans le paragraphe.
\end{corrige}

\begin{corrige}[Exercice 9 --- Sujet type Baccalauréat Technique : Composition guidée]
\textbf{1. Analyse des documents (6 points, barème indicatif : 1,5 pt par document).}
\begin{itemize}
\item[a.] \textbf{Document 1 :} carte historique de l'Europe et de la Méditerranée en 1942 (document cartographique), montrant l'extension maximale de l'Axe. \textbf{Document 2 :} extrait du discours officiel du général de Gaulle à Brazzaville, 30 janvier 1944 (texte politique). \textbf{Document 3 :} graphique statistique de la production de caoutchouc camerounais, 1938-1945 (source administrative coloniale). \textbf{Document 4 :} photographie de tirailleurs du BM13 en Italie, 1944 (document iconographique, archives ou propagande).
\item[b.] \textbf{Doc 1 :} la France métropolitaine est occupée ou vichyste, donc neutralisée ; l'Afrique équatoriale française et le Cameroun ralliés constituent la base territoriale de la France libre. \textbf{Doc 2 :} refus explicite de l'autonomie coloniale ; promesse de réformes dans le « bloc français ». \textbf{Doc 3 :} la production de caoutchouc quadruple (1 200 t en 1939, 4 800 t en 1944, soit +300 \%) : effort économique majeur ; baisse dès 1945 (3 900 t) avec la fin des réquisitions de guerre. \textbf{Doc 4 :} des soldats camerounais combattent en Europe sous uniforme et avec équipement alliés ; leur présence décorée atteste une contribution militaire directe à la libération.
\item[c.] \textbf{Contradiction :} le document 2 refuse toute autonomie aux colonisés au nom de la « mission civilisatrice », alors que le document 4 montre des Camerounais versant leur sang pour la liberté de la France. De même, le document 3 révèle une exploitation économique intensive (travail forcé encore pratiqué) au service d'une guerre menée au nom de la liberté : l'Empire combat le nazisme tout en niant les droits politiques de ses propres populations.
\end{itemize}

\textbf{2. Mise en relation et connaissances (10 points).}
\begin{itemize}
\item[a.] \textbf{Tournant stratégique du ralliement (août 1940) :} avec le Tchad (Éboué) et le Congo, le Cameroun (gouverneur Brunot, colonel Leclerc) donne à la France libre un territoire, une armée (tirailleurs), des ressources (caoutchouc, bois, or) et des infrastructures (aérodromes de Douala et Garoua, radios). Sans ce socle africain, la France libre serait restée un gouvernement en exil sans armée ni assise ; Brazzaville devient sa capitale symbolique et la colonne Leclerc part du Cameroun-Tchad vers Koufra (1941), le Fezzan puis la Tunisie.
\item[b.] \textbf{Contributions :} humaines --- environ 2 500 engagés camerounais (BM13, renforts de la 2\textsuperscript{e} DB, porteurs, marins) ; le BM13 combat en Italie (Garigliano, 1944), débarque en Provence (15 août 1944) et combat jusqu'en Alsace (187 pertes, soit 7,5 \%). Économiques --- caoutchouc multiplié par quatre, bois, cultures vivrières pour les troupes, or pour le financement ; bases aériennes et portuaires (Douala, escale transafricaine).
\item[c.] \textbf{Transformation des rapports coloniaux :} Brazzaville (1944) opère une rupture limitée (fin annoncée de l'indigénat et du travail forcé, citoyenneté, représentation) dans une continuité impériale (refus de l'autonomie). Le statut de territoire sous tutelle (13 décembre 1946) internationalise la question camerounaise : visites de l'ONU, droit de pétition. Le nationalisme s'organise : syndicats confédérés (grèves de 1945), puis UPC (1948) de Ruben Um Nyobé, qui réclame indépendance et réunification en s'appuyant sur la Charte de l'ONU et le sacrifice de guerre.
\end{itemize}

\textbf{3. Synthèse rédigée (4 points) :} la Seconde Guerre mondiale a été un triple basculement pour le Cameroun. Militairement, le ralliement de 1940 et les combats du BM13 ont fait du territoire un sauveur de la France libre. Économiquement, la mise à contribution intensive (caoutchouc, bois, vivriers) a révélé sa valeur stratégique. Politiquement, le sang versé et le nouveau droit international (Charte, tutelle, DUDH) ont rendu intenable l'ancien statut colonial : la France a dû réformer (Brazzaville, tutelle de 1946), mais son refus de l'autonomie a radicalisé le nationalisme (UPC), engageant la marche vers l'indépendance de 1960. Le Cameroun apparaît ainsi comme le révélateur de la mutation de tout l'Empire français.

\textbf{Points de vigilance :} chaque document doit être présenté (nature, auteur, date, contexte) avant exploitation ; les calculs du document 3 (quadruplement, +300 \%) doivent être explicités ; la synthèse doit répondre à la problématique et non répéter l'analyse.
\end{corrige}

\begin{corrige}[Exercice 10 --- Sujet type Probatoire : Analyse de document unique (Brazzaville)]
\textbf{1. Identification et contextualisation (4 points).}
\begin{itemize}
\item[a.] Extrait du discours d'ouverture de la Conférence de Brazzaville, prononcé par le général de Gaulle, chef du CFLN, le 30 janvier 1944 à Brazzaville (AEF). Nature : discours politique officiel.
\item[b.] Contexte : janvier 1944, la guerre bascule en faveur des Alliés (Stalingrad, Italie, débarquement de Normandie en préparation) ; la France est dirigée par le CFLN à Alger ; l'Empire a largement contribué à l'effort de guerre ; les États-Unis évoquent l'autodétermination des peuples (Charte de l'Atlantique, 1941). De Gaulle veut fixer l'avenir de l'Empire avant la conférence de paix.
\end{itemize}

\textbf{2. Analyse du contenu (8 points).}
\begin{itemize}
\item[a.] La « mission civilisatrice » est l'idéologie justifiant la colonisation : la France prétend élever les peuples colonisés vers le progrès. Dire qu'elle « exclut toute idée d'autonomie » signifie que l'indépendance est refusée par principe, car elle serait un abandon de cette mission ; « hors le bloc français » exclut aussi toute tutelle internationale : les colonies doivent évoluer uniquement dans l'ensemble impérial français (future Union française de 1946).
\item[b.] Réformes décidées : accès des colonisés à la citoyenneté française ; suppression annoncée du code de l'indigénat et du travail forcé (effectives en 1946) ; création d'assemblées locales élues ; représentation des territoires dans les assemblées françaises ; effort d'investissement (FIDES) et développement de l'enseignement.
\item[c.] Ambiguïté : rupture avec Vichy sur le plan symbolique et républicain (restauration des libertés, fin des législations d'exception, refus de la collaboration) ; continuité sur le plan colonial (maintien de la domination française, refus de l'indépendance, primat de l'administrateur). La France libre combat pour la liberté en Europe tout en la refusant aux colonisés.
\end{itemize}

\textbf{3. Ouverture : conséquences pour le Cameroun (8 points).}
\begin{itemize}
\item[a.] Conséquences immédiates : application des réformes (citoyenneté, fin de l'indigénat et du travail forcé en 1946) ; le Cameroun devient territoire sous tutelle de l'ONU par l'accord du 13 décembre 1946 ; création de l'Assemblée représentative du Cameroun (ARC) ; élection de représentants camerounais à Paris.
\item[b.] Réactions des élites : Louis-Paul Aujoulat travaille dans le cadre français (évolution au sein de l'Union française) ; les syndicalistes (USC) mènent des revendications sociales et politiques (grèves de 1945) ; Ruben Um Nyobé et l'UPC (1948) jugent les réformes insuffisantes, exigent l'indépendance immédiate et la réunification des deux Cameroun, et saisissent l'ONU par pétitions (1952, 1954).
\item[c.] Levier de la participation à la guerre : le ralliement de 1940 (choix libre du territoire), les combats du BM13 (sang versé pour la liberté) et l'effort économique (caoutchouc, bois, vivriers) fondent un argument moral et juridique : les droits proclamés (Charte de l'ONU, autodétermination ; DUDH de 1948) doivent s'appliquer aux Camerounais. Le droit international né de la guerre est ainsi retourné contre la puissance administrante.
\end{itemize}
\textbf{Points de vigilance :} ne pas écrire que Brazzaville « accorde l'autonomie » ; dater précisément la tutelle (13 décembre 1946) ; citer au moins deux réformes concrètes et deux personnalités camerounaises.
\end{corrige}

\begin{corrige}[Exercice 11 --- Évaluation formative : QCM et questions courtes]
\begin{enumerate}
\item 27 août 1940.
\item Environ 2 500 hommes (BM13, renforts de la 2\textsuperscript{e} DB, marins, porteurs).
\item Le Bataillon de Marche n\textsuperscript{o}13 (BM13).
\item La Conférence de San Francisco (25 avril -- 26 juin 1945).
\item Le 26 juin 1945.
\item L'article 27 (alinéa 3) de la Charte.
\item « L'assassinat, l'extermination, la réduction en esclavage, la déportation et tout autre acte inhumain commis contre toute population civile, avant ou pendant la guerre, ou les persécutions pour des motifs politiques, raciaux ou religieux » (article 6(c) du Statut de Nuremberg).
\item La Conférence de Yalta (4 -- 11 février 1945).
\item Le plan Marshall (\emph{European Recovery Program}, juin 1947).
\item Le 10 décembre 1948.
\end{enumerate}
\textbf{Points de vigilance :} une ligne maximum par réponse ; les dates doivent être complètes (jour, mois, année quand elles sont demandées) ; ne pas confondre la signature de la Charte (26 juin 1945) et son entrée en vigueur (24 octobre 1945).
\end{corrige}

\newpage

\coursec{Fiche bilan}

\begin{popfigure}
\begin{tikzpicture}[
    mindmap,
    grow cyclic,
    every node/.style={concept, execute at begin node=\hskip0pt},
    root concept/.append style={concept color=popblue, minimum size=3.5cm, text width=3cm, font=\bfseries\large, text=white},
    level 1 concept/.append style={level distance=4.5cm, sibling angle=360/5, minimum size=2.5cm, text width=2.2cm, font=\small\bfseries, text=white},
    level 2 concept/.append style={level distance=3.5cm, sibling angle=45, minimum size=1.8cm, text width=1.6cm, font=\footnotesize, text=white},
    concept color=popblue
]
\node[root concept, align=center]{Seconde\\ Guerre\\ Mondiale\\ 1939-1945}
    child[concept color=poporange] { node[level 1 concept, align=center]{1. Genèse\\ \& Extension\\ (1939-1941)}
        child[concept color=poporange!70!popgold] { node[level 2 concept] {Traité Versailles} }
        child[concept color=poporange!70!popgold] { node[level 2 concept] {Totalitarismes} }
        child[concept color=poporange!70!popgold] { node[level 2 concept] {Blitzkrieg} }
        child[concept color=poporange!70!popgold] { node[level 2 concept] {Barbarossa} }
        child[concept color=poporange!70!popgold] { node[level 2 concept] {Pearl Harbor} }
    }
    child[concept color=popred] { node[level 1 concept, align=center]{2. Violence\\ Paroxystique}
        child[concept color=popred!70!poppink] { node[level 2 concept] {Shoah} }
        child[concept color=popred!70!poppink] { node[level 2 concept] {Guerre anéantissement} }
        child[concept color=popred!70!poppink] { node[level 2 concept] {Société mobilisée} }
        child[concept color=popred!70!poppink] { node[level 2 concept] {Crimes humanité} }
        child[concept color=popred!70!poppink] { node[level 2 concept] {Hiroshima/Nagasaki} }
    }
    child[concept color=popgreen] { node[level 1 concept, align=center]{3. Afrique\\ \& Cameroun}
        child[concept color=popgreen!70!popblue] { node[level 2 concept] {Ralliement 1940} }
        child[concept color=popgreen!70!popblue] { node[level 2 concept] {Force publique} }
        child[concept color=popgreen!70!popblue] { node[level 2 concept] {Effort guerre} }
        child[concept color=popgreen!70!popblue] { node[level 2 concept] {Brazzaville 1944} }
        child[concept color=popgreen!70!popblue] { node[level 2 concept] {Conscience politique} }
    }
    child[concept color=poppurple] { node[level 1 concept, align=center]{4. 1945 : Monde\\ Nouveau \& ONU}
        child[concept color=poppurple!70!poppink] { node[level 2 concept] {Conférences Yalta/Potsdam} }
        child[concept color=poppurple!70!poppink] { node[level 2 concept] {Bipolarisation} }
        child[concept color=poppurple!70!poppink] { node[level 2 concept] {Charte San Francisco} }
        child[concept color=poppurple!70!poppink] { node[level 2 concept] {DUDH 1948} }
        child[concept color=poppurple!70!poppink] { node[level 2 concept] {ONU : CS, AG, CIJ} }
    }
    child[concept color=popgold] { node[level 1 concept, align=center]{5. Héritages\\ \& Mémoires}
        child[concept color=popgold!70!poporange] { node[level 2 concept] {Décolonisation} }
        child[concept color=popgold!70!poporange] { node[level 2 concept] {Construction UE} }
        child[concept color=popgold!70!poporange] { node[level 2 concept] {Droit humanitaire} }
        child[concept color=popgold!70!poporange] { node[level 2 concept] {Devoir mémoire} }
        child[concept color=popgold!70!poporange] { node[level 2 concept] {Guerre froide} }
    };
\end{tikzpicture}
\legende{Carte mentale de synthèse : La Deuxième Guerre mondiale (Programme Terminale Technique)}
\end{popfigure}

\begin{bilan}

\textbf{Définitions clés (à connaître par cœur)}
\begin{itemize}
    \item \cle{Guerre mondiale} : Conflit impliquant les grandes puissances sur plusieurs continents (1939-1945 : Axe vs Alliés).
    \item \cle{Guerre d'anéantissement} : Guerre visant la destruction physique de l'ennemi (populations, infrastructures), sans distinction combattants/civils (ex : front de l'Est).
    \item \cle{Génocide / Shoah} : Extermination systématique et planifiée d'un groupe humain (6 millions de Juifs par le IIIe Reich).
    \item \cle{Crimes contre l'humanité} : Notion juridique née à Nuremberg (1945) : meurtre, extermination, esclavage, déportation contre populations civiles.
    \item \cle{Société mobilisée / Économie de guerre} : Mise au service de l'effort militaire de l'ensemble des ressources humaines, industrielles, financières et propagandistes.
    \item \cle{ONU (Organisation des Nations Unies)} : Organisation internationale créée en 1945 (Charte de San Francisco) pour maintenir la paix et la sécurité internationales.
    \item \cle{Décolonisation} : Processus d'accession à l'indépendance des colonies (1945-1975), accéléré par l'affaiblissement des métropoles et la Charte de l'ONU (droit des peuples à disposer d'eux-mêmes).
\end{itemize}

\textbf{Repères chronologiques \& Acteurs majeurs}
\begin{center}
\begin{tabularx}{\linewidth}{|>{\bfseries}l|X|X|}\hline
\textbf{Période / Événement} & \textbf{Dates clés} & \textbf{Acteurs / Décisions} \\ \hline
Genèse & 1939 (1er sept) & Hitler (Allemagne), Staline (Pacte germano-soviétique), Daladier/Chamberlain \\ \hline
Extension & 1941 (Juin/Déc) & Opération Barbarossa ; Pearl Harbor (Roosevelt, Tōjō) → Mondialisation \\ \hline
Tournants & 1942-1943 & Stalingrad, El Alamein, Midway → Recul de l'Axe \\ \hline
Fin Europe & 8 mai 1945 & Capitulation Allemagne (Keitel) ; Suicide Hitler (30 avril) \\ \hline
Fin Pacifique & 2 sept 1945 & Bombardements atomiques (Truman) : Hiroshima (6 août), Nagasaki (9 août) \\ \hline
Conférences & 1943-1945 & Téhéran, Yalta (Fév 45), Potsdam (Juil 45) : Partage monde, ONU, Allemagne \\ \hline
Naissance ONU & 26 juin 1945 & Signature Charte San Francisco (51 États) ; Siège New York \\ \hline
Conf. Brazzaville & Janv-Fév 1944 & De Gaulle, Gouverneurs AEF : Fin code indigénat, citoyenneté, assemblées locales \\ \hline
\end{tabularx}
\end{center}

\textbf{Méthodes express (Bac)}
\begin{methode}[Commentaire de document historique (Type Bac)]
\begin{enumerate}
    \item \textbf{Identification} : Nature, auteur, date, contexte historique précis (quelques lignes).
    \item \textbf{Analyse} : Dégager 2-3 idées forces (thèmes) illustrées par des citations courtes (« ... »).
    \item \textbf{Explication / Mise en perspective} : Mobiliser les connaissances du cours (causes, conséquences, acteurs) pour expliquer le document. Nuancer (portée, limites).
    \item \textbf{Conclusion} : Répondre à la problématique implicite ; ouvrir sur la conséquence majeure ou l'héritage.
\end{enumerate}
\end{methode}

\begin{methode}[Dissertation / Sujet de réflexion (Type Bac)]
\begin{enumerate}
    \item \textbf{Analyse du sujet} : Définir termes, bornes chrono/spatiales, problématique (question tension).
    \item \textbf{Plan détaillé} : 2 ou 3 parties équilibrées (ex : I. Genèse/Extension ; II. Violence/Anéantissement ; III. Conséquences/Mondialisation). Sous-parties argumentées.
    \item \textbf{Rédaction} : Intro (accroche, définitions, sujet, problématique, annonce plan). Développement (arguments + exemples précis + connecteurs logiques). Conclusion (bilan, réponse problèmatique, ouverture).
\end{enumerate}
\end{methode}

\textbf{Focus Cameroun (Dossier 1 - Exigible)}
\begin{itemize}
    \item \textbf{Ralliement} : Août 1940 (Gouverneur Félix Éboué / Leclerc / De Gaulle) : Premier territoire français d'Afrique à rejoindre la France Libre.
    \item \textbf{Contribution militaire} : Tirailleurs camerounais (Régiment de marche du Cameroun) → Campagnes du Fezzan (Libye), Tunisie, France (Libération de Paris, Alsace), Allemagne.
    \item \textbf{Effort économique} : Production stratégique (caoutchouc, étain, café, cacao, bois) pour l'effort de guerre allié ; travail forcé (prestations) accentué.
    \item \textbf{Conséquences politiques} : Conférence de Brazzaville (1944) → Fin indigénat, citoyenneté française, députés à l'Assemblée constituante (1945), naissance partis politiques (UPC 1948).
\end{itemize}

\end{bilan}

\begin{aretenir}[Checklist avant le Bac]
$\square$ Je sais définir : Guerre mondiale, Génocide/Shoah, Crime contre l'humanité, Économie de guerre, ONU, Décolonisation.
$\square$ Je maîtrise la chronologie : 1939 (début), 1941 (mondialisation), 1942-43 (tournants), 8 mai 1945 (Europe), 2 sept 1945 (Japon).
$\square$ J'explique les causes : Traité de Versailles, montée totalitarismes (nazisme, fascisme, militarisme japonais), échec SDN, politique apaisement.
$\square$ Je caractérise la violence paroxystique : Shoah (processus : exclusion $\rightarrow$ ghettoïsation $\rightarrow$ solution finale), guerre d'anéantissement à l'Est, bombardements stratégiques, bombes atomiques.
$\square$ Je décris la société mobilisée : Propagande, rationnement, réquisition main-d'œuvre (STO), rôle femmes, résistance/collaboration.
$\square$ Je détaille le rôle de l'Afrique / Cameroun : Ralliement 1940, contribution militaire (Fezzan, France), effort économique, Conférence Brazzaville 1944.
$\square$ Je connais la naissance de l'ONU : Conférences alliées, Charte San Francisco (26 juin 1945), objectifs (paix, DDH, développement), organes (CS, AG, SG, CIJ).
$\square$ J'analyse les conséquences géopolitiques : Bipolarisation (USA/URSS), déclin Europe, décolonisation (Asie puis Afrique), construction européenne (CECA 1951).
$\square$ Je sais exploiter un document : Identifier (nature, auteur, date, contexte), Analyser (idées forces + citations), Expliquer (connaissances), Conclure (portée).
$\square$ Je structure une dissertation : Intro (amorce, définitions, problématique, plan), Développement (parties équilibrées, arguments + exemples), Conclusion (bilan, ouverture).
$\square$ Je connais les mots-clés du programme : \cle{Blitzkrieg}, \cle{Barbarossa}, \cle{Solution finale}, \cle{Nuremberg}, \cle{Yalta/Potsdam}, \cle{Brazzaville 1944}, \cle{DUDH 1948}.
$\square$ Je relie l'histoire à l'actualité : Héritages ONU/DDH, enjeux mémoire, construction européenne, racines décolonisation.
\end{aretenir}

\findecours

\end{document}
